\chapter{李群与李代数}
\label{ch:lie}

% ════════════════════════════════════════════════════════════════
%  黄金范本章：吸收 视觉SLAM十四讲 ch4 + Barfoot SER2 ch7/8 + SLAM Handbook ch2
%  + Solà micro-Lie(1812.01537)，自包含独立记录，右扰动为主。
%  教学层遵循 v5.0 理论教学规范（动机→反面→历史→理论→陷阱→练习 + 符号表/速查/故障排查）。
% ════════════════════════════════════════════════════════════════

\begin{practice}[前置自测]
开始之前，先试答下列问题。若已能流畅作答，可只速读本章表格与速查卡；若卡住，正是本章要为你解决的。
\begin{enumerate}
  \item 两个旋转矩阵 $\mathbf{R}_1,\mathbf{R}_2\in\SO(3)$ 相加，结果还是旋转矩阵吗？相乘呢？
  \item 优化一个相机位姿 $\mathbf{T}$ 时，为什么不能像优化普通向量那样对 $\mathbf{T}$ 的 $16$ 个元素直接做梯度下降？
  \item 旋转向量（轴角）$\boldsymbol{\phi}=\theta\mathbf{a}$ 与旋转矩阵 $\mathbf{R}$ 之间的换算公式叫什么？它和“矩阵指数”有何关系？
  \item 给定误差函数 $\mathbf{e}(\mathbf{R})=\mathbf{R}\mathbf{p}-\mathbf{z}$，你打算如何定义并计算 $\partial\mathbf{e}/\partial\mathbf{R}$？“对一个矩阵求导”在这里合法吗？
  \item 在 SLAM 代码里把位姿更新写成 $\mathbf{R}\leftarrow\mathbf{R}\,\mathrm{Exp}(\delta\boldsymbol{\phi})$ 与 $\mathbf{R}\leftarrow\mathrm{Exp}(\delta\boldsymbol{\phi})\,\mathbf{R}$，二者有何区别？哪个是“在机体系下扰动”？
\end{enumerate}
\end{practice}

\section*{本章目标}
学完本章，你应当能够：
\begin{enumerate}
  \item \textbf{说清}旋转/位姿为何构成\emph{群}而非向量空间，以及这给优化带来的根本困难；
  \item \textbf{推导} $\SO(3)$、$\SE(3)$ 的李代数 $\so(3)$、$\se(3)$，并解释“李代数 = 单位元处切空间”的含义；
  \item \textbf{掌握}指数/对数映射（罗德里格斯公式、$\SE(3)$ 指数与左雅可比 $\mathbf{J}$），并理解 $\boldsymbol{\rho}\neq\mathbf{t}$ 的来由；
  \item \textbf{运用} BCH 线性近似与左/右雅可比 $\mathbf{J}_l,\mathbf{J}_r$，在李代数上做加减与微分；
  \item \textbf{独立完成}右扰动模型下旋转点、位姿作用点对扰动的求导，并说明它为何优于直接李代数求导；
  \item \textbf{使用}伴随 $\mathrm{Ad}(\mathbf{T})$ 在左右扰动、不同坐标系之间精确转换；
  \item \textbf{解释} $\mathbf{J}$ 的运动学含义 $\boldsymbol{\omega}=\mathbf{J}\dot{\boldsymbol{\phi}}$，把指数映射里的雅可比与角速度/预积分焊接起来；
  \item \textbf{联系}流形 / retraction / $\boxplus\,\boxminus$ 的现代语言，理解右扰动如何对接因子图与 GTSAM/Ceres 求解器；
  \item \textbf{掌握}李群上的距离、体积元 $|\det\mathbf{J}|$ 与积分，并据此把高斯密度严格地定义在群上；
  \item \textbf{处理}李群上的高斯不确定度、位姿复合/求逆/相关差分的协方差传播，以及多位姿估计的融合；
  \item \textbf{实现}李群上的保群（测地）插值，并理解它为何等价于匀速运动、优于逐元素线性插值；
  \item \textbf{建立}扩展位姿群 $\mathrm{SE}_2(3)$（旋转 + 位置 + 速度）的群结构、指数映射与伴随，作为 InEKF/扩展位姿预积分的统一地基；
  \item \textbf{用群论语言}陈述对称、不变与等变，并据此理解“状态估计不依赖参考系选择”与 InEKF“线性化误差与状态无关”的一致性论断。
\end{enumerate}

\subsection*{本章知识导航}
本章是一条“\emph{把带约束的旋转/位姿，化为可自由加减、可微分的向量}”的主线。结构如下：

\begin{center}
\small
\begin{tikzpicture}[
  node distance=4mm and 6mm, >=Stealth,
  blk/.style={draw, rounded corners, align=center, font=\small, inner sep=3pt, fill=main!6, draw=main!60!black},
  grp/.style={draw, rounded corners, align=center, font=\small, inner sep=3pt, fill=second!6, draw=second!60!black},
  ln/.style={->, thick, gray!60!black}]
\node[blk] (group) {群 / 矩阵李群\\ $\SO(3),\SE(3)$};
\node[blk, right=of group] (alg) {李代数\\ $\so(3),\se(3)$};
\node[blk, right=of alg] (exp) {指数 / 对数\\ Rodrigues, $\mathbf{J}$};
\node[grp, below=8mm of group] (bch) {BCH\\ $\mathbf{J}_l,\mathbf{J}_r$};
\node[grp, right=of bch] (pert) {扰动模型\\（右扰动为主）};
\node[grp, right=of pert] (adj) {伴随\\ $\mathrm{Ad},\mathrm{ad}$};
\node[blk, below=8mm of bch, fill=third!8, draw=third!60!black] (man) {流形 / retraction\\ $\boxplus\,\boxminus$（对接因子图）};
\node[blk, right=of man, fill=third!8, draw=third!60!black] (unc) {李群上的\\不确定度};
\draw[ln] (group)--(alg); \draw[ln] (alg)--(exp);
\draw[ln] (exp.south) -- ++(0,-4.5mm) -| (bch.north);
\draw[ln] (bch)--(pert); \draw[ln] (pert)--(adj);
\draw[ln] (pert.south)-- (man.north);
\draw[ln] (man)--(unc);
\end{tikzpicture}
\end{center}

\noindent\textbf{推荐路径}：第 1--4 节是任何读者的主干（建立“群—代数—指数”三角）；第 5--7 节是\emph{优化所需}的核心工具（BCH、运动学 $\boldsymbol{\omega}=\mathbf{J}\dot{\boldsymbol{\phi}}$、扰动、伴随），SLAM 后端、VIO、标定章节会反复引用——其中 $\boldsymbol{\omega}=\mathbf{J}\dot{\boldsymbol{\phi}}$（\cref{thm:so3-kinematics}）是连续旋转预积分的前置定理；流形小节用现代语言把前面统一，并说明本书\emph{右扰动}约定的工程归宿（因子图）；不确定度一节面向状态估计，先补李群上的距离/体积/积分（密度归一化的地基），再讲切空间高斯、复合/求逆/相关差分的协方差传播与多位姿融合，可在读完滤波章后回看。带（选读）标记的插值、Sim(3)、$\mathrm{SE}_2(3)$、对称/不变/等变、不确定度四阶项、附录的最小多项式/Jordan 分解为进阶/选读——其中李群插值（\cref{sec:lie-interp}）是连续时间轨迹（STEAM）的几何地基；$\mathrm{SE}_2(3)$（\cref{sec:se23}）是全书扩展位姿与不变滤波的权威定义处；对称/不变/等变（\cref{sec:lie-symmetry}）为 InEKF 一致性论断提供群论语言。

\subsection*{前置知识桥接}
\cref{ch:rigid_body} 已给出旋转的多种\emph{表示}——旋转矩阵 $\mathbf{R}$、旋转向量（轴角）$\boldsymbol{\phi}=\theta\mathbf{a}$、欧拉角、四元数，以及罗德里格斯公式 $\mathbf{R}=\cos\theta\,\mathbf{I}+(1-\cos\theta)\mathbf{a}\mathbf{a}^\top+\sin\theta\,\mathbf{a}^\wedge$ 的“正向”用法。本章要解决的是上一章\emph{回避}的问题：当 $\mathbf{R},\mathbf{T}$ 是\textbf{待估计、待优化}的未知量时，如何对它们求导、做迭代更新。换言之，从“如何表示旋转”进到“\emph{如何估计旋转}”。

\subsection*{如果跳过本章会怎样}
\begin{itemize}
  \item 在 \cref{ch:nlopt}（非线性优化）里，你会写出形如 $\min_{\mathbf{T}}\sum_i\|\mathbf{z}_i-\mathbf{T}\mathbf{p}_i\|^2$ 的问题，却\emph{无法}计算 $\partial/\partial\mathbf{T}$——因为 $\SE(3)$ 没有加法，求导定义直接失效；
  \item 在 VIO / 位姿图优化里，你会发现 Ceres、GTSAM 都要求你提供一个“流形局部参数化（local parameterization / manifold）”，而不解释清楚它就是本章的\emph{右扰动 + 指数映射}，你将只能照抄而不会调试。
\end{itemize}

\begin{table}[H]\centering\small
\caption{本章阅读方式建议}
\begin{tabular}{lll}
\toprule
阅读方式 & 时间 & 适合谁 \\
\midrule
精读（含全部推导与练习） & 5--7 小时 & 首次系统学习、要做后端/标定推导 \\
速读（跳过冗长推导细节） & 2 小时 & 有李群基础、想对齐本书记号与右扰动约定 \\
速查（只看小结的符号表 + 定理速查表） & 15 分钟 & 写代码时回来查公式 \\
\bottomrule
\end{tabular}
\end{table}

\begin{note}
\textbf{本章记号与约定.}\;
旋转矩阵记 $\mathbf{R}\in\SO(3)$（Barfoot\cite{barfoot2024state} 用 $\mathbf{C}$，本书统一 $\mathbf{R}$，与十四讲\cite{gaoxiang2019slam14}、Eigen 一致）；位姿 $\mathbf{T}\in\SE(3)$。
$\so(3)$ 元素 $\boldsymbol{\phi}\in\mathbb{R}^3$，经反对称算子 $(\cdot)^\wedge$ 化为 $\boldsymbol{\phi}^\wedge\in\mathbb{R}^{3\times3}$，逆为 $(\cdot)^\vee$。
$\se(3)$ 元素 $\boldsymbol{\xi}=[\boldsymbol{\rho}^\top,\boldsymbol{\phi}^\top]^\top\in\mathbb{R}^6$（\textbf{平移分量 $\boldsymbol{\rho}$ 在前、旋转 $\boldsymbol{\phi}$ 在后}，与十四讲、Sophus、Barfoot 一致）。
记 $\mathrm{Exp}(\boldsymbol{\phi}):=\exp(\boldsymbol{\phi}^\wedge)$、$\mathrm{Log}(\mathbf{R}):=\ln(\mathbf{R})^\vee$（吃/吐向量的便捷算子，沿用 Handbook\cite{carlone2026handbook}、Solà\cite{sola2018micro}）。
\textbf{扰动方向：全书以右扰动 $\mathbf{R}\,\mathrm{Exp}(\delta\boldsymbol{\phi})$ 为主}（局部坐标，对接因子图/自动微分），右雅可比 $\mathbf{J}_r$ 为主；引用左扰动来源（Barfoot、十四讲）处并列给出经伴随的转换。\textbf{注意 $\boldsymbol{\rho}\neq\mathbf{t}$}（见 \cref{thm:se3-exp}）。
\end{note}

% ════════════════════════════════════════════════════════════════
\section{为什么需要李群与李代数}
\label{sec:why-lie}

\paragraph{动机：位姿是待估计量。}
在 SLAM 中，相机/机器人的位姿 $\mathbf{R},\mathbf{t}$ 并非已知，而是要从观测里\emph{估计}出来。典型做法是构造优化问题：寻找最优位姿使预测观测与实际观测的误差最小\cite{gaoxiang2019slam14}。例如已知世界系下路标 $\mathbf{p}$ 与含噪观测 $\mathbf{z}$，理想满足 $\mathbf{z}=\mathbf{T}\mathbf{p}$，于是最小化
\begin{equation}
  \min_{\mathbf{T}}\ J(\mathbf{T})=\sum_{i=1}^{N}\big\|\mathbf{z}_i-\mathbf{T}\mathbf{p}_i\big\|_2^2.
  \label{eq:pose-opt}
\end{equation}
求解 \cref{eq:pose-opt} 需要 $J$ 关于 $\mathbf{T}$ 的导数，以便迭代地调整当前估计。

\paragraph{反面：直接拿矩阵当优化变量会卡住。}
旋转矩阵自带约束：正交（$\mathbf{R}\mathbf{R}^\top=\mathbf{I}$）且 $\det\mathbf{R}=1$。若把 $\mathbf{R}$ 的 $9$ 个元素当成自由变量做梯度下降，每走一步都会\emph{破坏正交性}，迭代结果不再是旋转。强行带约束优化（如拉格朗日乘子）则代价高、易病态。问题的根子在于：$\SO(3)$、$\SE(3)$ \textbf{不是向量空间}，没有自由的加减法，求导的差商定义 $\lim(\,f(\mathbf{R}+\Delta\mathbf{R})-f(\mathbf{R})\,)/\Delta\mathbf{R}$ 里 “$\mathbf{R}+\Delta\mathbf{R}$” 根本不合法。

\paragraph{历史：从挪威数学家 Lie 到机器人学。}
$19$ 世纪挪威数学家 Sophus Lie 研究连续变换群（今称\emph{李群}）时发现：一个连续群在单位元附近的全部局部信息，可由一个\emph{线性}对象——它的\emph{李代数}——精确编码，二者经指数映射相连。机器人学自 Chirikjian\cite{chirikjian2011stochastic} 等人起系统地把这套工具引入位姿估计；今天它是 Barfoot《状态估计》\cite{barfoot2024state}、SLAM Handbook\cite{carlone2026handbook} 的共同语言，也是 GTSAM、Sophus、Ceres、manif 等库的数学内核。

\paragraph{核心思想。}
李群（一个\emph{流形}）在单位元附近的局部结构，由一个\emph{向量空间}——李代数——精确刻画；二者经指数/对数映射局部一一对应。于是：在李代数这一线性空间上自由地加减、求导，再用指数映射映回李群。这把\textbf{带约束优化转化为无约束优化}——正是全书反复使用的范式。

\begin{figure}[ht]
\centering
\begin{tikzpicture}[scale=1.5,>=Stealth]
  \draw[thick, main] (0,0) circle (1);
  \node[main] at (-1.15,0.95) {$\mathcal{M}$ (李群)};
  \coordinate (E) at (0,1);
  \fill (E) circle (0.8pt) node[above right] {$\mathcal{E}$ (单位元)};
  \draw[thick, second] (-1.1,1) -- (1.25,1) node[right] {$T_{\mathcal{E}}\mathcal{M}$ (李代数/切空间)};
  \draw[->, third, thick] (E) .. controls (0.5,0.95) and (0.85,0.6) .. (0.92,0.38);
  \node[third] at (1.2,0.55) {$\exp$};
  \fill[second] (0.6,1) circle (0.7pt) node[above] {$\boldsymbol{\phi}$};
\end{tikzpicture}
\caption{李群流形 $\mathcal{M}$ 与其在单位元 $\mathcal{E}$ 处的切空间（李代数）。切向量 $\boldsymbol{\phi}$ 经指数映射“贴合”到流形上。（仿 \cite{barfoot2024state} 重绘）}
\label{fig:lie-manifold}
\end{figure}

\begin{pitfall}[概念误区：“流形 = 弯曲，所以一定难算”]
新手常以为“活在曲面上就只能用复杂的微分几何”。实际上，流形的\emph{定义性质}恰恰是\textbf{局部像欧氏空间}（\cref{fig:lie-manifold} 的切平面）。李群之所以好用，正因为我们\emph{只在单位元附近的切空间（李代数）里做线性运算}，再映回去——难的部分被指数映射吸收了。理解这一点，本章后面所有“扰动 + 求导”的套路都顺理成章。
\end{pitfall}

% ════════════════════════════════════════════════════════════════
\section{群与矩阵李群}
\label{sec:group}

\paragraph{直觉：加法不封闭，只有乘法封闭。}
为什么旋转“只能是群、不能是向量空间”？一个朴素却关键的观察\cite{gaoxiang2019slam14}：任取两个旋转矩阵 $\mathbf{R}_1,\mathbf{R}_2$，它们的\emph{和}一般不再是旋转矩阵（例如 $\mathbf{R}+(-\mathbf{R})=\mathbf{0}\notin\SO(3)$）：
\begin{equation}
  \mathbf{R}_1+\mathbf{R}_2\notin\SO(3),\qquad \mathbf{T}_1+\mathbf{T}_2\notin\SE(3).
  \label{eq:not-vs}
\end{equation}
但它们的\emph{乘积}仍是旋转/位姿：$\mathbf{R}_1\mathbf{R}_2\in\SO(3)$，$\mathbf{T}_1\mathbf{T}_2\in\SE(3)$，对应“先后做两次变换”的复合。对于这种\emph{只配备一个良好运算（乘法）}的集合，数学上称之为\emph{群}。

\begin{definition}[群]\label{def:group}
群是集合 $A$ 与二元运算 $\cdot$ 构成的代数结构 $G=(A,\cdot)$，满足（助记“封结幺逆”）：
\begin{enumerate}
  \item \textbf{封}闭性：$\forall a_1,a_2\in A,\; a_1\cdot a_2\in A$；
  \item \textbf{结}合律：$\forall a_1,a_2,a_3\in A,\;(a_1\cdot a_2)\cdot a_3=a_1\cdot(a_2\cdot a_3)$；
  \item \textbf{幺}元：$\exists a_0\in A,\;\forall a\in A,\; a_0\cdot a=a\cdot a_0=a$；
  \item 逆：$\forall a\in A,\;\exists a^{-1}\in A,\; a\cdot a^{-1}=a_0$。
\end{enumerate}
\end{definition}

\noindent 常见群：整数加法 $(\mathbb{Z},+)$、非零有理数乘法 $(\mathbb{Q}\setminus\{0\},\cdot)$；矩阵中的一般线性群 $\mathrm{GL}(n)$、特殊正交群 $\SO(n)$、特殊欧氏群 $\SE(n)$。\textbf{李群}指\emph{具有连续（光滑）结构}的群：离散的 $(\mathbb{Z},+)$ 不是李群，而刚体可以在空间中\emph{连续}地运动，故 $\SO(n),\SE(n)$ 都是李群。三维旋转与位姿分别是：
\begin{align}
  \SO(3) &= \big\{\mathbf{R}\in\mathbb{R}^{3\times3}\,\big|\,\mathbf{R}\mathbf{R}^\top=\mathbf{I},\ \det\mathbf{R}=1\big\},\label{eq:so3}\\
  \SE(3) &= \Big\{\mathbf{T}=\begin{bmatrix}\mathbf{R}&\mathbf{t}\\\mathbf{0}^\top&1\end{bmatrix}\in\mathbb{R}^{4\times4}\,\Big|\,\mathbf{R}\in\SO(3),\ \mathbf{t}\in\mathbb{R}^3\Big\}.\label{eq:se3}
\end{align}

\begin{proposition}[$\SO(3),\SE(3)$ 是矩阵李群]\label{prop:matrix-lie}
$\SO(3),\SE(3)$ 关于矩阵乘法满足 \cref{def:group} 四条公理，且是光滑流形，故为矩阵李群；但二者对矩阵加法\emph{不封闭}（\cref{eq:not-vs}），故\emph{不是}向量空间。
\end{proposition}
\begin{proof}
结合律由矩阵乘法的结合律直接给出；幺元为 $\mathbf{I}$。\emph{封闭性}：$\SO(3)$ 中 $(\mathbf{R}_1\mathbf{R}_2)(\mathbf{R}_1\mathbf{R}_2)^\top=\mathbf{R}_1(\mathbf{R}_2\mathbf{R}_2^\top)\mathbf{R}_1^\top=\mathbf{I}$，$\det(\mathbf{R}_1\mathbf{R}_2)=\det\mathbf{R}_1\det\mathbf{R}_2=1$；$\SE(3)$ 中 $\mathbf{T}_1\mathbf{T}_2=\big[\begin{smallmatrix}\mathbf{R}_1\mathbf{R}_2&\mathbf{R}_1\mathbf{t}_2+\mathbf{t}_1\\\mathbf{0}^\top&1\end{smallmatrix}\big]$ 仍是位姿。\emph{逆元}：$\mathbf{R}^{-1}=\mathbf{R}^\top\in\SO(3)$；$\mathbf{T}^{-1}=\big[\begin{smallmatrix}\mathbf{R}^\top&-\mathbf{R}^\top\mathbf{t}\\\mathbf{0}^\top&1\end{smallmatrix}\big]$，验证 $\mathbf{T}\mathbf{T}^{-1}=\big[\begin{smallmatrix}\mathbf{R}\mathbf{R}^\top&-\mathbf{R}\mathbf{R}^\top\mathbf{t}+\mathbf{t}\\\mathbf{0}^\top&1\end{smallmatrix}\big]=\mathbf{I}$。\emph{加法不封闭}：零矩阵 $\mathbf{0}\notin\SO(3)$。
\end{proof}

\begin{note}
\textbf{$\SE(3)$ 的逆有闭式}（三书一致，常用）：
\[
  \mathbf{T}^{-1}=\begin{bmatrix}\mathbf{R}^\top&-\mathbf{R}^\top\mathbf{t}\\\mathbf{0}^\top&1\end{bmatrix},
\]
即先用 $\mathbf{R}^\top$ 反转旋转，再把平移反向并旋转回来。Barfoot 记旋转为 $\mathbf{C}$；SLAM Handbook 多从流形/因子图视角使用同一对象。
\end{note}

\begin{pitfall}[概念误区：把 $\mathbf{R}^{-1}=\mathbf{R}^\top$ 当成巧合]
“旋转矩阵的逆恰好等于转置，真方便”——这\emph{不是}巧合，而是正交矩阵的\textbf{定义性质}：$\mathbf{R}\in\SO(3)\iff\mathbf{R}^\top\mathbf{R}=\mathbf{I}$。几何上，旋转保长度（保范），所以它的逆就是“反着转”，即转置。理解了这点，就知道为什么 $\SE(3)$ 的逆\emph{不是}简单转置（平移项要额外处理）——$\SE(3)$ 不是正交群。
\end{pitfall}

\begin{practice}
\begin{enumerate}
  \item 验证 $\SO(3)$、$\SE(3)$ 关于矩阵乘法构成群（逐条核对“封结幺逆”）。
  \item 由 $\mathbf{T}=\big[\begin{smallmatrix}\mathbf{R}&\mathbf{t}\\\mathbf{0}^\top&1\end{smallmatrix}\big]$ 直接验证上面 $\mathbf{T}^{-1}$ 的闭式（提示：算 $\mathbf{T}\mathbf{T}^{-1}$）。
  \item 举一个反例说明 $\SO(3)$ 对加法不封闭。
\end{enumerate}
\end{practice}

% ════════════════════════════════════════════════════════════════
\section{李代数}
\label{sec:lie-algebra}

\paragraph{动机与招牌推导：从 $\mathbf{R}\mathbf{R}^\top=\mathbf{I}$ 自然“挤出”李代数。}
我们不直接“天降”李代数这个符号，而是顺着\emph{旋转随时间变化}把它逼出来\cite{gaoxiang2019slam14}。设旋转 $\mathbf{R}(t)$ 连续变化，恒满足约束 $\mathbf{R}(t)\mathbf{R}(t)^\top=\mathbf{I}$。两边对时间求导：
\begin{equation}
  \dot{\mathbf{R}}\mathbf{R}^\top+\mathbf{R}\dot{\mathbf{R}}^\top=\mathbf{0}
  \;\Longrightarrow\;
  \dot{\mathbf{R}}\mathbf{R}^\top=-\big(\dot{\mathbf{R}}\mathbf{R}^\top\big)^\top.
  \label{eq:RRt-deriv}
\end{equation}
\cref{eq:RRt-deriv} 说明 $\dot{\mathbf{R}}\mathbf{R}^\top$ 是\textbf{反对称矩阵}。任一反对称矩阵都对应唯一三维向量 $\boldsymbol{\phi}(t)\in\mathbb{R}^3$，记 $\dot{\mathbf{R}}\mathbf{R}^\top=\boldsymbol{\phi}^\wedge$，其中
\begin{equation}
  \boldsymbol{\phi}^\wedge=\begin{bmatrix}0&-\phi_3&\phi_2\\\phi_3&0&-\phi_1\\-\phi_2&\phi_1&0\end{bmatrix},\qquad (\boldsymbol{\phi}^\wedge)^\vee=\boldsymbol{\phi}.
  \label{eq:hat}
\end{equation}
右乘 $\mathbf{R}$（正交，$\mathbf{R}^\top\mathbf{R}=\mathbf{I}$）即得旋转的\emph{微分方程}：
\begin{equation}
  \dot{\mathbf{R}}(t)=\boldsymbol{\phi}(t)^\wedge\,\mathbf{R}(t).
  \label{eq:Rdot}
\end{equation}
直觉是：\emph{每对旋转求一次导，相当于左乘一个反对称矩阵 $\boldsymbol{\phi}^\wedge$}。为看清 $\boldsymbol{\phi}$ 的几何含义，不妨设 $t_0=0$、$\mathbf{R}(0)=\mathbf{I}$，把 $\mathbf{R}(t)$ 在 $t=0$ 处做一阶泰勒展开\cite{gaoxiang2019slam14}：
\begin{equation}
  \mathbf{R}(t)\approx\mathbf{R}(0)+\dot{\mathbf{R}}(0)\,t=\mathbf{I}+\boldsymbol{\phi}(0)^\wedge\,t.
  \label{eq:R-taylor}
\end{equation}
\cref{eq:R-taylor} 把抽象的“切向量”落到了实处：在单位元附近，旋转矩阵\emph{偏离单位阵的那一点点}恰由反对称矩阵 $\boldsymbol{\phi}^\wedge$ 沿时间线性给出。换言之，$\boldsymbol{\phi}$ 刻画了 $\mathbf{R}$ 在单位元的\emph{一阶变化率}，正是 $\mathbf{R}$ 在单位元处的\emph{切向量}，而它生活的那个由反对称矩阵张成的空间，就是 $\SO(3)$ 在单位元的\emph{切空间}（Tangent Space）。这正是“李代数 = 单位元处切空间”这句话的具体由来。若进一步设 $\boldsymbol{\phi}$ 在该处取常值 $\boldsymbol{\phi}_0$，\cref{eq:Rdot} 退化为线性常微分方程，结合初值 $\mathbf{R}(0)=\mathbf{I}$ 解得
\begin{equation}
  \mathbf{R}(t)=\exp(\boldsymbol{\phi}_0^\wedge\, t),
  \label{eq:exp-emerge}
\end{equation}
——旋转矩阵与反对称矩阵\emph{经矩阵指数发生了联系}。这就引出两个待解决的问题：(i) 与 $\mathbf{R}$ 对应的 $\boldsymbol{\phi}$ 是什么？它就是 $\SO(3)$ 的李代数 $\so(3)$ 元素；(ii) 矩阵指数 $\exp(\boldsymbol{\phi}^\wedge)$ 怎么算、能否反算？这就是指数/对数映射（\cref{sec:exp-log}）。

\begin{definition}[李代数]\label{def:lie-algebra}
李代数由向量空间 $\mathbb{V}$、数域 $\mathbb{F}$ 与二元运算（李括号）$[\cdot,\cdot]:\mathbb{V}\times\mathbb{V}\to\mathbb{V}$ 组成，满足：
\begin{enumerate}
  \item 封闭性：$\forall \mathbf{X},\mathbf{Y}\in\mathbb{V},\;[\mathbf{X},\mathbf{Y}]\in\mathbb{V}$；
  \item 双线性：$[a\mathbf{X}+b\mathbf{Y},\mathbf{Z}]=a[\mathbf{X},\mathbf{Z}]+b[\mathbf{Y},\mathbf{Z}]$（对第二变元同理）；
  \item 自反性：$\forall\mathbf{X}\in\mathbb{V},\;[\mathbf{X},\mathbf{X}]=\mathbf{0}$；
  \item 雅可比恒等式：$[\mathbf{X},[\mathbf{Y},\mathbf{Z}]]+[\mathbf{Z},[\mathbf{X},\mathbf{Y}]]+[\mathbf{Y},[\mathbf{Z},\mathbf{X}]]=\mathbf{0}$。
\end{enumerate}
李括号“度量两元素的差异”，不要求结合律。三维向量空间配叉积 $(\mathbb{R}^3,\mathbb{R},\times)$ 就是一个李代数。
\end{definition}

\begin{definition}[$\so(3)$ 与 $\se(3)$]\label{def:so3-se3}
\begin{align}
  \so(3) &= \big\{\boldsymbol{\phi}\in\mathbb{R}^3,\ \boldsymbol{\Phi}=\boldsymbol{\phi}^\wedge\in\mathbb{R}^{3\times3}\big\},\quad
  [\boldsymbol{\phi}_1,\boldsymbol{\phi}_2]=(\boldsymbol{\phi}_1^\wedge\boldsymbol{\phi}_2^\wedge-\boldsymbol{\phi}_2^\wedge\boldsymbol{\phi}_1^\wedge)^\vee;\\
  \se(3) &= \Big\{\boldsymbol{\xi}=\begin{bmatrix}\boldsymbol{\rho}\\\boldsymbol{\phi}\end{bmatrix}\in\mathbb{R}^6,\ \boldsymbol{\rho}\in\mathbb{R}^3,\ \boldsymbol{\phi}\in\so(3),\ \boldsymbol{\xi}^\wedge=\begin{bmatrix}\boldsymbol{\phi}^\wedge&\boldsymbol{\rho}\\\mathbf{0}^\top&0\end{bmatrix}\in\mathbb{R}^{4\times4}\Big\},\\
  &\qquad [\boldsymbol{\xi}_1,\boldsymbol{\xi}_2]=(\boldsymbol{\xi}_1^\wedge\boldsymbol{\xi}_2^\wedge-\boldsymbol{\xi}_2^\wedge\boldsymbol{\xi}_1^\wedge)^\vee\quad(\se(3)\text{ 的李括号，与 }\so(3)\text{ 同构造}).\notag
\end{align}
$\so(3)$ 与 $\SO(3)$ 的关系由 $\mathbf{R}=\exp(\boldsymbol{\phi}^\wedge)$ 给定。$\so(3)$ 恰是“由旋转向量组成、用来表达旋转矩阵导数”的空间。$\se(3)$ 则可朴素地理解为“一个平移分量 $\boldsymbol{\rho}$ 加上一个 $\so(3)$ 元素 $\boldsymbol{\phi}$ 拼成的六维向量”——其李括号与 $\so(3)$ 完全照搬“矩阵对易子再 $\vee$ 回向量”的构造，验证它满足四条公理留作本节习题。
\end{definition}

\begin{note}
此处 $(\cdot)^\wedge$ 被\emph{重载}：在 $\so(3)$ 是反对称化 $\mathbb{R}^3\to\mathbb{R}^{3\times3}$；在 $\se(3)$ 是 $\mathbb{R}^6\to\mathbb{R}^{4\times4}$（右上块直接放 $\boldsymbol{\rho}$，\textbf{不再反对称}）。$\boldsymbol{\rho}$ 是 $\se(3)$ 的“平移分量”，但\textbf{并非}位姿的真实平移 $\mathbf{t}$；二者相差一个左雅可比 $\mathbf{t}=\mathbf{J}\boldsymbol{\rho}$（见 \cref{thm:se3-exp}）。
\end{note}

\begin{pitfall}[概念误区：把“李代数”当成“切空间本身”或把 $\boldsymbol{\phi}^\wedge$ 当成切空间]
两个常见混淆：(1) $\so(3)$ 的\emph{元素}是切向量 $\boldsymbol{\phi}$（或其反对称化 $\boldsymbol{\phi}^\wedge$），而\emph{整个} $\so(3)$ 才是切空间 $T_{\mathcal{E}}\SO(3)$——元素与空间别混。(2) $\boldsymbol{\phi}$ 与 $\boldsymbol{\phi}^\wedge$ 只是同一对象的“向量形式”与“矩阵形式”，靠 $\wedge/\vee$ 互换，做线性代数时常不加区分，但维度要心里有数（$3$ 维向量 vs $3\times3$ 矩阵）。
\end{pitfall}

\begin{practice}
\begin{enumerate}
  \item 验证 $(\mathbb{R}^3,\mathbb{R},\times)$ 满足 \cref{def:lie-algebra} 四条（尤其用 $\mathbf{a}\times\mathbf{a}=\mathbf{0}$ 验自反性、用叉积的雅可比恒等式验第四条）。
  \item 验证 $\so(3)$、$\se(3)$ 在各自李括号下满足李代数公理。
  \item （概念）解释为什么 \cref{eq:RRt-deriv} 一定给出反对称矩阵，而非任意矩阵。
\end{enumerate}
\end{practice}

% ════════════════════════════════════════════════════════════════
\section{指数与对数映射}
\label{sec:exp-log}

李代数到李群的桥是\emph{指数映射} $\mathbf{R}=\exp(\boldsymbol{\phi}^\wedge)$。直接按矩阵指数的级数 $\exp(\mathbf{A})=\sum_{n\ge0}\mathbf{A}^n/n!$ 计算要算无穷次幂，并不现实；幸好 $\so(3)$ 有闭式。

\begin{theorem}[$\SO(3)$ 指数映射：罗德里格斯公式]\label{thm:so3-exp}
设 $\boldsymbol{\phi}=\theta\mathbf{a}$，$\theta=\|\boldsymbol{\phi}\|$，$\mathbf{a}$ 为单位轴，则
\begin{equation}
  \mathbf{R}=\exp(\boldsymbol{\phi}^\wedge)=\cos\theta\,\mathbf{I}+(1-\cos\theta)\,\mathbf{a}\mathbf{a}^\top+\sin\theta\,\mathbf{a}^\wedge.
  \label{eq:rodrigues}
\end{equation}
\end{theorem}

\begin{proof}
关键是用两条单位向量恒等式（$\mathbf{a}$ 为单位向量）把高阶幂降阶：
\begin{equation}
  \mathbf{a}^\wedge\mathbf{a}^\wedge=\mathbf{a}\mathbf{a}^\top-\mathbf{I},\qquad \mathbf{a}^\wedge\mathbf{a}^\wedge\mathbf{a}^\wedge=-\mathbf{a}^\wedge,
  \label{eq:a-identities}
\end{equation}
于是 $\mathbf{a}^\wedge$ 的幂在 $\{\mathbf{a}^\wedge,(\mathbf{a}^\wedge)^2\}$ 间循环。（第二式 $(\mathbf{a}^\wedge)^3=-\mathbf{a}^\wedge$ 并非偶然，它正是 $\so(3)$ 的\emph{最小多项式}，是本章一切闭式得以收敛的代数源头；统一的出处见 \cref{thm:minimal-poly}。）代入指数级数并按奇偶项归并：
\begin{align*}
  \exp(\theta\mathbf{a}^\wedge)
  &=\mathbf{I}+\theta\mathbf{a}^\wedge+\tfrac{\theta^2}{2!}(\mathbf{a}^\wedge)^2+\tfrac{\theta^3}{3!}(\mathbf{a}^\wedge)^3+\tfrac{\theta^4}{4!}(\mathbf{a}^\wedge)^4+\cdots\\
  &=\mathbf{a}\mathbf{a}^\top+\Big(\theta-\tfrac{\theta^3}{3!}+\tfrac{\theta^5}{5!}-\cdots\Big)\mathbf{a}^\wedge
   -\Big(1-\tfrac{\theta^2}{2!}+\tfrac{\theta^4}{4!}-\cdots\Big)(\mathbf{a}^\wedge)^2\\
  &=\mathbf{a}\mathbf{a}^\top+\sin\theta\,\mathbf{a}^\wedge-\cos\theta\,(\mathbf{a}\mathbf{a}^\top-\mathbf{I})
   =\cos\theta\,\mathbf{I}+(1-\cos\theta)\,\mathbf{a}\mathbf{a}^\top+\sin\theta\,\mathbf{a}^\wedge,
\end{align*}
末步用 $(\mathbf{a}^\wedge)^2=\mathbf{a}\mathbf{a}^\top-\mathbf{I}$ 回代，即 \cref{eq:rodrigues}。
\end{proof}
\noindent 注意 \cref{eq:rodrigues} 与 \cref{ch:rigid_body} 的罗德里格斯公式\emph{一模一样}——这正说明 $\so(3)$ 就是旋转向量的空间，指数映射就是罗德里格斯公式。上面的证明是\emph{代数}的（级数按奇偶项归并）；其\emph{几何}诠释见 \cref{fig:rodrigues}。

\begin{figure}[ht]
  \centering
  \includegraphics[width=.52\linewidth]{figures/pdf/rodrigues.pdf}
  \caption{罗德里格斯公式的几何诠释：向量 $\mathbf{v}$ 绕单位轴 $\mathbf{k}$ 转 $\theta$。把 $\mathbf{v}$ 分解为平行分量 $\mathbf{v}_\parallel$（沿轴、旋转不变）与垂直分量 $\mathbf{v}_\perp$；后者在垂直平面内转 $\theta$ 得 $\mathbf{v}_\perp'=\cos\theta\,\mathbf{v}_\perp+\sin\theta\,(\mathbf{k}\times\mathbf{v}_\perp)$，再与 $\mathbf{v}_\parallel$ 叠加即得 $\mathbf{v}_{\mathrm{rot}}$。这正是 \cref{eq:rodrigues} 的向量（几何）形式。}
  \label{fig:rodrigues}
\end{figure}

\paragraph{对数映射（反向）。}与指数映射对称，由 $\mathbf{R}$ 反解 $\boldsymbol{\phi}=\mathrm{Log}(\mathbf{R})$ 也有一个原则性的级数定义\cite{gaoxiang2019slam14}
\begin{equation}
  \boldsymbol{\phi}=\ln(\mathbf{R})^\vee=\Big(\sum_{n=0}^{\infty}\frac{(-1)^n}{n+1}(\mathbf{R}-\mathbf{I})^{n+1}\Big)^{\!\vee},
  \label{eq:so3-log-series}
\end{equation}
但与指数映射一样，我们并不真去算这个无穷级数：用迹与特征向量更省——转角由迹给出，转轴是特征值 $1$ 的特征向量：
\begin{equation}
  \theta=\cos^{-1}\!\frac{\operatorname{tr}(\mathbf{R})-1}{2},\qquad \mathbf{R}\mathbf{a}=\mathbf{a}\ \ (\mathbf{a}\text{ 为特征值 }1\text{ 的单位特征向量}).
  \label{eq:so3-log}
\end{equation}
\cref{eq:so3-log} 即 \cref{ch:rigid_body} 由旋转矩阵反求轴角的公式——再次印证“$\so(3)$ 就是旋转向量空间”。

\begin{proposition}[满射非单射]\label{prop:surjective}
指数映射 $\so(3)\to\SO(3)$ 是\emph{满射但非单射}：$\theta$ 与 $\theta+2\pi m$ 给出同一 $\mathbf{R}$。若限制 $|\theta|<\pi$，则李群与李代数元素一一对应。
\end{proposition}

\begin{pitfall}[思维陷阱：忽视周期性，优化跨越 $\pm\pi$ 处“跳变”]
因为指数映射非单射（\cref{prop:surjective}），$\mathrm{Log}$ 通常取主值 $|\theta|\le\pi$。当真实旋转接近 $\pi$、或优化迭代跨过该边界时，$\mathrm{Log}$ 会发生 $2\pi$ 跳变，残差和雅可比出现不连续。实现里常见后果是 BA/位姿图在大旋转处不收敛。对策：增量更新（扰动）每步都是小量，远离 $\pi$；必要时对相对旋转取 $\mathrm{Log}$ 而非绝对旋转。
\end{pitfall}

\begin{theorem}[$\SE(3)$ 指数映射]\label{thm:se3-exp}
$\mathbf{T}=\exp(\boldsymbol{\xi}^\wedge)=\big[\begin{smallmatrix}\mathbf{R}&\mathbf{J}\boldsymbol{\rho}\\\mathbf{0}^\top&1\end{smallmatrix}\big]$，其中 $\mathbf{R}=\exp(\boldsymbol{\phi}^\wedge)$，\textbf{真实平移} $\mathbf{t}=\mathbf{J}\boldsymbol{\rho}$，左雅可比
\begin{equation}
  \mathbf{J}=\sum_{n=0}^{\infty}\frac{1}{(n+1)!}(\boldsymbol{\phi}^\wedge)^n
  =\frac{\sin\theta}{\theta}\mathbf{I}+\Big(1-\frac{\sin\theta}{\theta}\Big)\mathbf{a}\mathbf{a}^\top+\frac{1-\cos\theta}{\theta}\mathbf{a}^\wedge.
  \label{eq:left-jacobian}
\end{equation}
\end{theorem}
\begin{proof}
$\exp(\boldsymbol{\xi}^\wedge)$ 按 $4\times4$ 分块展开：左上块是 $\mathbf{R}=\exp(\boldsymbol{\phi}^\wedge)$，右上块为 $\big(\sum_{n\ge0}\tfrac{1}{(n+1)!}(\boldsymbol{\phi}^\wedge)^n\big)\boldsymbol{\rho}=\mathbf{J}\boldsymbol{\rho}$。令 $\boldsymbol{\phi}=\theta\mathbf{a}$，用 $(\mathbf{a}^\wedge)^2=\mathbf{a}\mathbf{a}^\top-\mathbf{I}$、$(\mathbf{a}^\wedge)^3=-\mathbf{a}^\wedge$ 归并：
\begin{align*}
  \mathbf{J}&=\mathbf{I}+\tfrac{1}{2!}\theta\mathbf{a}^\wedge+\tfrac{1}{3!}\theta^2(\mathbf{a}^\wedge)^2+\tfrac{1}{4!}\theta^3(\mathbf{a}^\wedge)^3+\cdots\\
  &=\mathbf{I}+\frac{1-\cos\theta}{\theta}\mathbf{a}^\wedge+\frac{\theta-\sin\theta}{\theta}(\mathbf{a}\mathbf{a}^\top-\mathbf{I})
   =\frac{\sin\theta}{\theta}\mathbf{I}+\Big(1-\frac{\sin\theta}{\theta}\Big)\mathbf{a}\mathbf{a}^\top+\frac{1-\cos\theta}{\theta}\mathbf{a}^\wedge,
\end{align*}
即 \cref{eq:left-jacobian}。其逆见 \cref{eq:Jl}，在 $\theta=2\pi m\,(m\neq0)$ 处奇异。
\end{proof}

\begin{insight}
$\mathbf{t}=\mathbf{J}\boldsymbol{\rho}$ 正是“$\se(3)$ 平移分量 $\boldsymbol{\rho}\neq$ 真实平移 $\mathbf{t}$”的根源：只有当 $\boldsymbol{\phi}\to\mathbf{0}$（即 $\mathbf{J}\to\mathbf{I}$）时二者才相等。这个 $\mathbf{J}$ 不是巧合地“顺便出现”，它就是下一节 BCH 里的\emph{左雅可比}，会在扰动、伴随、不确定度里反复登场——记住它。
\end{insight}

\begin{practice}
\begin{enumerate}
  \item 用 \cref{eq:a-identities} 验证 $(\mathbf{a}^\wedge)^4=-(\mathbf{a}^\wedge)^2$，并据此说明为何级数只剩 $\sin/\cos$ 两族。
  \item （推导，在草稿纸上完成）从级数出发独立推出 \cref{eq:rodrigues}，对照正文证明。
  \item 取 $\theta\to0$ 对 \cref{eq:left-jacobian} 做泰勒展开，验证 $\mathbf{J}\to\mathbf{I}+\tfrac12\boldsymbol{\phi}^\wedge+O(\theta^2)$，从而 $\mathbf{t}\approx\boldsymbol{\rho}$。
\end{enumerate}
\end{practice}

% ════════════════════════════════════════════════════════════════
\subsection{雅可比的运动学含义：\texorpdfstring{$\boldsymbol{\omega}=\mathbf{J}\dot{\boldsymbol{\phi}}$}{omega = J phidot}}
\label{subsec:kinematics-jacobian}

\paragraph{动机：把 $\mathbf{J}$ 从“指数映射的副产品”升格为“运动学映射”。}
上一节 $\mathbf{J}$ 是在算 $\SE(3)$ 指数时\emph{顺带}冒出的（$\mathbf{t}=\mathbf{J}\boldsymbol{\rho}$），看起来像个纯代数对象。但它还有一个\emph{物理}身份，这个身份才是它在状态估计里无处不在的真正原因。回忆 \cref{sec:rb-kinematics} 的 Poisson 方程：旋转随时间变化时，机体角速度 $\boldsymbol{\omega}$ 与旋转的导数由 $\dot{\mathbf{R}}=\mathbf{R}\,\boldsymbol{\omega}^\wedge$（机体系）或 $\dot{\mathbf{R}}=\boldsymbol{\omega}^\wedge\mathbf{R}$（世界系）相联。现在换一个问法：如果我们用\emph{李代数坐标} $\boldsymbol{\phi}(t)$（即 $\mathbf{R}(t)=\mathrm{Exp}(\boldsymbol{\phi}(t))$）来描述这条旋转轨迹，那么 $\boldsymbol{\phi}$ 的变化率 $\dot{\boldsymbol{\phi}}$ 与角速度 $\boldsymbol{\omega}$ 是什么关系？答案恰好是\emph{左/右雅可比}——这就把 \cref{thm:se3-exp} 的 $\mathbf{J}$ 与 \cref{sec:rb-kinematics} 的角速度直接焊在一起。

\paragraph{反面：$\dot{\boldsymbol{\phi}}\neq\boldsymbol{\omega}$（轴角率不是角速度）。}
一个极常见的误解是“旋转向量的导数就是角速度”。\emph{只有瞬时}（$\boldsymbol{\phi}=\mathbf{0}$）才成立。一般地，沿固定轴 $\mathbf{a}$ 匀速转动 $\boldsymbol{\phi}(t)=\theta(t)\mathbf{a}$ 时二者恰好相等（见下方推论），但只要转轴在动，$\dot{\boldsymbol{\phi}}$ 与 $\boldsymbol{\omega}$ 就相差一个 $\mathbf{J}$。把二者混为一谈，是连续时间预积分、姿态滤波里系统性偏差的来源之一。

\begin{theorem}[旋转运动学的李代数形式：$\boldsymbol{\omega}=\mathbf{J}\dot{\boldsymbol{\phi}}$]\label{thm:so3-kinematics}
设旋转轨迹 $\mathbf{R}(t)=\mathrm{Exp}(\boldsymbol{\phi}(t))\in\SO(3)$，李代数坐标 $\boldsymbol{\phi}(t)$ 可微。则\emph{李代数坐标率} $\dot{\boldsymbol{\phi}}$ 与角速度的关系为
\begin{equation}
  \boldsymbol{\omega}_w=\mathbf{J}_l(\boldsymbol{\phi})\,\dot{\boldsymbol{\phi}},\qquad
  \boldsymbol{\omega}_b=\mathbf{J}_r(\boldsymbol{\phi})\,\dot{\boldsymbol{\phi}},
  \label{eq:so3-kinematics}
\end{equation}
其中 $\boldsymbol{\omega}_w$ 为\emph{世界系}角速度（$\dot{\mathbf{R}}=\boldsymbol{\omega}_w^\wedge\mathbf{R}$）、$\boldsymbol{\omega}_b$ 为\emph{机体系}角速度（$\dot{\mathbf{R}}=\mathbf{R}\,\boldsymbol{\omega}_b^\wedge$），$\mathbf{J}_l,\mathbf{J}_r$ 是 \cref{eq:Jl,eq:Jr} 的左/右雅可比。其逆（本书右扰动主线最常用）为
\begin{equation}
  \boxed{\;\dot{\boldsymbol{\phi}}=\mathbf{J}_r(\boldsymbol{\phi})^{-1}\,\boldsymbol{\omega}_b\;}\qquad\big(\text{亦}\ \dot{\boldsymbol{\phi}}=\mathbf{J}_l(\boldsymbol{\phi})^{-1}\,\boldsymbol{\omega}_w\big).
  \label{eq:so3-kinematics-inv}
\end{equation}
\end{theorem}

\begin{proof}
对 $\mathbf{R}=\mathrm{Exp}(\boldsymbol{\phi})=\exp(\boldsymbol{\phi}^\wedge)$ 求时间导数。矩阵指数的导数\emph{不能}简单写成 $\dot{\boldsymbol{\phi}}^\wedge\exp(\boldsymbol{\phi}^\wedge)$（因 $\boldsymbol{\phi}(t)$ 与 $\dot{\boldsymbol{\phi}}(t)$ 一般不交换），正确的公式是\emph{含参矩阵指数的导数}\cite{barfoot2024state}
\begin{equation}
  \frac{d}{dt}\exp(\mathbf{A}(t))=\int_{0}^{1}\exp(\alpha\mathbf{A})\,\dot{\mathbf{A}}\,\exp\big((1-\alpha)\mathbf{A}\big)\,d\alpha,
  \label{eq:matexp-deriv}
\end{equation}
（它本身可由 $\exp$ 的级数逐项求导、再对各项里 $\dot{\mathbf{A}}$ 的位置积分得到。）取 $\mathbf{A}=\boldsymbol{\phi}^\wedge$、$\dot{\mathbf{A}}=\dot{\boldsymbol{\phi}}^\wedge$，并右乘 $\mathbf{R}^\top=\exp(-\boldsymbol{\phi}^\wedge)$：
\begin{align}
  \dot{\mathbf{R}}\mathbf{R}^\top
  &=\int_0^1\exp(\alpha\boldsymbol{\phi}^\wedge)\,\dot{\boldsymbol{\phi}}^\wedge\,\exp(-\alpha\boldsymbol{\phi}^\wedge)\,d\alpha
   =\int_0^1\mathbf{R}^\alpha\,\dot{\boldsymbol{\phi}}^\wedge\,\mathbf{R}^{-\alpha}\,d\alpha\notag\\
  &=\int_0^1\big(\mathbf{R}^\alpha\dot{\boldsymbol{\phi}}\big)^\wedge\,d\alpha
   =\Big(\underbrace{\Big(\int_0^1\mathbf{R}^\alpha\,d\alpha\Big)}_{=\,\mathbf{J}_l}\dot{\boldsymbol{\phi}}\Big)^\wedge
   =\big(\mathbf{J}_l\dot{\boldsymbol{\phi}}\big)^\wedge,
  \label{eq:kinematics-Jl-integral}
\end{align}
第三步用了伴随恒等式 $\mathbf{R}^\alpha\mathbf{v}^\wedge\mathbf{R}^{-\alpha}=(\mathbf{R}^\alpha\mathbf{v})^\wedge$（\cref{sec:adjoint}，因 $\mathbf{R}^\alpha\in\SO(3)$），第四步用了左雅可比的\emph{积分表示}
\begin{equation}
  \mathbf{J}_l(\boldsymbol{\phi})=\int_0^1\mathbf{R}^\alpha\,d\alpha=\int_0^1\exp(\alpha\boldsymbol{\phi}^\wedge)\,d\alpha,
  \label{eq:Jl-integral}
\end{equation}
此式与级数式 \cref{eq:left-jacobian} 等价：逐项积分 $\int_0^1\alpha^n d\alpha=\tfrac{1}{n+1}$ 恰好把 $\exp$ 的级数 $\sum_n\tfrac{(\boldsymbol{\phi}^\wedge)^n}{n!}\alpha^n$ 变成 $\sum_n\tfrac{(\boldsymbol{\phi}^\wedge)^n}{(n+1)!}=\mathbf{J}_l$。对照 Poisson 方程 $\dot{\mathbf{R}}\mathbf{R}^\top=\boldsymbol{\omega}_w^\wedge$ 即得 $\boldsymbol{\omega}_w=\mathbf{J}_l\dot{\boldsymbol{\phi}}$。\par
机体系版本同理：把 $\mathbf{R}^\top$ 改为\emph{左}乘，$\mathbf{R}^\top\dot{\mathbf{R}}=\int_0^1\mathbf{R}^{-(1-\alpha)}\dot{\boldsymbol{\phi}}^\wedge\mathbf{R}^{1-\alpha}d\alpha$，换元 $\beta=1-\alpha$ 得 $\big((\int_0^1\mathbf{R}^{-\beta}d\beta)\dot{\boldsymbol{\phi}}\big)^\wedge=(\mathbf{J}_l(-\boldsymbol{\phi})\dot{\boldsymbol{\phi}})^\wedge=(\mathbf{J}_r\dot{\boldsymbol{\phi}})^\wedge$，用 $\mathbf{J}_r(\boldsymbol{\phi})=\mathbf{J}_l(-\boldsymbol{\phi})$（\cref{eq:Jr}）。对照 $\mathbf{R}^\top\dot{\mathbf{R}}=\boldsymbol{\omega}_b^\wedge$ 即得 $\boldsymbol{\omega}_b=\mathbf{J}_r\dot{\boldsymbol{\phi}}$。两式在 $\mathbf{J}$ 可逆区（$|\boldsymbol{\phi}|<\pi$）取逆即 \cref{eq:so3-kinematics-inv}。\par
两者也可由伴随直接互验：$\boldsymbol{\omega}_w=\mathrm{Ad}(\mathbf{R})\boldsymbol{\omega}_b=\mathbf{R}\boldsymbol{\omega}_b$，配合 $\mathbf{J}_l=\mathbf{R}\,\mathbf{J}_r$（同一积分两种写法，亦见 \cref{eq:Jl} 与 $\mathrm{Ad}$ 的关系）即自洽。
\end{proof}

\begin{corollary}[匀轴转动时 $\dot{\boldsymbol{\phi}}=\boldsymbol{\omega}$]\label{cor:fixed-axis}
若转轴固定 $\boldsymbol{\phi}(t)=\theta(t)\mathbf{a}$（$\mathbf{a}$ 常向量），则 $\dot{\boldsymbol{\phi}}=\dot{\theta}\mathbf{a}$ 与 $\mathbf{a}$ 共线。而 $\mathbf{J}_l,\mathbf{J}_r$ 作用在\emph{沿轴}向量上均为恒等（由 \cref{eq:Jl,eq:Jr}：$\mathbf{J}\mathbf{a}=\mathbf{a}$，因 $\mathbf{a}^\wedge\mathbf{a}=\mathbf{0}$、$\mathbf{a}\mathbf{a}^\top\mathbf{a}=\mathbf{a}$ 使三项系数和为 $1$）。故 $\boldsymbol{\omega}_w=\boldsymbol{\omega}_b=\dot{\theta}\mathbf{a}=\dot{\boldsymbol{\phi}}$——此时（也\emph{仅}此时）轴角率等于角速度。这解释了为何“$\dot{\boldsymbol{\phi}}=\boldsymbol{\omega}$”在教科书的单轴例子里总成立，却在一般三维转动里失效。
\end{corollary}

\begin{theorem}[位姿运动学的李代数形式：$\boldsymbol{\varpi}=\boldsymbol{\mathcal{J}}\dot{\boldsymbol{\xi}}$]\label{thm:se3-kinematics}
设位姿轨迹 $\mathbf{T}(t)=\mathrm{Exp}(\boldsymbol{\xi}(t))\in\SE(3)$，则广义速度 $\boldsymbol{\varpi}=[\boldsymbol{\nu};\boldsymbol{\omega}]$（\cref{eq:rb-se3-kin}）与李代数坐标率的关系为
\begin{equation}
  \boldsymbol{\varpi}_w=\boldsymbol{\mathcal{J}}_l(\boldsymbol{\xi})\,\dot{\boldsymbol{\xi}},\qquad
  \boldsymbol{\varpi}_b=\boldsymbol{\mathcal{J}}_r(\boldsymbol{\xi})\,\dot{\boldsymbol{\xi}},\qquad
  \dot{\boldsymbol{\xi}}=\boldsymbol{\mathcal{J}}_r(\boldsymbol{\xi})^{-1}\boldsymbol{\varpi}_b,
  \label{eq:se3-kinematics}
\end{equation}
其中 $\boldsymbol{\varpi}_w^\wedge=\dot{\mathbf{T}}\mathbf{T}^{-1}$（世界系）、$\boldsymbol{\varpi}_b^\wedge=\mathbf{T}^{-1}\dot{\mathbf{T}}$（机体系），$\boldsymbol{\mathcal{J}}_l,\boldsymbol{\mathcal{J}}_r$ 是 \cref{eq:se3-6x6} 的 $6\times6$ 雅可比。
\end{theorem}
\begin{proof}
与 \cref{thm:so3-kinematics} 逐字平行，只是把 $\boldsymbol{\phi}^\wedge$ 换成 $\boldsymbol{\xi}^\wedge$、把伴随恒等式换成 $\mathbf{T}^\alpha\boldsymbol{\xi}^\wedge\mathbf{T}^{-\alpha}=(\boldsymbol{\mathcal{T}}^\alpha\dot{\boldsymbol{\xi}})^\wedge$（$\boldsymbol{\mathcal{T}}=\mathrm{Ad}(\mathbf{T})$）、把积分表示换成 $\boldsymbol{\mathcal{J}}_l=\int_0^1\boldsymbol{\mathcal{T}}^\alpha d\alpha=\sum_{n\ge0}\tfrac{(\boldsymbol{\xi}^\curlywedge)^n}{(n+1)!}$（\cref{eq:se3-6x6}）。对照 $\dot{\mathbf{T}}\mathbf{T}^{-1}=\boldsymbol{\varpi}_w^\wedge$ 即得世界系式，机体系式由换元同理。
\end{proof}

\begin{insight}[这就是连续旋转预积分所引用的前置定理]
\cref{thm:so3-kinematics} 的逆式 \cref{eq:so3-kinematics-inv}，$\dot{\boldsymbol{\phi}}=\mathbf{J}_r^{-1}\boldsymbol{\omega}_b$，正是 VIO 与 IMU 预积分章里\emph{连续旋转预积分}所依赖的旋转向量动力学（彼处把旋转向量记作 $\mathbf{r}$，写成 $\dot{\mathbf{r}}=\mathbf{J}_r(\mathbf{r})^{-1}\boldsymbol{\omega}_b^b$）：陀螺直接测的是\emph{机体系}角速度 $\boldsymbol{\omega}_b$，而要在\emph{线性}空间 $\boldsymbol{\phi}$（或 $\mathbf{r}$）里做高斯过程插值/解析积分，就必须先用 $\mathbf{J}_r^{-1}$ 把 $\boldsymbol{\omega}_b$ 转成李代数坐标率 $\dot{\boldsymbol{\phi}}$。这条式子也是离散预积分里一阶近似 $\mathrm{Exp}(\boldsymbol{\phi}+\delta\boldsymbol{\phi})\approx\mathrm{Exp}(\boldsymbol{\phi})\mathrm{Exp}(\mathbf{J}_r\delta\boldsymbol{\phi})$（见 \cref{sec:bch}）的连续时间对应：把“小增量 $\delta\boldsymbol{\phi}$”换成“无穷小时间步 $\dot{\boldsymbol{\phi}}\,dt$”即得。\textbf{至此 $\mathbf{J}$ 在本书有了双重身份：指数映射里的代数因子（$\mathbf{t}=\mathbf{J}\boldsymbol{\rho}$）与运动学里的速度映射（$\boldsymbol{\omega}=\mathbf{J}\dot{\boldsymbol{\phi}}$）——同一个 $\mathbf{J}$，两处现身绝非巧合，根源都是 \cref{eq:Jl-integral} 的积分 $\int_0^1\mathbf{R}^\alpha d\alpha$。}
\end{insight}

\begin{pitfall}[陷阱：世界系/机体系角速度与左/右雅可比张冠李戴]
\cref{eq:so3-kinematics} 有四种组合，配错会引入难查的符号/方向错误。铁律：\textbf{机体系角速度 $\boldsymbol{\omega}_b$ 配右雅可比 $\mathbf{J}_r$，世界系角速度 $\boldsymbol{\omega}_w$ 配左雅可比 $\mathbf{J}_l$}（与“右扰动在机体系、左扰动在世界系”的图像 \cref{eq:right-pert-so3} 一脉相承）。陀螺测的是机体系 $\boldsymbol{\omega}_b$，故工程里几乎总用 $\dot{\boldsymbol{\phi}}=\mathbf{J}_r^{-1}\boldsymbol{\omega}_b$。Barfoot 以\emph{左}扰动为主线，故其正文写 $\boldsymbol{\omega}=\mathbf{J}_l\dot{\boldsymbol{\phi}}$（世界系）\cite{barfoot2024state}；本书右扰动主线下机体系版 $\boldsymbol{\omega}_b=\mathbf{J}_r\dot{\boldsymbol{\phi}}$ 才是常用形态，两者经 $\boldsymbol{\omega}_w=\mathbf{R}\boldsymbol{\omega}_b$、$\mathbf{J}_l=\mathbf{R}\mathbf{J}_r$ 精确互换。
\end{pitfall}

\begin{practice}
\begin{enumerate}
  \item 由 \cref{eq:Jl,eq:Jr} 直接验证 $\mathbf{J}_l\mathbf{a}=\mathbf{J}_r\mathbf{a}=\mathbf{a}$（沿轴方向），从而印证 \cref{cor:fixed-axis}。
  \item 用 \cref{eq:matexp-deriv} 推 \cref{eq:kinematics-Jl-integral}（提示：$\mathbf{R}^\alpha\dot{\boldsymbol{\phi}}^\wedge\mathbf{R}^{-\alpha}$ 用伴随化成 $(\mathbf{R}^\alpha\dot{\boldsymbol{\phi}})^\wedge$）。
  \item （连接 \cref{sec:bch}）把 \cref{eq:so3-kinematics-inv} 在一个小时间步 $\Delta t$ 上做一阶离散：$\boldsymbol{\phi}(t+\Delta t)\approx\boldsymbol{\phi}+\mathbf{J}_r^{-1}\boldsymbol{\omega}_b\Delta t$，并说明它如何对应 IMU 预积分里的 $\mathrm{Exp}(\boldsymbol{\phi})\mathrm{Exp}(\mathbf{J}_r(\boldsymbol{\phi})\delta\boldsymbol{\phi})$ 形式。
\end{enumerate}
\end{practice}

% ════════════════════════════════════════════════════════════════
\section{BCH 公式与左右雅可比}
\label{sec:bch}

\paragraph{动机：李代数上的“加法”对应李群上的什么？}
优化要在李代数里做加减，但李群里是\emph{乘法}。于是要问：$\exp(\boldsymbol{\phi}_1^\wedge)\exp(\boldsymbol{\phi}_2^\wedge)$ 是否等于 $\exp((\boldsymbol{\phi}_1+\boldsymbol{\phi}_2)^\wedge)$？对\emph{标量}成立，但对\emph{矩阵}一般\textbf{不成立}（因为矩阵不交换）。差多少，由 Baker–Campbell–Hausdorff（BCH）公式给出：
\begin{equation}
  \ln\big(\exp(\mathbf{A})\exp(\mathbf{B})\big)=\mathbf{A}+\mathbf{B}+\tfrac12[\mathbf{A},\mathbf{B}]+\tfrac1{12}[\mathbf{A},[\mathbf{A},\mathbf{B}]]-\tfrac1{12}[\mathbf{B},[\mathbf{A},\mathbf{B}]]+\cdots
  \label{eq:bch}
\end{equation}
即乘积的对数 = 两者之和 + 一串\emph{李括号}余项。

\paragraph{线性近似——本节真正要用的。}
当其中一个为小量时，二阶以上李括号可略，BCH 退化为线性形式，并\emph{自然分出左、右两种}\cite{gaoxiang2019slam14}：
\begin{equation}
  \ln\!\big(\exp(\boldsymbol{\phi}_1^\wedge)\exp(\boldsymbol{\phi}_2^\wedge)\big)^\vee\approx
  \begin{cases}
    \mathbf{J}_l(\boldsymbol{\phi}_2)^{-1}\,\boldsymbol{\phi}_1+\boldsymbol{\phi}_2, & \boldsymbol{\phi}_1\to\mathbf{0}\ (\text{左乘小量}),\\[2pt]
    \mathbf{J}_r(\boldsymbol{\phi}_1)^{-1}\,\boldsymbol{\phi}_2+\boldsymbol{\phi}_1, & \boldsymbol{\phi}_2\to\mathbf{0}\ (\text{右乘小量}).
  \end{cases}
  \label{eq:bch-linear}
\end{equation}
解读：给 $\mathbf{R}_2$（李代数 $\boldsymbol{\phi}_2$）\emph{左乘}一个微小旋转（李代数 $\boldsymbol{\phi}_1$），在李代数上约等于在 $\boldsymbol{\phi}_2$ 上加 $\mathbf{J}_l^{-1}\boldsymbol{\phi}_1$；\emph{右乘}小量则加 $\mathbf{J}_r^{-1}$ 项。\textbf{用左乘模型还是右乘模型，须自始至终一致}——这正对应本书选定的\emph{右扰动}（\cref{sec:perturb}）。

\paragraph{左/右雅可比的闭式。}左雅可比 $\mathbf{J}_l$ 即 \cref{eq:left-jacobian} 的 $\mathbf{J}$；其逆与右版本为（$\boldsymbol{\phi}=\theta\mathbf{a}$）：
\begin{align}
  \mathbf{J}_l &= \frac{\sin\theta}{\theta}\mathbf{I}+\Big(1-\frac{\sin\theta}{\theta}\Big)\mathbf{a}\mathbf{a}^\top+\frac{1-\cos\theta}{\theta}\mathbf{a}^\wedge,\label{eq:Jl}\\
  \mathbf{J}_l^{-1} &= \frac{\theta}{2}\cot\frac{\theta}{2}\,\mathbf{I}+\Big(1-\frac{\theta}{2}\cot\frac{\theta}{2}\Big)\mathbf{a}\mathbf{a}^\top-\frac{\theta}{2}\mathbf{a}^\wedge,\\
  \mathbf{J}_r(\boldsymbol{\phi}) &= \mathbf{J}_l(-\boldsymbol{\phi})=\frac{\sin\theta}{\theta}\mathbf{I}+\Big(1-\frac{\sin\theta}{\theta}\Big)\mathbf{a}\mathbf{a}^\top-\frac{1-\cos\theta}{\theta}\mathbf{a}^\wedge,\label{eq:Jr}\\
  \mathbf{J}_r^{-1} &= \frac{\theta}{2}\cot\frac{\theta}{2}\,\mathbf{I}+\Big(1-\frac{\theta}{2}\cot\frac{\theta}{2}\Big)\mathbf{a}\mathbf{a}^\top+\frac{\theta}{2}\mathbf{a}^\wedge.
\end{align}
关键关系（贯穿全章）：$\boxed{\,\mathbf{J}_r(\boldsymbol{\phi})=\mathbf{J}_l(-\boldsymbol{\phi})=\mathbf{J}_l(\boldsymbol{\phi})^\top\,}$（最后一个等号因 $\mathbf{a}\mathbf{a}^\top$、$\mathbf{I}$ 对称而 $\mathbf{a}^\wedge$ 反对称）。当 $\theta\to0$，$\mathbf{J}_l,\mathbf{J}_r\to\mathbf{I}$，左右扰动趋于一致。

\paragraph{把 BCH 翻译成“群乘 $\leftrightarrow$ 代数加”的可操作规则。}\cref{eq:bch-linear} 写成对数形式不够直观；下面把它\emph{重新叙述}成实际优化里会反复套用的两条对照规则\cite{gaoxiang2019slam14}，读者务必把它们当成“词典”记牢。设旋转 $\mathbf{R}$ 的李代数为 $\boldsymbol{\phi}$，给它\emph{左乘}一个微小旋转 $\Delta\mathbf{R}=\mathrm{Exp}(\Delta\boldsymbol{\phi})$。在李群上结果是 $\Delta\mathbf{R}\cdot\mathbf{R}$，在李代数上则等价于在 $\boldsymbol{\phi}$ 上加一项 $\mathbf{J}_l^{-1}(\boldsymbol{\phi})\Delta\boldsymbol{\phi}$，合并为
\begin{equation}
  \exp(\Delta\boldsymbol{\phi}^\wedge)\exp(\boldsymbol{\phi}^\wedge)=\exp\!\big((\boldsymbol{\phi}+\mathbf{J}_l^{-1}(\boldsymbol{\phi})\Delta\boldsymbol{\phi})^\wedge\big).
  \label{eq:bch-group-to-alg}
\end{equation}
反过来，若我们\emph{先}在李代数上做加法 $\boldsymbol{\phi}+\Delta\boldsymbol{\phi}$，再映回群，则它\emph{不}等于朴素的群乘，而要带上左/右雅可比：
\begin{equation}
  \exp\!\big((\boldsymbol{\phi}+\Delta\boldsymbol{\phi})^\wedge\big)=\exp\!\big((\mathbf{J}_l\Delta\boldsymbol{\phi})^\wedge\big)\exp(\boldsymbol{\phi}^\wedge)=\exp(\boldsymbol{\phi}^\wedge)\exp\!\big((\mathbf{J}_r\Delta\boldsymbol{\phi})^\wedge\big).
  \label{eq:bch-alg-to-group}
\end{equation}
一句话记忆：\textbf{“代数加一点”到“群上乘一点”之间，差一个雅可比 $\mathbf{J}$；放左边用 $\mathbf{J}_l$，放右边用 $\mathbf{J}_r$。} \cref{eq:bch-group-to-alg,eq:bch-alg-to-group} 正是后面在李代数上做微积分（\cref{sec:perturb} 的导数模型）所依赖的桥梁。

\paragraph{$\SE(3)$ 的 BCH。}对位姿同样成立，只是 $3\times3$ 的 $\mathbf{J}_l,\mathbf{J}_r$ 换成 $6\times6$ 的 $\boldsymbol{\mathcal{J}}_l,\boldsymbol{\mathcal{J}}_r$：
\begin{align}
  \exp(\Delta\boldsymbol{\xi}^\wedge)\exp(\boldsymbol{\xi}^\wedge)&\approx\exp\!\big((\boldsymbol{\mathcal{J}}_l^{-1}\Delta\boldsymbol{\xi}+\boldsymbol{\xi})^\wedge\big)\quad(\text{左乘小量}),\label{eq:bch-se3-left}\\
  \exp(\boldsymbol{\xi}^\wedge)\exp(\Delta\boldsymbol{\xi}^\wedge)&\approx\exp\!\big((\boldsymbol{\mathcal{J}}_r^{-1}\Delta\boldsymbol{\xi}+\boldsymbol{\xi})^\wedge\big)\quad(\text{右乘小量}).\label{eq:bch-se3-right}
\end{align}
$\SE(3)$ 的 $6\times6$ 雅可比 $\boldsymbol{\mathcal{J}}_l$ 形式较复杂，部分教材（如视觉SLAM十四讲\cite{gaoxiang2019slam14}）因计算中用不到而略去；本书出于自包含，在下面把它的闭式补全。

\paragraph{$\SE(3)$ 的 $6\times6$ 雅可比（Barfoot 闭式）。}
这里补上 $\boldsymbol{\mathcal{J}}_l$ 的闭式\cite{barfoot2024state}。记 $\boldsymbol{\mathcal{J}}_l(\boldsymbol{\xi})=\sum_{n\ge0}\frac{1}{(n+1)!}(\boldsymbol{\xi}^\curlywedge)^n$（$\boldsymbol{\xi}^\curlywedge$ 见 \cref{sec:adjoint}），其分块为
\begin{equation}
  \boldsymbol{\mathcal{J}}_l(\boldsymbol{\xi})=\begin{bmatrix}\mathbf{J}_l&\mathbf{Q}_l\\\mathbf{0}&\mathbf{J}_l\end{bmatrix},\qquad
  \boldsymbol{\mathcal{J}}_l^{-1}=\begin{bmatrix}\mathbf{J}_l^{-1}&-\mathbf{J}_l^{-1}\mathbf{Q}_l\mathbf{J}_l^{-1}\\\mathbf{0}&\mathbf{J}_l^{-1}\end{bmatrix},
  \label{eq:se3-6x6}
\end{equation}
其中 $\mathbf{J}_l$ 是 \cref{eq:Jl} 的 $3\times3$ 左雅可比，$\mathbf{Q}_l(\boldsymbol{\rho},\boldsymbol{\phi})$ 的完整闭式（Barfoot 式 8.91）见 \cref{sec:lie-appendix}。右版本 $\boldsymbol{\mathcal{J}}_r(\boldsymbol{\xi})=\boldsymbol{\mathcal{J}}_l(-\boldsymbol{\xi})$。一般 SLAM 旋转较小时常用近似 $\boldsymbol{\mathcal{J}}\approx\mathbf{I}$；需要高精度协方差传播时才用完整式。

\begin{pitfall}[思维陷阱：把 BCH 当成精确等式滥用]
\cref{eq:bch-linear} 是\emph{一阶近似}，只在扰动\emph{小}时成立。有人把它当精确公式用于大角度，会引入系统误差。正确用法：扰动量每步都保持小（这正是迭代优化的天然条件）；要精确合并两个有限旋转，请用群乘法 $\mathrm{Exp}(\boldsymbol{\phi}_1)\mathrm{Exp}(\boldsymbol{\phi}_2)$ 本身，或伴随恒等式（\cref{sec:adjoint}），而非 BCH 截断。
\end{pitfall}

\begin{practice}
\begin{enumerate}
  \item 由 $\mathbf{J}_r(\boldsymbol{\phi})=\mathbf{J}_l(-\boldsymbol{\phi})$ 推出 \cref{eq:Jr}，并验证 $\mathbf{J}_r=\mathbf{J}_l^\top$。
  \item （开放思考）当 $\theta\to0$ 时 $\tfrac{\theta}{2}\cot\tfrac{\theta}{2}\to1$，验证 $\mathbf{J}_l^{-1}\to\mathbf{I}-\tfrac12\boldsymbol{\phi}^\wedge$；解释这与“小扰动下左右一致”如何吻合。
\end{enumerate}
\end{practice}

% ════════════════════════════════════════════════════════════════
\section{扰动模型（右扰动为主）}
\label{sec:perturb}

要解 \cref{eq:pose-opt}，须计算误差对位姿的导数，例如旋转后点 $\mathbf{R}\mathbf{p}$ 对旋转的导数。$\SO(3)$ 无加法，差商定义失效，有两条出路\cite{gaoxiang2019slam14}：

\begin{enumerate}
  \item \textbf{李代数求导（导数模型）}：用 $\boldsymbol{\phi}$ 表示 $\mathbf{R}=\exp(\boldsymbol{\phi}^\wedge)$，对 $\boldsymbol{\phi}$ 求导。
  \item \textbf{扰动模型}：对 $\mathbf{R}$ 左乘或右乘一个微小扰动 $\mathrm{Exp}(\delta\boldsymbol{\phi})$，对扰动 $\delta\boldsymbol{\phi}$ 求导。
\end{enumerate}

\paragraph{导数模型为何不实用。}对 $\partial(\exp(\boldsymbol{\phi}^\wedge)\mathbf{p})/\partial\boldsymbol{\phi}$，因 $\SO(3)$ 无加法，对李代数加扰动 $\delta\boldsymbol{\phi}$ 并用 BCH 的左乘形式 $\exp((\boldsymbol{\phi}+\delta\boldsymbol{\phi})^\wedge)\approx\exp((\mathbf{J}_l\delta\boldsymbol{\phi})^\wedge)\exp(\boldsymbol{\phi}^\wedge)$：
\begin{align}
  \frac{\partial(\mathbf{R}\mathbf{p})}{\partial\boldsymbol{\phi}}
  &=\lim_{\delta\boldsymbol{\phi}\to\mathbf{0}}\frac{\exp((\boldsymbol{\phi}+\delta\boldsymbol{\phi})^\wedge)\mathbf{p}-\exp(\boldsymbol{\phi}^\wedge)\mathbf{p}}{\delta\boldsymbol{\phi}}
   =\lim_{\delta\boldsymbol{\phi}\to\mathbf{0}}\frac{\big(\mathbf{I}+(\mathbf{J}_l\delta\boldsymbol{\phi})^\wedge\big)\mathbf{R}\mathbf{p}-\mathbf{R}\mathbf{p}}{\delta\boldsymbol{\phi}}\notag\\
  &=\lim_{\delta\boldsymbol{\phi}\to\mathbf{0}}\frac{(\mathbf{J}_l\delta\boldsymbol{\phi})^\wedge\mathbf{R}\mathbf{p}}{\delta\boldsymbol{\phi}}
   =\lim_{\delta\boldsymbol{\phi}\to\mathbf{0}}\frac{-(\mathbf{R}\mathbf{p})^\wedge\mathbf{J}_l\delta\boldsymbol{\phi}}{\delta\boldsymbol{\phi}}
   =-(\mathbf{R}\mathbf{p})^\wedge\,\mathbf{J}_l,
  \label{eq:deriv-model}
\end{align}
第三步用 $(\mathbf{J}_l\delta\boldsymbol{\phi})^\wedge\mathbf{R}\mathbf{p}=-(\mathbf{R}\mathbf{p})^\wedge\mathbf{J}_l\delta\boldsymbol{\phi}$（叉积反交换）。结果里\emph{残留一个左雅可比 $\mathbf{J}_l$}——形式复杂、还要算 $\mathbf{J}_l$，不划算。扰动模型则把它消掉。

\paragraph{右扰动模型（本书主线）。}给 $\mathbf{R}$ \textbf{右乘}小扰动 $\mathrm{Exp}(\boldsymbol{\varphi})$（$\boldsymbol{\varphi}\to\mathbf{0}$，李代数记 $\boldsymbol{\varphi}$ 以区别于位姿的 $\boldsymbol{\phi}$），对 $\boldsymbol{\varphi}$ 求导：
\begin{align}
  \frac{\partial(\mathbf{R}\mathbf{p})}{\partial\boldsymbol{\varphi}}
  &=\lim_{\boldsymbol{\varphi}\to\mathbf{0}}\frac{\mathbf{R}\,\mathrm{Exp}(\boldsymbol{\varphi})\,\mathbf{p}-\mathbf{R}\mathbf{p}}{\boldsymbol{\varphi}}
   =\lim_{\boldsymbol{\varphi}\to\mathbf{0}}\frac{\mathbf{R}(\mathbf{I}+\boldsymbol{\varphi}^\wedge)\mathbf{p}-\mathbf{R}\mathbf{p}}{\boldsymbol{\varphi}}\notag\\
  &=\lim_{\boldsymbol{\varphi}\to\mathbf{0}}\frac{\mathbf{R}\,\boldsymbol{\varphi}^\wedge\mathbf{p}}{\boldsymbol{\varphi}}
   =\lim_{\boldsymbol{\varphi}\to\mathbf{0}}\frac{-\mathbf{R}\,\mathbf{p}^\wedge\boldsymbol{\varphi}}{\boldsymbol{\varphi}}
   =\boxed{-\mathbf{R}\,\mathbf{p}^\wedge}.
  \label{eq:right-pert-so3}
\end{align}
其中第三步用一阶泰勒 $\mathrm{Exp}(\boldsymbol{\varphi})\approx\mathbf{I}+\boldsymbol{\varphi}^\wedge$，第四步用 $\boldsymbol{\varphi}^\wedge\mathbf{p}=-\mathbf{p}^\wedge\boldsymbol{\varphi}$（叉积反交换）。结果 \cref{eq:right-pert-so3} \textbf{不含 $\mathbf{J}_l$}，干净——这就是扰动模型的优势。

\paragraph{$\SE(3)$ 右扰动。}对位姿作用齐次点 $\tilde{\mathbf{p}}=[\mathbf{p}^\top,1]^\top$，给 $\mathbf{T}$ 右乘 $\mathrm{Exp}(\delta\boldsymbol{\xi})$（$\delta\boldsymbol{\xi}=[\delta\boldsymbol{\rho};\delta\boldsymbol{\phi}]$）：
\begin{equation}
  \frac{\partial(\mathbf{T}\tilde{\mathbf{p}})}{\partial\delta\boldsymbol{\xi}}
  =\mathbf{T}\,\tilde{\mathbf{p}}^{\,\odot}
  =\begin{bmatrix}\mathbf{R}&-\mathbf{R}\,\mathbf{p}^\wedge\\\mathbf{0}^\top&\mathbf{0}^\top\end{bmatrix},
  \qquad
  \tilde{\mathbf{p}}^{\,\odot}=\begin{bmatrix}\mathbf{p}\\1\end{bmatrix}^{\!\odot}=\begin{bmatrix}\mathbf{I}&-\mathbf{p}^\wedge\\\mathbf{0}^\top&\mathbf{0}^\top\end{bmatrix},
  \label{eq:right-pert-se3}
\end{equation}
其中 $(\cdot)^\odot$ 是把齐次点变成 $4\times6$ 矩阵的算子（Handbook\cite{carlone2026handbook} 记号），满足恒等式 $\boldsymbol{\xi}^\wedge\tilde{\mathbf{p}}=\tilde{\mathbf{p}}^{\,\odot}\boldsymbol{\xi}$（该算子的完整代数——含与之配对的 $6\times4$ 算子 $(\cdot)^\circledcirc$ 及把 $\odot$ 搬过位姿的伴随恒等式 $(\mathbf{T}\mathbf{p})^\odot=\mathbf{T}\mathbf{p}^\odot\boldsymbol{\mathcal{T}}^{-1}$——见 \cref{subsec:odot-algebra}）。推导与 \cref{eq:right-pert-so3} 同理：$\mathbf{T}\,\mathrm{Exp}(\delta\boldsymbol{\xi})\tilde{\mathbf{p}}\approx\mathbf{T}(\mathbf{I}+\delta\boldsymbol{\xi}^\wedge)\tilde{\mathbf{p}}=\mathbf{T}\tilde{\mathbf{p}}+\mathbf{T}\,\delta\boldsymbol{\xi}^\wedge\tilde{\mathbf{p}}$，故增量为 $\mathbf{T}\,\delta\boldsymbol{\xi}^\wedge\tilde{\mathbf{p}}=\mathbf{T}\,\tilde{\mathbf{p}}^{\,\odot}\delta\boldsymbol{\xi}$，对 $\delta\boldsymbol{\xi}$ 求导即得 \cref{eq:right-pert-se3}（其中 $\tilde{\mathbf{p}}^{\,\odot}$ 的右上块 $-\mathbf{p}^\wedge$ 来自 $\boldsymbol{\phi}^\wedge\mathbf{p}=-\mathbf{p}^\wedge\boldsymbol{\phi}$，左块 $\mathbf{I}$ 来自 $\boldsymbol{\rho}$）。

\paragraph{与左扰动的对照。}部分文献（如十四讲、Barfoot）惯用\emph{左}扰动；为便于对照，这里把 $\SE(3)$ 左扰动 $\partial(\mathbf{T}\tilde{\mathbf{p}})/\partial\delta\boldsymbol{\xi}$ 的分块推导补全：给 $\mathbf{T}$ \emph{左}乘 $\mathrm{Exp}(\delta\boldsymbol{\xi})\approx\mathbf{I}+\delta\boldsymbol{\xi}^\wedge$，增量为 $\delta\boldsymbol{\xi}^\wedge(\mathbf{T}\tilde{\mathbf{p}})$，按 $4\times4$ 分块展开（记 $\mathbf{T}\tilde{\mathbf{p}}=[\mathbf{R}\mathbf{p}+\mathbf{t};1]$）：
\begin{align}
  \delta\boldsymbol{\xi}^\wedge(\mathbf{T}\tilde{\mathbf{p}})
  &=\begin{bmatrix}\delta\boldsymbol{\phi}^\wedge&\delta\boldsymbol{\rho}\\\mathbf{0}^\top&0\end{bmatrix}
    \begin{bmatrix}\mathbf{R}\mathbf{p}+\mathbf{t}\\1\end{bmatrix}
   =\begin{bmatrix}\delta\boldsymbol{\phi}^\wedge(\mathbf{R}\mathbf{p}+\mathbf{t})+\delta\boldsymbol{\rho}\\0\end{bmatrix}\notag\\
  &=\begin{bmatrix}-(\mathbf{R}\mathbf{p}+\mathbf{t})^\wedge\delta\boldsymbol{\phi}+\delta\boldsymbol{\rho}\\0\end{bmatrix}
   =\begin{bmatrix}\mathbf{I}&-(\mathbf{R}\mathbf{p}+\mathbf{t})^\wedge\\\mathbf{0}^\top&\mathbf{0}^\top\end{bmatrix}
    \begin{bmatrix}\delta\boldsymbol{\rho}\\\delta\boldsymbol{\phi}\end{bmatrix},
  \label{eq:left-pert-se3}
\end{align}
故 $\partial(\mathbf{T}\tilde{\mathbf{p}})/\partial\delta\boldsymbol{\xi}=(\mathbf{T}\tilde{\mathbf{p}})^\odot=\big[\begin{smallmatrix}\mathbf{I}&-(\mathbf{R}\mathbf{p}+\mathbf{t})^\wedge\\\mathbf{0}^\top&\mathbf{0}^\top\end{smallmatrix}\big]$，旋转结果同理为 $\partial(\mathbf{R}\mathbf{p})/\partial\boldsymbol{\varphi}=-(\mathbf{R}\mathbf{p})^\wedge$。左、右结果相差一个 $\mathbf{R}$（伴随）的位置，可用 $(\mathbf{R}\mathbf{p})^\wedge=\mathbf{R}\mathbf{p}^\wedge\mathbf{R}^\top$ 互验：$-(\mathbf{R}\mathbf{p})^\wedge=-\mathbf{R}\mathbf{p}^\wedge\mathbf{R}^\top=(-\mathbf{R}\mathbf{p}^\wedge)\mathbf{R}^\top$。二者通过伴随 $\mathrm{Ad}(\mathbf{R})=\mathbf{R}$ 精确转换（\cref{sec:adjoint}）。

\begin{note}
\textbf{向量对向量求导的排布约定.}\;上面把分子 $[\mathbf{a};\mathbf{b}]$ 对分母 $[\mathbf{x};\mathbf{y}]$ 求导写成分块矩阵时，本书取分子布局 $\dfrac{\mathrm{d}[\mathbf{a};\mathbf{b}]}{\mathrm{d}[\mathbf{x};\mathbf{y}]}=\big[\begin{smallmatrix}\partial\mathbf{a}/\partial\mathbf{x}&\partial\mathbf{a}/\partial\mathbf{y}\\\partial\mathbf{b}/\partial\mathbf{x}&\partial\mathbf{b}/\partial\mathbf{y}\end{smallmatrix}\big]$（行随分子、列随分母，与十四讲\cite{gaoxiang2019slam14} 一致）。这正是为何 \cref{eq:left-pert-se3} 中平移块 $\delta\boldsymbol{\rho}$ 对应单位阵 $\mathbf{I}$（在 $\delta\boldsymbol{\rho}$ 列）、旋转块对应 $-(\cdot)^\wedge$（在 $\delta\boldsymbol{\phi}$ 列）。雅可比维度排布弄反是常见 bug，务必与本约定核对。
\end{note}

\begin{insight}[为什么右扰动 = 在“机体/局部系”下扰动]
Solà\cite{sola2018micro} 给了一个极清晰的图像：写成 $\mathbf{R}\,\mathrm{Exp}(\boldsymbol{\varphi})$（扰动在\emph{右}），增量 $\boldsymbol{\varphi}$ 表达在 $\mathbf{R}$ \emph{所在的局部坐标系}（机体系）；写成 $\mathrm{Exp}(\boldsymbol{\varphi})\,\mathbf{R}$（扰动在\emph{左}），$\boldsymbol{\varphi}$ 表达在\emph{全局/世界系}（单位元处）。SLAM 里 IMU 角速度、相机增量观测天然在机体系给出，故\emph{右扰动}与传感器/求解器（GTSAM、Ceres、Sophus 的右乘更新）天然对齐——这就是本书选右扰动为主的根本理由。
\end{insight}

\paragraph{优化中的位姿更新。}于是迭代里位姿的“无约束更新”写成右乘指数：
\begin{equation}
  \mathbf{R}\leftarrow\mathbf{R}\,\mathrm{Exp}(\delta\boldsymbol{\phi}^*),\qquad
  \mathbf{T}\leftarrow\mathbf{T}\,\mathrm{Exp}(\delta\boldsymbol{\xi}^*),
  \label{eq:right-update}
\end{equation}
其中 $\delta\boldsymbol{\phi}^*,\delta\boldsymbol{\xi}^*$ 是在切空间（李代数）解出的增量。这保证更新后仍严格在 $\SO(3)/\SE(3)$ 上——把带约束优化变成了切空间里的\emph{无约束}最小二乘。

\begin{pitfall}[编程陷阱：左右扰动混用导致雅可比与更新不一致]
最常见的实现 bug：雅可比按\emph{右}扰动推（\cref{eq:right-pert-so3}），更新却写成\emph{左}乘 $\mathrm{Exp}(\delta)\mathbf{R}$（或反之）。后果是收敛极慢甚至发散，且难定位。铁律：\textbf{雅可比的扰动侧、迭代更新的乘法侧、协方差定义的扰动侧，三者必须同侧}。本书统一右侧；引用左扰动公式时务必先用伴随转成右侧再用。
\end{pitfall}

\begin{practice}
\begin{enumerate}
  \item （推导，在草稿纸上完成）仿照 \cref{eq:right-pert-so3}，推导\emph{左}扰动下 $\partial(\mathbf{R}\mathbf{p})/\partial\boldsymbol{\varphi}=-(\mathbf{R}\mathbf{p})^\wedge$，并用 $(\mathbf{R}\mathbf{p})^\wedge=\mathbf{R}\mathbf{p}^\wedge\mathbf{R}^\top$ 与右扰动结果对照。
  \item 推导 \cref{eq:right-pert-se3}（提示：$\mathrm{Exp}(\delta\boldsymbol{\xi})\approx\mathbf{I}+\delta\boldsymbol{\xi}^\wedge$，按齐次坐标分块，用 $\boldsymbol{\xi}^\wedge\tilde{\mathbf{p}}=\tilde{\mathbf{p}}^{\,\odot}\boldsymbol{\xi}$）。
  \item （开放思考）若观测误差定义为 $\mathbf{e}=\mathrm{Log}(\mathbf{T}^{-1}\mathbf{Z})$（位姿测量 $\mathbf{Z}$），用右扰动 $\mathbf{T}\leftarrow\mathbf{T}\,\mathrm{Exp}(\delta\boldsymbol{\xi})$ 推 $\partial\mathbf{e}/\partial\delta\boldsymbol{\xi}$（结果含 $\boldsymbol{\mathcal{J}}_r^{-1}$）。
\end{enumerate}
\end{practice}

% ════════════════════════════════════════════════════════════════
\section{伴随：左右扰动的精确桥梁}
\label{sec:adjoint}

伴随回答一个反复出现的问题：把一个李代数元素“搬过”一个群元素，会变成什么？它是\emph{左右扰动转换、坐标系变换、协方差传播}的统一工具。

\paragraph{群伴随 $\mathrm{Ad}$。}定义 $\mathrm{Ad}(\mathbf{T})$ 使共轭 $\mathbf{T}\,\boldsymbol{\xi}^\wedge\mathbf{T}^{-1}=(\mathrm{Ad}(\mathbf{T})\boldsymbol{\xi})^\wedge$。其闭式为（Barfoot\cite{barfoot2024state}、Handbook\cite{carlone2026handbook}）
\begin{equation}
  \mathrm{Ad}(\mathbf{R})=\mathbf{R}\ \ (\SO(3)),\qquad
  \mathrm{Ad}(\mathbf{T})=\begin{bmatrix}\mathbf{R}&\mathbf{t}^\wedge\mathbf{R}\\\mathbf{0}&\mathbf{R}\end{bmatrix}\in\mathbb{R}^{6\times6}\ \ (\SE(3)).
  \label{eq:Ad}
\end{equation}
最有用的是\textbf{精确}（无近似）的“移动指数”恒等式：
\begin{equation}
  \mathbf{T}\,\mathrm{Exp}(\boldsymbol{\xi})=\mathrm{Exp}\big(\mathrm{Ad}(\mathbf{T})\boldsymbol{\xi}\big)\,\mathbf{T}.
  \label{eq:Ad-move}
\end{equation}
\cref{eq:Ad-move} 把\emph{右}扰动 $\boldsymbol{\xi}$ 精确换成\emph{左}扰动 $\mathrm{Ad}(\mathbf{T})\boldsymbol{\xi}$（反之亦然）。这正是“引用左扰动公式后转成右扰动”的依据，也是把扰动跨越已知变换搬移的标准手段。

\paragraph{小伴随 $\mathrm{ad}$（$\boldsymbol{\xi}^\curlywedge$）。}群伴随在单位元的导数即\emph{小伴随}，等价于李括号 $\mathrm{ad}(\boldsymbol{\xi}^\wedge)\mathbf{x}^\wedge=[\boldsymbol{\xi}^\wedge,\mathbf{x}^\wedge]$，记输出向量为 $\boldsymbol{\xi}^\curlywedge\mathbf{x}$：
\begin{equation}
  \boldsymbol{\phi}^\curlywedge=\boldsymbol{\phi}^\wedge\ \ (\so(3)),\qquad
  \boldsymbol{\xi}^\curlywedge=\begin{bmatrix}\boldsymbol{\rho}\\\boldsymbol{\phi}\end{bmatrix}^{\!\curlywedge}=\begin{bmatrix}\boldsymbol{\phi}^\wedge&\boldsymbol{\rho}^\wedge\\\mathbf{0}&\boldsymbol{\phi}^\wedge\end{bmatrix}\in\mathbb{R}^{6\times6}\ \ (\se(3)).
  \label{eq:ad}
\end{equation}
两者经指数相连：$\mathrm{Ad}(\mathrm{Exp}(\boldsymbol{\xi}))=\exp(\mathrm{ad}(\boldsymbol{\xi}))$，即 $\mathrm{Ad}(\mathbf{T})=\exp(\boldsymbol{\xi}^\curlywedge)$。这也解释了上节 $6\times6$ 雅可比 \cref{eq:se3-6x6} 里 $\boldsymbol{\xi}^\curlywedge$ 的来历。

\begin{pitfall}[概念误区：以为 $\mathrm{Ad}(\mathbf{T})$ 也是个“变换矩阵”]
$\mathrm{Ad}(\mathbf{T})$ 是 $6\times6$、作用在\emph{李代数向量 $\boldsymbol{\xi}\in\mathbb{R}^6$}上的线性映射，不是作用在点上的 $4\times4$ 位姿。它把“切空间里的扰动”从一个坐标系搬到另一个。混淆二者会在协方差传播（\cref{sec:uncertainty}）里把维度和对象用错。
\end{pitfall}

\begin{practice}
\begin{enumerate}
  \item 证明 $\SO(3)$ 伴随性质 $\mathbf{R}\,\mathrm{Exp}(\boldsymbol{\phi})\,\mathbf{R}^\top=\mathrm{Exp}(\mathbf{R}\boldsymbol{\phi})$，即 $\mathbf{R}\boldsymbol{\phi}^\wedge\mathbf{R}^\top=(\mathbf{R}\boldsymbol{\phi})^\wedge$。
  \item 由 \cref{eq:Ad-move} 推出：右扰动雅可比 $\mathbf{J}^{\text{右}}$ 与左扰动雅可比 $\mathbf{J}^{\text{左}}$ 满足 $\mathbf{J}^{\text{左}}=\mathbf{J}^{\text{右}}\,\mathrm{Ad}(\mathbf{T})^{-1}$（对 $\mathbf{T}\tilde{\mathbf{p}}$ 的两种结果做对照）。
\end{enumerate}
\end{practice}

% ════════════════════════════════════════════════════════════════
\section{流形视角：retraction 与 \texorpdfstring{$\boxplus\,\boxminus$}{boxplus/boxminus}（对接因子图）}
\label{sec:manifold}

前面用“群 + 李代数 + 指数”讲清了机理。现代 SLAM（因子图、GTSAM、Ceres）用一套更抽象、更可复用的语言把这一切统一起来——这套语言也解释了你在求解器里必须填的“流形局部参数化”到底是什么\cite{carlone2026handbook,hertzberg2013integrating}。

\paragraph{流形、切空间、局部坐标。}状态变量带约束（如 $\mathbf{R}\mathbf{R}^\top=\mathbf{I}$），约束把变量限制在一个\emph{流形} $\mathcal{M}$ 上。流形在每点 $\boldsymbol{\chi}$ 附近\emph{局部像欧氏空间}：其\emph{切空间} $T_{\boldsymbol{\chi}}\mathcal{M}$ 配一组基，给出该点的\emph{局部坐标}。优化时，我们不动 $\mathcal{M}$ 上的点，而是在切空间里解一个普通的欧氏增量 $\boldsymbol{\xi}$，再“贴回”流形。

\paragraph{retraction（收缩映射）。}把切空间增量映回流形的映射称为 retraction。指数映射 $\mathrm{Exp}$ 是\emph{首选}的 retraction，但\textbf{并非唯一}——任何满足一阶相切条件的映射都可用（Handbook、Boumal）。这把李群优化纳入更一般的\emph{黎曼优化}框架。

\paragraph{$\boxplus/\boxminus$ 算子（求解器接口）。}定义“李群元素 $\boxplus$ 切向量”与“两元素相 $\boxminus$ 得切向量”（本书取右乘约定，与 GTSAM 默认一致；Handbook 记作 $\oplus/\ominus$）：
\begin{equation}
  \mathbf{T}\boxplus\boldsymbol{\xi}:=\mathbf{T}\,\mathrm{Exp}(\boldsymbol{\xi}),\qquad
  \mathbf{T}_2\boxminus\mathbf{T}_1:=\mathrm{Log}\big(\mathbf{T}_1^{-1}\mathbf{T}_2\big),
  \label{eq:boxplus}
\end{equation}
满足互逆 $(\mathbf{T}\boxplus\boldsymbol{\xi})\boxminus\mathbf{T}=\boldsymbol{\xi}$。有了 $\boxplus/\boxminus$，\emph{流形优化在形式上与欧氏优化无异}：误差线性化只剩对切向量 $\boldsymbol{\xi}$ 的标准雅可比，解出 $\boldsymbol{\xi}^*$ 后用 $\boxplus$ 收缩回流形（\cref{eq:right-update}）。这正是 GTSAM 的 \texttt{retract}/\texttt{localCoordinates}、Ceres 的 \texttt{Manifold::Plus/Minus}、Sophus/manif 的右乘更新所封装的接口。

\begin{insight}
一句话主线：\textbf{流形} $\to$ \textbf{切空间局部坐标 $\boldsymbol{\xi}$} $\to$ \textbf{retraction $\mathrm{Exp}$（即 $\boxplus$）} $\to$ 误差线性化后只剩欧氏 $\boldsymbol{\xi}$ $\to$ 套标准 Gauss–Newton/LM $\to$ $\boxplus$ 收缩回流形。本书的\emph{右扰动}约定，就是为了让推导出的雅可比与这些求解器的 retract 约定\emph{逐位对齐}，使手推公式可以直接喂给 GTSAM/Ceres。
\end{insight}

\begin{pitfall}[编程陷阱：误把 $9$（或 $16$）维矩阵当优化维度]
直接把 $\mathbf{R}$ 的 $9$ 个元素交给优化器，会出现“过参数化 + 约束破坏”：维度多了、还得维持正交。$\boxplus$ 的意义就是把局部增量维度压到\emph{真实自由度}（$\SO(3)$ 为 $3$、$\SE(3)$ 为 $6$），从而\emph{无需}约束。在 Ceres/GTSAM 里这对应“给状态注册一个 Manifold/LocalParameterization”，漏掉它是新手最常见的收敛性 bug。
\end{pitfall}

% ════════════════════════════════════════════════════════════════
\section{（选读）李群上的插值}
\label{sec:lie-interp}

\paragraph{动机：两个关键帧之间的位姿怎么“插”出来？}
连续时间 SLAM、轨迹平滑、关键帧间补帧、IMU/相机的时间对齐，都要在两个已知位姿 $\mathbf{T}_1,\mathbf{T}_2$ 之间取一个“中间位姿”——给定插值参数 $\alpha\in[0,1]$，要一条从 $\mathbf{T}_1$（$\alpha{=}0$）平滑走到 $\mathbf{T}_2$（$\alpha{=}1$）的曲线。本节给出李群上做这件事的\emph{几何正确}方式，它正是连续时间轨迹表示（STEAM / GP 轨迹，\cref{sec:est-steam}）与 LiDAR 逐点去畸变的几何地基。

\paragraph{反面：逐元素线性插值不保群。}
向量空间里的线性插值 $\mathbf{x}=(1-\alpha)\mathbf{x}_1+\alpha\mathbf{x}_2$ 直接搬到旋转/位姿上立刻失效\cite{barfoot2024state}：
\begin{equation}
  (1-\alpha)\mathbf{R}_1+\alpha\mathbf{R}_2\notin\SO(3),\qquad
  (1-\alpha)\mathbf{T}_1+\alpha\mathbf{T}_2\notin\SE(3)
  \label{eq:interp-linear-fail}
\end{equation}
对一般 $\alpha\in(0,1)$ 成立——根子还是 \cref{eq:not-vs}：群对加法不封闭，线性组合落到流形\emph{外}。LiDAR 里常用的 LOAM“匀速线性插值”其实是这一错误的小角度一阶近似（$\boldsymbol{\phi}$ 很小时尚可，大转动则破坏正交、引入畸变）。要既插值又留在群上，必须换一种“乘法式”的定义。

\begin{theorem}[李群上的测地插值（保群）]\label{thm:lie-interp}
给定 $\mathbf{R}_1,\mathbf{R}_2\in\SO(3)$、$\mathbf{T}_1,\mathbf{T}_2\in\SE(3)$ 与 $\alpha\in[0,1]$，定义
\begin{align}
  \mathbf{R}(\alpha)&=\big(\mathbf{R}_2\mathbf{R}_1^\top\big)^{\alpha}\,\mathbf{R}_1,
  &\big(\mathbf{R}_2\mathbf{R}_1^\top\big)^{\alpha}&:=\mathrm{Exp}\big(\alpha\,\mathrm{Log}(\mathbf{R}_2\mathbf{R}_1^\top)\big),\label{eq:interp-so3}\\
  \mathbf{T}(\alpha)&=\big(\mathbf{T}_2\mathbf{T}_1^{-1}\big)^{\alpha}\,\mathbf{T}_1,
  &\big(\mathbf{T}_2\mathbf{T}_1^{-1}\big)^{\alpha}&:=\mathrm{Exp}\big(\alpha\,\mathrm{Log}(\mathbf{T}_2\mathbf{T}_1^{-1})\big).\label{eq:interp-se3}
\end{align}
则 $\mathbf{R}(\alpha)\in\SO(3)$、$\mathbf{T}(\alpha)\in\SE(3)$ 对一切 $\alpha\in[0,1]$ 成立，且 $\mathbf{R}(0)=\mathbf{R}_1$、$\mathbf{R}(1)=\mathbf{R}_2$（位姿同理）。其中“群元素的 $\alpha$ 次幂”由指数/对数定义：先 $\mathrm{Log}$ 到李代数、乘以 $\alpha$、再 $\mathrm{Exp}$ 回群。
\end{theorem}
\begin{proof}
记相对旋转 $\mathbf{R}_{21}=\mathbf{R}_2\mathbf{R}_1^\top=\mathrm{Exp}(\boldsymbol{\phi})\in\SO(3)$（由群封闭性，它确是旋转）。则 $\mathbf{R}_{21}^{\alpha}=\mathrm{Exp}(\alpha\boldsymbol{\phi})\in\SO(3)$（指数映射的像总在群内），再与 $\mathbf{R}_1$ 复合仍在 $\SO(3)$（封闭性）。端点：$\alpha{=}0$ 时 $\mathbf{R}_{21}^0=\mathbf{I}$，得 $\mathbf{R}_1$；$\alpha{=}1$ 时 $\mathbf{R}_{21}\mathbf{R}_1=\mathbf{R}_2$。$\SE(3)$ 把 $\mathbf{R}_1^\top$ 换成 $\mathbf{T}_1^{-1}$、$\boldsymbol{\phi}$ 换成 $\boldsymbol{\xi}$ 逐字平行。
\end{proof}

\paragraph{为何这叫“李代数里的线性插值”。}
把 \cref{eq:interp-so3} 取 $\mathrm{Log}$，在小角度（$\boldsymbol{\varphi}=\mathrm{Log}(\mathbf{R}(\alpha))$ 距 $\mathbf{R}_1$ 不远）下用 BCH（\cref{eq:bch-alg-to-group}）得
\begin{equation}
  \boldsymbol{\varphi}=\mathrm{Log}\big(\mathrm{Exp}(\alpha\boldsymbol{\phi})\,\mathrm{Exp}(\boldsymbol{\phi}_1)\big)
  \approx\alpha\,\mathbf{J}_l(\boldsymbol{\phi}_1)^{-1}\boldsymbol{\phi}+\boldsymbol{\phi}_1,
  \label{eq:interp-loglinear}
\end{equation}
形如向量插值 $\mathbf{x}=\alpha(\mathbf{x}_2-\mathbf{x}_1)+\mathbf{x}_1$——这正说明 \cref{eq:interp-so3} 是把插值搬进\emph{李代数}做线性，再映回群。一个无近似的特例：当 $\mathbf{R}_1=\mathbf{I}$ 时 $\mathbf{R}(\alpha)=\mathbf{R}_2^{\alpha}$、$\boldsymbol{\varphi}=\alpha\boldsymbol{\phi}_2$，干净的“按比例缩放转角”。

\begin{insight}[测地插值 = 匀速运动：与 Poisson 方程的精确联系]
\cref{eq:interp-so3} 不只是“一种可行方案”，它有强物理含义：它\emph{恰好}是匀角速度运动。把 $\alpha=(t-t_1)/(t_2-t_1)$、并令常值角速度 $\boldsymbol{\omega}=\boldsymbol{\phi}/(t_2-t_1)$（$\boldsymbol{\phi}=\mathrm{Log}(\mathbf{R}_2\mathbf{R}_1^\top)$），则
\[
  \mathbf{R}(t)=\mathrm{Exp}\big((t-t_1)\boldsymbol{\omega}^\wedge\big)\,\mathbf{R}(t_1)
\]
正是 Poisson 方程 $\dot{\mathbf{R}}=\boldsymbol{\omega}^\wedge\mathbf{R}$（\cref{eq:Rdot} 在常 $\boldsymbol{\omega}$ 下的解，见 \cref{ch:rigid_body} 运动学）在常角速度下的解。换言之，群上测地插值 = 两端点间“匀速转过去”，这也是连续时间轨迹去畸变比线性插值更准的根源：它尊重刚体的真实运动学，而非在矩阵元素上硬拉直线。$\SE(3)$ 版同理对应广义速度 $\boldsymbol{\varpi}$ 恒定的螺旋运动。
\end{insight}

\paragraph{扰动下的插值（供协方差/雅可比传播）。}
当端点带不确定度或要对端点求雅可比时，需知道端点的小扰动 $\delta\boldsymbol{\phi}_1,\delta\boldsymbol{\phi}_2$ 如何传到插值结果 $\delta\boldsymbol{\varphi}$。此处的一阶结果在\emph{左扰动}约定下形式最干净（即 $\mathrm{Exp}(\delta\boldsymbol{\phi}_i)\mathbf{R}_i$ 扰端点、$\mathrm{Exp}(\delta\boldsymbol{\varphi})\mathbf{R}(\alpha)$ 扰结果，$\boldsymbol{\phi}=\mathrm{Log}(\mathbf{R}_2\mathbf{R}_1^\top)$），仍像线性插值\cite{barfoot2024state}：
\begin{equation}
  \delta\boldsymbol{\varphi}=\big(\mathbf{I}-\mathbf{A}(\alpha,\boldsymbol{\phi})\big)\,\delta\boldsymbol{\phi}_1+\mathbf{A}(\alpha,\boldsymbol{\phi})\,\delta\boldsymbol{\phi}_2,
  \qquad
  \mathbf{A}(\alpha,\boldsymbol{\phi})=\alpha\,\mathbf{J}(\alpha\boldsymbol{\phi})\,\mathbf{J}(\boldsymbol{\phi})^{-1},
  \label{eq:interp-perturb}
\end{equation}
其中 $\mathbf{J}$ 为 $\SO(3)$ 左雅可比（\cref{eq:Jl}）。小转动时 $\mathbf{A}\to\alpha\mathbf{I}$，回到普通线性插值的权重 $(1-\alpha,\alpha)$。把 $\boldsymbol{\phi}\to\boldsymbol{\xi}$、$\mathbf{J}\to\boldsymbol{\mathcal{J}}_l$（\cref{eq:se3-6x6}）即得 $\SE(3)$ 版 $\boldsymbol{\mathcal{A}}(\alpha,\boldsymbol{\xi})=\alpha\boldsymbol{\mathcal{J}}_l(\alpha\boldsymbol{\xi})\boldsymbol{\mathcal{J}}_l(\boldsymbol{\xi})^{-1}$（$6\times6$）。\emph{约定提醒}：这干净的 $(\mathbf{I}-\mathbf{A},\mathbf{A})$ 分拆只在上述\emph{左}扰动下成立；若按本书主线的\emph{右}扰动 $\mathbf{R}_i\mathrm{Exp}(\delta\boldsymbol{\phi}_i)$、$\mathbf{R}(\alpha)\mathrm{Exp}(\delta\boldsymbol{\varphi})$ 重述，需用伴随把左、右扰动互换（$\delta\boldsymbol{\phi}^{\,L}=\mathrm{Ad}(\mathbf{R})\,\delta\boldsymbol{\phi}^{\,R}=\mathbf{R}\,\delta\boldsymbol{\phi}^{\,R}$），得 $\delta\boldsymbol{\varphi}=\mathbf{R}(\alpha)^\top\big[(\mathbf{I}-\mathbf{A})\mathbf{R}_1\,\delta\boldsymbol{\phi}_1+\mathbf{A}\mathbf{R}_2\,\delta\boldsymbol{\phi}_2\big]$——权重不再是单纯的 $(\mathbf{I}-\mathbf{A},\mathbf{A})$ 分拆，而夹了 $\mathbf{R}_1,\mathbf{R}_2,\mathbf{R}(\alpha)$ 的旋转因子。两式等价，按所需扰动侧择一即可。系数矩阵 $\mathbf{A}$ 还可展成 Faulhaber 级数 $\mathbf{A}(\alpha,\boldsymbol{\phi})=\sum_{n\ge0}\frac{F_n(\alpha)}{n!}(\boldsymbol{\phi}^\wedge)^n$，$F_n(\alpha)=\sum_{\beta=0}^{\alpha-1}\beta^n$ 为 Faulhaber 系数（$F_0=\alpha$、$F_1=\tfrac{\alpha(\alpha-1)}{2}$、$\dots$）；小转动只需头几项。

\begin{note}
\textbf{与本书其它插值的关系。}\;
(i) 球面线性插值（SLERP，见 \cref{ch:rigid_body} 四元数）是 \cref{eq:interp-so3} 在单位四元数上的等价写法——同一条测地线，换了参数化。(ii) 三次样条/B-spline 与连续时间 GP 轨迹（\cref{ch:slam_est} 的 STEAM）需要在多个控制位姿间做\emph{高阶}光滑插值，本节的两点测地插值是其最基本的构件（相邻控制点之间的“一段”）。(iii) LiDAR 章（\cref{ch:lidar_slam}）的运动去畸变可用本节插值替代匀速线性近似以提精度。
\end{note}

\begin{practice}
\begin{enumerate}
  \item 证明 $\mathbf{R}(\tfrac12)=(\mathbf{R}_2\mathbf{R}_1^\top)^{1/2}\mathbf{R}_1$ 是 $\mathbf{R}_1,\mathbf{R}_2$ 的“中点旋转”，并说明它一般\emph{不}等于 $\tfrac12(\mathbf{R}_1+\mathbf{R}_2)$ 投影回 $\SO(3)$（除非两旋转很近）。
  \item 用 \cref{eq:interp-perturb} 验证 $\alpha=0$ 时 $\delta\boldsymbol{\varphi}=\delta\boldsymbol{\phi}_1$、$\alpha=1$ 时 $\delta\boldsymbol{\varphi}=\delta\boldsymbol{\phi}_2$（提示：$\mathbf{A}(0,\cdot)=\mathbf{0}$、$\mathbf{A}(1,\cdot)=\mathbf{I}$）。
  \item （概念）解释为什么测地插值“匀速”这一性质，使它比逐点线性插值更适合 LiDAR 运动去畸变。
\end{enumerate}
\end{practice}

% ════════════════════════════════════════════════════════════════
\section{李群上的不确定度}
\label{sec:uncertainty}

状态估计要带\emph{协方差}。但给 $\mathbf{R}$ 直接加高斯噪声会破坏正交性。Barfoot\cite{barfoot2024state} 的标准做法是：\textbf{把高斯放在切空间（李代数），再用指数映射注入群}。要把这件事做\emph{严格}——说清“分布”到底是定义在哪、如何归一化、均值/协方差怎么定义——先得回答一个看似纯几何的问题：李群上怎么量距离、怎么算体积、怎么积分？这正是本节先补的地基。

\subsection{李群上的距离、体积元与积分}
\label{subsec:lie-measure}

\paragraph{动机：要在群上定义概率密度，先得有“体积”。}
概率密度 $p(\cdot)$ 的归一化 $\int p=1$ 里那个“$\int$”，在 $\mathbb{R}^n$ 上默认是 Lebesgue 测度 $d\mathbf{x}$。可 $\SO(3)/\SE(3)$ 是\emph{弯曲}的流形，没有现成的 $d\mathbf{x}$。要把高斯密度合法地放到群上、再谈“均值是不是众数”“归一化常数是多少”，就必须先在群上立一套\emph{距离}与\emph{体积元}。本小节给出三件事：群上两点的距离度量、无穷小体积元 $|\det\mathbf{J}|\,d\boldsymbol{\phi}$、以及把群上积分化为李代数上普通积分的换元公式。

\paragraph{$\so(3)$ 内积与两旋转的距离。}先在李代数上立内积。对 $\boldsymbol{\phi}_1^\wedge,\boldsymbol{\phi}_2^\wedge\in\so(3)$ 定义
\begin{equation}
  \langle\boldsymbol{\phi}_1^\wedge,\boldsymbol{\phi}_2^\wedge\rangle
  :=\tfrac12\operatorname{tr}\big(\boldsymbol{\phi}_1^\wedge\boldsymbol{\phi}_2^{\wedge\top}\big)=\boldsymbol{\phi}_1^\top\boldsymbol{\phi}_2,
  \label{eq:so3-inner}
\end{equation}
末等号可由 $\operatorname{tr}(\boldsymbol{\phi}_1^\wedge\boldsymbol{\phi}_2^{\wedge\top})=2\boldsymbol{\phi}_1^\top\boldsymbol{\phi}_2$ 直接验证（反对称阵的迹内积退化为向量点积）。两个旋转 $\mathbf{R}_1,\mathbf{R}_2$ 的\emph{距离}则取其相对旋转的对数之长：
\begin{equation}
  d(\mathbf{R}_1,\mathbf{R}_2)=\big\|\mathrm{Log}(\mathbf{R}_1^\top\mathbf{R}_2)\big\|=\big|\boldsymbol{\phi}_{12}\big|,\qquad
  \boldsymbol{\phi}_{12}=\mathrm{Log}(\mathbf{R}_1^\top\mathbf{R}_2).
  \label{eq:so3-distance}
\end{equation}
这正是 \cref{sec:traj-error} 里 ATE/RPE 反复用的“相对位姿 $\to\mathrm{Log}\to$ 取范数”，现在有了内积\cref{eq:so3-inner} 它是一个名正言顺的\emph{黎曼距离}（右不变度量下的测地距离）。$\SE(3)$ 同理，用 $\boldsymbol{\xi}_{12}=\mathrm{Log}(\mathbf{T}_1^{-1}\mathbf{T}_2)$、$d(\mathbf{T}_1,\mathbf{T}_2)=|\boldsymbol{\xi}_{12}|$；其内积须对平移/旋转块加权（否则量纲不齐），取带权迹形式 $\langle\boldsymbol{\xi}_1^\wedge,\boldsymbol{\xi}_2^\wedge\rangle=\operatorname{tr}\big(\boldsymbol{\xi}_1^\wedge\operatorname{diag}(\tfrac12\mathbf{I},1)\,\boldsymbol{\xi}_2^{\wedge\top}\big)=\boldsymbol{\xi}_1^\top\boldsymbol{\xi}_2$，即恰好诱导出欧氏点积\cite{barfoot2024state}。

\paragraph{无穷小体积元：$|\det\mathbf{J}|$ 因子。}关键来了。把 $\mathbf{R}=\mathrm{Exp}(\boldsymbol{\phi})$ 用李代数坐标 $\boldsymbol{\phi}$ 参数化，给 $\boldsymbol{\phi}$ 一个微小变化 $\boldsymbol{\phi}\to\boldsymbol{\phi}+\delta\boldsymbol{\phi}$，对应的群元素变化用右扰动度量为 $\mathrm{Log}(\mathbf{R}^\top\mathbf{R}')\approx\mathbf{J}_r\delta\boldsymbol{\phi}$（左扰动则 $\mathbf{J}_l\delta\boldsymbol{\phi}$，由 \cref{eq:bch-alg-to-group} 立得）。于是群上的“无穷小体积”与李代数上的 $d\boldsymbol{\phi}$ 相差一个雅可比行列式：
\begin{equation}
  d\mathbf{R}=|\det\mathbf{J}|\,d\boldsymbol{\phi},\qquad
  d\mathbf{T}=|\det\boldsymbol{\mathcal{J}}|\,d\boldsymbol{\xi}.
  \label{eq:volume-element}
\end{equation}
这里出现了一个\emph{漂亮且重要}的事实：用左还是右扰动，体积元\emph{一样}。因为
\begin{equation}
  \det\mathbf{J}_l=\det(\mathbf{R}\,\mathbf{J}_r)=\underbrace{\det\mathbf{R}}_{=1}\det\mathbf{J}_r=\det\mathbf{J}_r,
  \label{eq:unimodular}
\end{equation}
（用了 $\mathbf{J}_l=\mathbf{R}\,\mathbf{J}_r$，见 \cref{thm:so3-kinematics} 证明末）——这称为 $\SO(3)$ 的\textbf{幺模性}（unimodular，左右不变测度一致）。$\SE(3)$ 同理 $\det\boldsymbol{\mathcal{J}}_l=\det(\boldsymbol{\mathcal{T}}\boldsymbol{\mathcal{J}}_r)=\det\boldsymbol{\mathcal{J}}_r$（$\det\boldsymbol{\mathcal{T}}=(\det\mathbf{R})^2=1$），故 \cref{eq:volume-element} 不必区分左右，这个统一的测度即李群的\emph{Haar 测度}。

\begin{theorem}[$\SO(3)$ 体积元的闭式]\label{thm:detJ}
左/右雅可比的行列式只依赖转角 $\phi=\|\boldsymbol{\phi}\|$，且
\begin{equation}
  |\det\mathbf{J}|=2\,\frac{1-\cos\phi}{\phi^2}
  =1-\frac{1}{12}\phi^2+\frac{1}{360}\phi^4-\cdots\ \xrightarrow{\phi\to0}1.
  \label{eq:detJ}
\end{equation}
$\SE(3)$ 则 $|\det\boldsymbol{\mathcal{J}}|=|\det\mathbf{J}|^2=\big(2\tfrac{1-\cos\phi}{\phi^2}\big)^2$。
\end{theorem}
\begin{proof}
$\mathbf{J}$ 的三个特征值由 \cref{eq:Jl} 读出：沿轴方向 $\mathbf{a}$ 是特征向量，特征值 $1$（前已述 $\mathbf{J}\mathbf{a}=\mathbf{a}$）；垂直于轴的平面内，$\mathbf{J}$ 在基 $\{\mathbf{I},\mathbf{a}^\wedge\}$ 上作用为复数 $\tfrac{\sin\phi}{\phi}+i\tfrac{1-\cos\phi}{\phi}$ 的乘法（把 $\mathbf{a}^\wedge$ 当虚单位 $i$，因 $(\mathbf{a}^\wedge)^2=\mathbf{a}\mathbf{a}^\top-\mathbf{I}$ 在垂直子空间上 $=-\mathbf{I}$），故另两个特征值是这对共轭复数。三者之积
\[
  \det\mathbf{J}=1\cdot\Big(\tfrac{\sin\phi}{\phi}+i\tfrac{1-\cos\phi}{\phi}\Big)\Big(\tfrac{\sin\phi}{\phi}-i\tfrac{1-\cos\phi}{\phi}\Big)
  =\frac{\sin^2\phi+(1-\cos\phi)^2}{\phi^2}=\frac{2(1-\cos\phi)}{\phi^2},
\]
末步用 $\sin^2\phi+(1-\cos\phi)^2=2-2\cos\phi$。$\SE(3)$ 的 $\boldsymbol{\mathcal{J}}=\big[\begin{smallmatrix}\mathbf{J}&\mathbf{Q}\\\mathbf{0}&\mathbf{J}\end{smallmatrix}\big]$（\cref{eq:se3-6x6}）是分块上三角，故 $\det\boldsymbol{\mathcal{J}}=(\det\mathbf{J})^2$。
\end{proof}

\paragraph{李群上的积分。}有了体积元，群上积分即化为李代数上的普通 Lebesgue 积分（对 $\SO(3)$ 因指数映射满射非单射，须限 $|\boldsymbol{\phi}|<\pi$ 以恰好扫过一次）：
\begin{equation}
  \int_{\SO(3)}f(\mathbf{R})\,d\mathbf{R}=\int_{|\boldsymbol{\phi}|<\pi}f\big(\mathrm{Exp}(\boldsymbol{\phi})\big)\,|\det\mathbf{J}|\,d\boldsymbol{\phi},\qquad
  \int_{\SE(3)}f\,d\mathbf{T}=\int_{\boldsymbol{\rho}\in\mathbb{R}^3,\,|\boldsymbol{\phi}|<\pi}f\,|\det\boldsymbol{\mathcal{J}}|\,d\boldsymbol{\xi}.
  \label{eq:lie-integral}
\end{equation}
（$\SE(3)$ 的平移分量 $\boldsymbol{\rho}\in\mathbb{R}^3$ 无界，只有旋转受 $|\boldsymbol{\phi}|<\pi$ 约束。）

\begin{insight}[为什么这个 $|\det\mathbf{J}|$ 因子值得专门立一节]
\cref{eq:lie-integral} 的 $|\det\mathbf{J}|$ 看似只是换元的雅可比，却是把“切空间高斯”升级为“群上严格概率密度”的\emph{唯一}缺失环节（下面马上用到）。它还解释了一个反直觉现象：当不确定度大时，群上分布的\emph{众数}（峰值）与\emph{均值}（一阶矩）不再重合——正是因为 $|\det\mathbf{J}|<1$（$\phi\neq0$ 时）这个随 $\phi$ 衰减的因子把密度往大角度方向“压”，使峰值偏离切空间高斯的中心。这与后面的“香蕉形”分布同源。
\end{insight}

\subsection{切空间高斯与诱导密度}
\label{subsec:tangent-gaussian}

\paragraph{切空间高斯。}以右扰动定义 $\SE(3)$ 上的高斯：
\begin{equation}
  \mathbf{T}=\bar{\mathbf{T}}\,\mathrm{Exp}(\boldsymbol{\epsilon}),\qquad \boldsymbol{\epsilon}\sim\mathcal{N}(\mathbf{0},\boldsymbol{\Sigma}),\ \boldsymbol{\epsilon}\in\mathbb{R}^6,
  \label{eq:lie-gaussian}
\end{equation}
$\bar{\mathbf{T}}$ 为均值（一个确定位姿），$\boldsymbol{\Sigma}$ 是\emph{切空间里}的协方差。如此 $\mathbf{T}$ 恒在 $\SE(3)$ 上，而“分布”由 $\mathbb{R}^6$ 上的标准高斯诱导。$\SO(3)$ 同理 $\mathbf{R}=\bar{\mathbf{R}}\,\mathrm{Exp}(\boldsymbol{\epsilon})$，$\boldsymbol{\epsilon}\in\mathbb{R}^3$。

\paragraph{诱导到群上的密度（用上 $|\det\mathbf{J}|$）。}\cref{eq:lie-gaussian} 把高斯放在切向量 $\boldsymbol{\epsilon}$ 上：$p(\boldsymbol{\epsilon})=\eta\exp(-\tfrac12\boldsymbol{\epsilon}^\top\boldsymbol{\Sigma}^{-1}\boldsymbol{\epsilon})$。但若问“位姿 $\mathbf{T}$ 本身的密度 $p(\mathbf{T})$ 是多少”，就要用 \cref{subsec:lie-measure} 的体积元做变量替换 $d\mathbf{T}=|\det\boldsymbol{\mathcal{J}}(\boldsymbol{\epsilon})|\,d\boldsymbol{\epsilon}$：
\begin{equation}
  p(\mathbf{T})=\frac{p(\boldsymbol{\epsilon})}{|\det\boldsymbol{\mathcal{J}}(\boldsymbol{\epsilon})|}
  =\frac{\eta}{|\det\boldsymbol{\mathcal{J}}(\boldsymbol{\epsilon})|}\exp\!\Big(-\tfrac12\boldsymbol{\epsilon}^\top\boldsymbol{\Sigma}^{-1}\boldsymbol{\epsilon}\Big),\qquad
  \boldsymbol{\epsilon}=\mathrm{Log}(\bar{\mathbf{T}}^{-1}\mathbf{T}).
  \label{eq:induced-density}
\end{equation}
当协方差\emph{小}（$\boldsymbol{\epsilon}$ 集中在 $\mathbf{0}$ 附近、$|\det\boldsymbol{\mathcal{J}}|\to1$）时，$p(\mathbf{T})$ 与切空间高斯几乎无异——这正是绝大多数 SLAM 把协方差直接放切空间、忽略 $|\det\boldsymbol{\mathcal{J}}|$ 的合法性依据。\textbf{均值的隐式定义}也由此而来：把 $\bar{\mathbf{T}}$ 定义为使 $E[\boldsymbol{\epsilon}]=E[\mathrm{Log}(\bar{\mathbf{T}}^{-1}\mathbf{T})]=\mathbf{0}$ 成立的那个位姿（即“切空间一阶矩为零点”），协方差则 $\boldsymbol{\Sigma}=E[\boldsymbol{\epsilon}\boldsymbol{\epsilon}^\top]$。这是一个\emph{隐式}方程（$\bar{\mathbf{T}}$ 出现在 $\boldsymbol{\epsilon}$ 的定义里），实际由迭代求解，与下文的位姿融合（\cref{subsec:pose-fusion}）是同一套 Gauss--Newton。

\paragraph{左/中/右三种扰动的取法。}\cref{eq:lie-gaussian} 把扰动放在\emph{右}（$\bar{\mathbf{T}}\,\mathrm{Exp}(\boldsymbol{\epsilon})$）。也可放\emph{左}（$\mathrm{Exp}(\boldsymbol{\epsilon})\,\bar{\mathbf{T}}$，Barfoot 主线）或\emph{中}（对称地分摊到 $\bar{\mathbf{T}}$ 两侧）。三者的协方差经伴随精确互换（右 $\to$ 左：$\boldsymbol{\Sigma}_l=\mathrm{Ad}(\bar{\mathbf{T}})\boldsymbol{\Sigma}_r\mathrm{Ad}(\bar{\mathbf{T}})^\top$）。本书统一取\emph{右}（机体系扰动，与 \cref{sec:perturb} 的右扰动雅可比、求解器的右乘更新一致）；Barfoot 取左，故其协方差公式里的伴随位置与本书镜像，引用时按此换算即可。

\paragraph{确定变换下分布的精确传播。}一个常被忽略却\emph{无近似}的好性质：若对随机位姿左乘一个\emph{确定}的变换 $\mathbf{T}'=\mathbf{T}_0\mathbf{T}$（$\mathbf{T}_0$ 无噪），则 $\mathbf{T}'=\mathbf{T}_0\bar{\mathbf{T}}\,\mathrm{Exp}(\boldsymbol{\epsilon})$ 仍是同一形式，均值变 $\bar{\mathbf{T}}'=\mathbf{T}_0\bar{\mathbf{T}}$ 而\emph{协方差不变}（右扰动下 $\boldsymbol{\Sigma}'=\boldsymbol{\Sigma}$）；对 $\SO(3)$ 旋转点 $\mathbf{R}'=\mathbf{R}_0\bar{\mathbf{R}}\mathrm{Exp}(\boldsymbol{\epsilon})$ 亦然。这说明“左乘确定旋转”只搬动均值、不改变切空间协方差形状——是相机外参、坐标系变换下协方差传播无损的根据。

\paragraph{位姿复合的协方差传播。}两个相互独立的含噪位姿复合 $\mathbf{T}=\mathbf{T}_1\mathbf{T}_2$，均值为 $\bar{\mathbf{T}}=\bar{\mathbf{T}}_1\bar{\mathbf{T}}_2$，协方差到\emph{一阶}（二阶矩）为
\begin{equation}
  \boldsymbol{\Sigma}\approx\boldsymbol{\Sigma}_1+\mathrm{Ad}(\bar{\mathbf{T}}_1)\,\boldsymbol{\Sigma}_2\,\mathrm{Ad}(\bar{\mathbf{T}}_1)^\top.
  \label{eq:compound-cov}
\end{equation}
这里伴随 $\mathrm{Ad}(\bar{\mathbf{T}}_1)$ 把第二个位姿的切空间协方差“搬运”到复合后的切空间——\textbf{协方差传播的本质，就是用伴随在不同均值点之间搬运切空间协方差}。\cref{eq:compound-cov} 是 VIO 预积分、位姿图边协方差的基石。

\paragraph{（选读）四阶修正与“香蕉形”分布。}当不确定度较大、或多段位姿链式复合时，一阶传播低估了真实弥散，分布呈现弯曲的“香蕉形”。Barfoot 给出在扰动上展开到\emph{四阶}的修正（见 \cref{sec:lie-appendix} 的指引）：
\begin{equation}
  \boldsymbol{\Sigma}_{4\text{th}}\approx\boldsymbol{\Sigma}_1+\boldsymbol{\Sigma}_2'+\tfrac14\boldsymbol{\mathcal{B}}+\tfrac{1}{12}\big(\boldsymbol{\mathcal{A}}_1\boldsymbol{\Sigma}_2'+\boldsymbol{\Sigma}_2'\boldsymbol{\mathcal{A}}_1^\top+\boldsymbol{\mathcal{A}}_2'\boldsymbol{\Sigma}_1+\boldsymbol{\Sigma}_1\boldsymbol{\mathcal{A}}_2'^\top\big),
  \label{eq:fourth-order}
\end{equation}
其中 $\boldsymbol{\Sigma}_2'=\mathrm{Ad}(\bar{\mathbf{T}}_1)\boldsymbol{\Sigma}_2\mathrm{Ad}(\bar{\mathbf{T}}_1)^\top$，算子 $\boldsymbol{\mathcal{A}},\boldsymbol{\mathcal{B}}$ 的定义见附录指引。本书在状态估计部分（\cref{ch:slam_est}）会再用到；初读可只记一阶式 \cref{eq:compound-cov}。

\subsection{位姿求逆与相关位姿的复合 / 差分}
\label{subsec:correlated-pose}

\paragraph{位姿求逆：协方差获精确闭式。}\cref{eq:compound-cov} 处理的是位姿\emph{复合}。同样常用的是位姿\emph{求逆}的协方差——例如把 $\mathbf{T}_{wb}$（机体在世界系）的不确定度换算到 $\mathbf{T}_{bw}=\mathbf{T}_{wb}^{-1}$。设 $\mathbf{T}=\bar{\mathbf{T}}\,\mathrm{Exp}(\boldsymbol{\epsilon})$、$\boldsymbol{\epsilon}\sim\mathcal{N}(\mathbf{0},\boldsymbol{\Sigma})$（右扰动），则
\begin{equation}
  \mathbf{T}^{-1}=\mathrm{Exp}(\boldsymbol{\epsilon})^{-1}\bar{\mathbf{T}}^{-1}=\mathrm{Exp}(-\boldsymbol{\epsilon})\,\bar{\mathbf{T}}^{-1}
  =\bar{\mathbf{T}}^{-1}\,\mathrm{Exp}\big(\underbrace{-\mathrm{Ad}(\bar{\mathbf{T}})\,\boldsymbol{\epsilon}}_{=:\,\boldsymbol{\epsilon}'}\big),
  \label{eq:pose-inv-derivation}
\end{equation}
末步把左扰动 $-\boldsymbol{\epsilon}$ 经移指数恒等式 \cref{eq:Ad-move} 搬到右侧（$\mathrm{Exp}(-\boldsymbol{\epsilon})\bar{\mathbf{T}}^{-1}=\bar{\mathbf{T}}^{-1}\mathrm{Exp}(-\mathrm{Ad}(\bar{\mathbf{T}})\boldsymbol{\epsilon})$）。故 $\mathbf{T}^{-1}$ 仍是以 $\bar{\mathbf{T}}^{-1}$ 为均值的右扰动高斯，协方差\textbf{无近似}地为
\begin{equation}
  \boxed{\;\boldsymbol{\Sigma}_{\text{inv}}=\mathrm{Ad}(\bar{\mathbf{T}})\,\boldsymbol{\Sigma}\,\mathrm{Ad}(\bar{\mathbf{T}})^\top\;}
  \label{eq:pose-inverse-cov}
\end{equation}
（因 $\boldsymbol{\epsilon}'=-\mathrm{Ad}(\bar{\mathbf{T}})\boldsymbol{\epsilon}$ 是 $\boldsymbol{\epsilon}$ 的\emph{线性}变换，协方差传播精确）。这是 \cref{eq:compound-cov} 之外伴随作为“协方差搬运器”的又一干净例证：求逆即用 $\mathrm{Ad}(\bar{\mathbf{T}})$ 把切空间协方差搬到逆位姿的切空间。

\paragraph{相关位姿的复合与差分。}\cref{eq:compound-cov} 假设两位姿\emph{独立}。但位姿图里相邻节点、同一里程计链上的位姿往往\emph{相关}（共享了同一段噪声测量），其联合分布带交叉协方差 $\boldsymbol{\Sigma}_{12}=E[\boldsymbol{\epsilon}_1\boldsymbol{\epsilon}_2^\top]\neq\mathbf{0}$：
\begin{equation}
  \begin{bmatrix}\boldsymbol{\epsilon}_1\\\boldsymbol{\epsilon}_2\end{bmatrix}\sim\mathcal{N}\!\left(\mathbf{0},\,\begin{bmatrix}\boldsymbol{\Sigma}_1&\boldsymbol{\Sigma}_{12}\\\boldsymbol{\Sigma}_{12}^\top&\boldsymbol{\Sigma}_2\end{bmatrix}\right).
  \label{eq:joint-corr}
\end{equation}
此时\emph{复合} $\mathbf{T}=\mathbf{T}_1\mathbf{T}_2$ 的协方差比 \cref{eq:compound-cov} 多两项交叉贡献（到二阶）：
\begin{equation}
  \boldsymbol{\Sigma}\approx\boldsymbol{\Sigma}_1+\bar{\boldsymbol{\mathcal{T}}}_1\boldsymbol{\Sigma}_2\bar{\boldsymbol{\mathcal{T}}}_1^\top
  +\boldsymbol{\Sigma}_{12}\bar{\boldsymbol{\mathcal{T}}}_1^\top+\bar{\boldsymbol{\mathcal{T}}}_1\boldsymbol{\Sigma}_{12}^\top,
  \qquad \bar{\boldsymbol{\mathcal{T}}}_1=\mathrm{Ad}(\bar{\mathbf{T}}_1).
  \label{eq:corr-compound}
\end{equation}
而\emph{差分}（相对位姿）$\mathbf{T}=\mathbf{T}_1\mathbf{T}_2^{-1}$ 的协方差则\emph{减去}交叉项（注意 $\bar{\boldsymbol{\mathcal{T}}}=\mathrm{Ad}(\bar{\mathbf{T}}_1\bar{\mathbf{T}}_2^{-1})$）：
\begin{equation}
  \boldsymbol{\Sigma}\approx\boldsymbol{\Sigma}_1+\bar{\boldsymbol{\mathcal{T}}}\boldsymbol{\Sigma}_2\bar{\boldsymbol{\mathcal{T}}}^\top
  -\boldsymbol{\Sigma}_{12}\bar{\boldsymbol{\mathcal{T}}}^\top-\bar{\boldsymbol{\mathcal{T}}}\boldsymbol{\Sigma}_{12}^\top.
  \label{eq:corr-diff}
\end{equation}

\begin{pitfall}[陷阱：把相关位姿当独立，回环/位姿图协方差系统性失真]
位姿图边、回环约束常由“两个绝对位姿估计相减”得到，而这两个估计若来自同一段里程计就\emph{高度相关}。此时若漏掉 \cref{eq:corr-diff} 的交叉项 $-\boldsymbol{\Sigma}_{12}\bar{\boldsymbol{\mathcal{T}}}^\top-\bar{\boldsymbol{\mathcal{T}}}\boldsymbol{\Sigma}_{12}^\top$，正相关时会\emph{高估}相对位姿不确定度（边权过小、优化对该约束信任不足），反之亦可能\emph{低估}导致过自信。差分减交叉、复合加交叉——符号正好相反，记混会让协方差失去正定性。本书后续位姿图章节（\cref{ch:slam_est}）默认独立处理是\emph{简化}；严格相关情形须用 \cref{eq:corr-compound,eq:corr-diff}。
\end{pitfall}

\subsection{位姿融合：李群上的 Gauss--Newton}
\label{subsec:pose-fusion}

\paragraph{动机：把 $K$ 个带噪位姿估计合成一个。}多传感器（视觉 + 激光 + 轮速）或多次观测常给出对\emph{同一}位姿的若干带协方差估计 $\{\bar{\mathbf{T}}_k,\boldsymbol{\Sigma}_k\}_{k=1}^K$。如何把它们最优融合成单一 $\{\bar{\mathbf{T}},\boldsymbol{\Sigma}\}$？在 $\mathbb{R}^n$ 上这是经典的\emph{信息形式加权平均} $\bar{\mathbf{x}}=(\sum_k\boldsymbol{\Sigma}_k^{-1})^{-1}\sum_k\boldsymbol{\Sigma}_k^{-1}\mathbf{x}_k$。但位姿不能直接加权平均（$\SE(3)$ 无加法），须把它升级为\emph{李群上的 Gauss--Newton}——这恰好把本章 \cref{sec:perturb}（右扰动雅可比）、\cref{subsec:tangent-gaussian}（切空间协方差）、\cref{subsec:correlated-pose}（伴随搬运）三者合流。

\paragraph{误差、雅可比与正规方程。}以待求融合位姿 $\mathbf{T}$ 为优化变量，定义第 $k$ 个估计的残差为相对位姿的对数（取 $\mathbf{T}^{-1}\bar{\mathbf{T}}_k$ 的次序使扰动落在末端、最契合右扰动）
\begin{equation}
  \mathbf{e}_k(\mathbf{T})=\mathrm{Log}\big(\mathbf{T}^{-1}\,\bar{\mathbf{T}}_k\big)\in\mathbb{R}^6,
  \label{eq:fusion-error}
\end{equation}
完全一致时 $\mathbf{e}_k=\mathbf{0}$。代价为各残差按各自信息阵 $\boldsymbol{\Sigma}_k^{-1}$ 加权的马氏距离之和 $J(\mathbf{T})=\tfrac12\sum_k\mathbf{e}_k^\top\boldsymbol{\Sigma}_k^{-1}\mathbf{e}_k$。用\emph{右}扰动 $\mathbf{T}\leftarrow\mathbf{T}\,\mathrm{Exp}(\boldsymbol{\epsilon})$ 在当前估计处线性化 $\mathbf{e}_k$。把扰动搬入对数的自变量：
\[
  \big(\mathbf{T}\,\mathrm{Exp}(\boldsymbol{\epsilon})\big)^{-1}\bar{\mathbf{T}}_k
  =\mathrm{Exp}(-\boldsymbol{\epsilon})\,\mathbf{T}^{-1}\bar{\mathbf{T}}_k
  =\mathrm{Exp}(-\boldsymbol{\epsilon})\,\underbrace{\mathbf{T}^{-1}\bar{\mathbf{T}}_k}_{=\,\mathrm{Exp}(\mathbf{e}_k)},
\]
取对数并用 BCH 一阶式的\emph{左乘}形式（\cref{eq:bch-se3-left}，$\mathrm{Log}(\mathrm{Exp}(\boldsymbol{\delta})\mathrm{Exp}(\mathbf{e}_k))\approx\mathbf{e}_k+\boldsymbol{\mathcal{J}}_l(\mathbf{e}_k)^{-1}\boldsymbol{\delta}$，此处 $\boldsymbol{\delta}=-\boldsymbol{\epsilon}$），得
\begin{equation}
  \mathbf{e}_k(\mathbf{T}\,\mathrm{Exp}(\boldsymbol{\epsilon}))\approx\mathbf{e}_k(\mathbf{T})+\mathbf{G}_k\,\boldsymbol{\epsilon},
  \qquad
  \mathbf{G}_k=-\boldsymbol{\mathcal{J}}_l\!\big(\mathbf{e}_k(\mathbf{T})\big)^{-1},
  \label{eq:fusion-jacobian}
\end{equation}
其中 $\boldsymbol{\mathcal{J}}_l$ 为 $\SE(3)$ 的 $6\times6$ 左雅可比（\cref{eq:se3-6x6}），负号来自扰动 $\mathrm{Exp}(-\boldsymbol{\epsilon})$ 落在对数自变量\emph{左}侧。令 $\partial J/\partial\boldsymbol{\epsilon}=\mathbf{0}$ 得\textbf{正规方程}与最优增量
\begin{equation}
  \Big(\sum_{k=1}^{K}\mathbf{G}_k^\top\boldsymbol{\Sigma}_k^{-1}\mathbf{G}_k\Big)\boldsymbol{\epsilon}^\star
  =-\sum_{k=1}^{K}\mathbf{G}_k^\top\boldsymbol{\Sigma}_k^{-1}\mathbf{e}_k(\mathbf{T}),
  \label{eq:fusion-normal}
\end{equation}
随后右乘更新 $\mathbf{T}\leftarrow\mathbf{T}\,\mathrm{Exp}(\boldsymbol{\epsilon}^\star)$ 并\emph{迭代}至收敛。收敛后取 $\bar{\mathbf{T}}=\mathbf{T}$，融合协方差即正规方程左端（Hessian）之逆：
\begin{equation}
  \boxed{\;\boldsymbol{\Sigma}=\Big(\sum_{k=1}^{K}\mathbf{G}_k^\top\boldsymbol{\Sigma}_k^{-1}\mathbf{G}_k\Big)^{-1}\;}
  \label{eq:fusion-cov}
\end{equation}

\begin{insight}[融合 = 李群上的“信息相加”]
对照 $\mathbb{R}^n$ 的信息形式：\cref{eq:fusion-cov} 的 $\sum_k\mathbf{G}_k^\top\boldsymbol{\Sigma}_k^{-1}\mathbf{G}_k$ 就是把各估计的信息阵 $\boldsymbol{\Sigma}_k^{-1}$ 经雅可比 $\mathbf{G}_k$ 搬到\emph{公共切空间}再相加——这正是欧氏“信息相加”在弯曲流形上的推广。当所有 $\bar{\mathbf{T}}_k$ 接近、$\mathbf{e}_k$ 小（$\boldsymbol{\mathcal{J}}_l\to\mathbf{I}$、$\mathbf{G}_k\to-\mathbf{I}$）时，$\mathbf{G}_k^\top\boldsymbol{\Sigma}_k^{-1}\mathbf{G}_k\to\boldsymbol{\Sigma}_k^{-1}$，\cref{eq:fusion-cov} 退化为 $\boldsymbol{\Sigma}=(\sum_k\boldsymbol{\Sigma}_k^{-1})^{-1}$、增量退化为 $\boldsymbol{\epsilon}^\star=\boldsymbol{\Sigma}\sum_k\boldsymbol{\Sigma}_k^{-1}\mathbf{e}_k$，与欧氏加权平均\emph{逐字一致}。这就是“位姿融合 = 协方差加权平均”在小不确定度下的严格依据，也说明为何 $\mathbf{G}_k$（即 $-\boldsymbol{\mathcal{J}}_l^{-1}$）是把残差从“相对位姿切空间”搬到“融合位姿切空间”的桥——又一次雅可比作为协方差搬运器。\textbf{注意}：Barfoot 以\emph{左}扰动推导本算法（误差 $\mathbf{e}_k=\ln(\bar{\mathbf{T}}_k\mathbf{T}^{-1})^\vee$、更新 $\mathbf{T}\leftarrow\mathrm{Exp}(\boldsymbol{\epsilon}^\star)\mathbf{T}$、$\mathbf{G}_k=\boldsymbol{\mathcal{J}}_l(-\mathbf{e}_k)^{-1}$）\cite{barfoot2024state}；本书改写为右扰动版（误差取 $\mathbf{T}^{-1}\bar{\mathbf{T}}_k$、更新右乘、$\mathbf{G}_k=-\boldsymbol{\mathcal{J}}_l(\mathbf{e}_k)^{-1}$），数学内容等价，仅扰动侧与本书主线对齐。
\end{insight}

\begin{insight}
本节是把“位姿带不确定度”讲清的\emph{关键一步}：核心思想是\textbf{位姿不确定度 = 切空间里的高斯 + 伴随搬运}（\cref{eq:lie-gaussian,eq:compound-cov}），由此协方差传播、位姿求逆（\cref{eq:pose-inverse-cov}）、相关位姿差分（\cref{eq:corr-diff}）、多估计融合（\cref{eq:fusion-cov}）都获得干净闭式\cite{barfoot2024state}。它把 \cref{sec:adjoint} 的伴随从“代数技巧”升格为“概率工具”，是后续滤波/平滑章节处理位姿协方差的前置。
\end{insight}

% ════════════════════════════════════════════════════════════════
\section{（选读）相似变换群 Sim(3)}
\label{sec:sim3}

单目 SLAM 存在\emph{尺度不确定性与漂移}，$\SE(3)$ 表达不出尺度变化，需显式带一个尺度因子 $s$。相似变换 $\mathbf{p}'=s\mathbf{R}\mathbf{p}+\mathbf{t}$ 构成相似变换群
\begin{equation}
  \mathrm{Sim}(3)=\Big\{\mathbf{S}=\begin{bmatrix}s\mathbf{R}&\mathbf{t}\\\mathbf{0}^\top&1\end{bmatrix}\Big\},\qquad
  \mathfrak{sim}(3)\ni\boldsymbol{\zeta}=[\boldsymbol{\rho}^\top,\boldsymbol{\phi}^\top,\sigma]^\top\in\mathbb{R}^7,
  \label{eq:sim3}
\end{equation}
其李代数多一维尺度 $\sigma$，对应 $s=e^{\sigma}$；$\boldsymbol{\zeta}$ 的反对称化在旋转块上多了一项 $\sigma\mathbf{I}$：
\begin{equation}
  \boldsymbol{\zeta}^\wedge=\begin{bmatrix}\sigma\mathbf{I}+\boldsymbol{\phi}^\wedge&\boldsymbol{\rho}\\\mathbf{0}^\top&0\end{bmatrix}\in\mathbb{R}^{4\times4}.
  \label{eq:sim3-hat}
\end{equation}
指数映射形如
\begin{equation}
  \exp(\boldsymbol{\zeta}^\wedge)=\begin{bmatrix}e^{\sigma}\mathbf{R}&\mathbf{J}_s\boldsymbol{\rho}\\\mathbf{0}^\top&1\end{bmatrix},
  \qquad s=e^\sigma,\ \mathbf{R}=\exp(\boldsymbol{\phi}^\wedge),\ \mathbf{t}=\mathbf{J}_s\boldsymbol{\rho},
  \label{eq:sim3-exp}
\end{equation}
其中尺度版雅可比 $\mathbf{J}_s(\sigma,\boldsymbol{\phi})$ 比 $\SE(3)$ 的左雅可比 $\mathbf{J}$（\cref{eq:left-jacobian}）多耦合一维尺度：它由级数 $\mathbf{J}_s=\sum_{n\ge0}\frac{1}{(n+1)!}(\sigma\mathbf{I}+\boldsymbol{\phi}^\wedge)^n$ 归并而来（把 \cref{eq:sim3-hat} 的右上块 $\big(\sum_{n\ge0}\frac{1}{(n+1)!}(\boldsymbol{\zeta}^\wedge)^n\big)$ 的旋转生成元由 $\boldsymbol{\phi}^\wedge$ 换成 $\sigma\mathbf{I}+\boldsymbol{\phi}^\wedge$），故同时含 $\sigma$ 的指数项与 $\theta$ 的三角项。其完整闭式 \cref{eq:sim3-Js} 见下面的推导盒，本书出于自包含在此补全；它与 SE(3) 情形一样可写成“$\mathbf{I}$、$\mathbf{a}^\wedge$、$\mathbf{a}^\wedge\mathbf{a}^\wedge$ 三族基的系数组合”，只是系数升级为 $\sigma,\theta$ 的混合函数。

\begin{derivation}[$\mathrm{Sim}(3)$ 的尺度雅可比 $\mathbf{J}_s$ 完整闭式]
设 $\boldsymbol{\phi}=\theta\mathbf{a}$（$\theta=\|\boldsymbol{\phi}\|$，$\mathbf{a}$ 单位轴），$s=e^{\sigma}$。把 $\mathbf{a}^\wedge$ 的幂用 \cref{eq:a-identities} 降阶、再按 $\sigma$ 与 $\theta$ 双重展开归并，得三族基的系数 $A_s,B_s,C_s$（与 Sophus 的 \texttt{calcW}、Eade\cite{eade2017lie} 的 $W$ 矩阵同构，仅基的标号不同）：
\begin{equation}
  \mathbf{J}_s
  =C_s\,\mathbf{I}+A_s\,\mathbf{a}^\wedge+B_s\,\mathbf{a}^\wedge\mathbf{a}^\wedge,
  \label{eq:sim3-Js}
\end{equation}
\begin{align*}
  C_s&=\frac{e^{\sigma}-1}{\sigma},\\
  A_s&=\frac{\sigma\,e^{\sigma}\sin\theta+(1-e^{\sigma}\cos\theta)\,\theta}{\sigma^2+\theta^2},\\
  B_s&=C_s-\frac{(e^{\sigma}\cos\theta-1)\,\sigma+(e^{\sigma}\sin\theta)\,\theta}{\sigma^2+\theta^2}.
\end{align*}
\textbf{两类退化}（实现里务必分支，否则 $0/0$ 出 \texttt{NaN}）：
\begin{itemize}
  \item $\sigma\to0$（无尺度，退回 $\SE(3)$）：令 $e^{\sigma}\to1$ 代入，并用 $\lim_{\sigma\to0}\tfrac{e^\sigma-1}{\sigma}=1$，三系数取极限为
  \[
    C_s\to1,\qquad
    A_s\to\frac{(1-\cos\theta)\,\theta}{\theta^2}=\frac{1-\cos\theta}{\theta},\qquad
    B_s\to1-\frac{\sin\theta\cdot\theta}{\theta^2}=\frac{\theta-\sin\theta}{\theta}.
  \]
  代回 \cref{eq:sim3-Js} 并用 $\mathbf{a}^\wedge\mathbf{a}^\wedge=\mathbf{a}\mathbf{a}^\top-\mathbf{I}$，恰好还原 $\SE(3)$ 左雅可比 $\mathbf{J}=\tfrac{\sin\theta}{\theta}\mathbf{I}+(1-\tfrac{\sin\theta}{\theta})\mathbf{a}\mathbf{a}^\top+\tfrac{1-\cos\theta}{\theta}\mathbf{a}^\wedge$（\cref{eq:left-jacobian}）。
  \item $\theta\to0$（无旋转，纯缩放）：$A_s\to0$、$B_s\to0$，故 $\mathbf{J}_s\to C_s\mathbf{I}=\dfrac{e^{\sigma}-1}{\sigma}\mathbf{I}$（再令 $\sigma\to0$ 则 $\mathbf{J}_s\to\mathbf{I}$）。此时 $\mathbf{t}=\tfrac{e^\sigma-1}{\sigma}\boldsymbol{\rho}$，正是“沿直线缩放”的几何意义。
\end{itemize}
\end{derivation}

\noindent 对一个点 $\mathbf{p}$ 经相似变换 $\mathbf{S}\mathbf{p}$，扰动模型同样比导数模型简单：记 $\mathbf{S}\mathbf{p}$ 的前三维为 $\mathbf{q}$，其对 $7$ 维李代数 $\boldsymbol{\zeta}$ 扰动的雅可比为 $4\times7$ 矩阵
\begin{equation}
  \frac{\partial(\mathbf{S}\mathbf{p})}{\partial\boldsymbol{\zeta}}
  =\begin{bmatrix}\mathbf{I}&-\mathbf{q}^\wedge&\mathbf{q}\\\mathbf{0}^\top&\mathbf{0}^\top&0\end{bmatrix},
  \label{eq:sim3-pert}
\end{equation}
其中末列 $\mathbf{q}$ 对应尺度自由度 $\sigma$ 的导数——这是 $\SE(3)$ 所没有的“缩放方向”。$\mathrm{Sim}(3)$ 在单目回环纠正（如 ORB-SLAM 的尺度感知位姿图）中必备；其指数映射与雅可比的系统推导见 Eade\cite{eade2017lie} 与 Strasdat 的工作\cite{strasdat2012thesis}，本节闭式已与之对照核验。不涉及单目可跳过。

% ════════════════════════════════════════════════════════════════
\section{（选读）扩展位姿群 \texorpdfstring{$\mathrm{SE}_2(3)$}{SE2(3)}：旋转 + 位置 + 速度}
\label{sec:se23}

\paragraph{动机：惯导/VIO 需要把速度也“装进群”。}
$\SE(3)$ 只装得下姿态 $\mathbf{R}$ 与位置 $\mathbf{r}$。但 IMU 测的是\emph{加速度}，积分一次得\emph{速度} $\mathbf{v}$、再积分才得位置；于是 VIO/惯导的状态天然是“姿态 + 速度 + 位置”三件套。最朴素的做法是把 $\mathbf{v}$ 当成一个普通的 $\mathbb{R}^3$ 向量挂在 $\SE(3)$ 旁边，但这样\emph{割裂}了三者的几何耦合——尤其在做协方差传播与一致性分析时会丢掉关键结构。更优雅的做法（Barrau--Bonnabel\cite{barrau2017invariant}、Brossard 等\cite{brossard2021associating}、Hartley 等\cite{hartley2020contact}）是把三者打包进\emph{一个}更大的矩阵李群——\textbf{扩展位姿群} $\mathrm{SE}_2(3)$。这一节给出它的\emph{权威定义}（生成元、$\wedge/\vee$、指数/对数、伴随、与 $\SE(3)$ 的关系），全书凡涉及 $\mathrm{SE}_2(3)$（误差状态滤波章 \cref{ch:eskf} 的不变滤波、\cref{ch:imu}/\cref{ch:vio} 的扩展位姿预积分）均回指本节，不再各自重述。

\begin{definition}[扩展位姿群 $\mathrm{SE}_2(3)$ 与其李代数 $\mathfrak{se}_2(3)$]\label{def:se23}
\textbf{扩展位姿群}是 $5\times5$ 矩阵的集合
\begin{equation}
  \mathrm{SE}_2(3)=\left\{\mathbf{T}=\begin{bmatrix}\mathbf{R}&\mathbf{r}&\mathbf{v}\\\mathbf{0}^\top&1&0\\\mathbf{0}^\top&0&1\end{bmatrix}\in\mathbb{R}^{5\times5}\ \middle|\ \mathbf{R}\in\SO(3),\ \mathbf{r},\mathbf{v}\in\mathbb{R}^3\right\},
  \label{eq:se23-group}
\end{equation}
即在 $\SE(3)$ 的齐次矩阵右侧再\emph{并排}一列速度 $\mathbf{v}$（并把右下角扩成 $2\times2$ 单位块）。它关于矩阵乘法构成 $9$ 维矩阵李群。其\textbf{李代数} $\mathfrak{se}_2(3)$ 的元素是 $9$ 维向量 $\boldsymbol{\xi}=[\boldsymbol{\rho}^\top,\boldsymbol{\nu}^\top,\boldsymbol{\phi}^\top]^\top$（\textbf{位置分量 $\boldsymbol{\rho}$、速度分量 $\boldsymbol{\nu}$、旋转 $\boldsymbol{\phi}$；沿用本书“平移类在前、旋转在后”约定}），其 $\wedge$ 算子为
\begin{equation}
  \boldsymbol{\xi}^\wedge=\begin{bmatrix}\boldsymbol{\rho}\\\boldsymbol{\nu}\\\boldsymbol{\phi}\end{bmatrix}^{\!\wedge}
  =\begin{bmatrix}\boldsymbol{\phi}^\wedge&\boldsymbol{\rho}&\boldsymbol{\nu}\\\mathbf{0}^\top&0&0\\\mathbf{0}^\top&0&0\end{bmatrix}\in\mathbb{R}^{5\times5},
  \label{eq:se23-hat}
\end{equation}
逆算子 $(\cdot)^\vee$ 把该矩阵读回向量。（注意 $\wedge$ 再次被\emph{重载}：左上仍是反对称化 $\boldsymbol{\phi}^\wedge$，两个并排列分别直接放 $\boldsymbol{\rho},\boldsymbol{\nu}$，不反对称——与 $\se(3)$ 的右上块同理。）
\end{definition}

\begin{proposition}[$\mathrm{SE}_2(3)$ 是矩阵李群]\label{prop:se23-group}
$\mathrm{SE}_2(3)$ 满足群公理：单位元为 $\mathbf{I}_5$；乘法封闭
\[
  \mathbf{T}_1\mathbf{T}_2=\begin{bmatrix}\mathbf{R}_1\mathbf{R}_2&\mathbf{R}_1\mathbf{r}_2+\mathbf{r}_1&\mathbf{R}_1\mathbf{v}_2+\mathbf{v}_1\\\mathbf{0}^\top&1&0\\\mathbf{0}^\top&0&1\end{bmatrix}\in\mathrm{SE}_2(3);
\]
逆为
\begin{equation}
  \mathbf{T}^{-1}=\begin{bmatrix}\mathbf{R}^\top&-\mathbf{R}^\top\mathbf{r}&-\mathbf{R}^\top\mathbf{v}\\\mathbf{0}^\top&1&0\\\mathbf{0}^\top&0&1\end{bmatrix}.
  \label{eq:se23-inverse}
\end{equation}
\end{proposition}
\begin{proof}
结合律继承自矩阵乘法；上式乘积直接按 $5\times5$ 分块相乘得到，右下 $2\times2$ 单位块在乘法下保持，故结果仍是 \cref{eq:se23-group} 的形态（封闭）。逆元验证 $\mathbf{T}\mathbf{T}^{-1}$：旋转块 $\mathbf{R}\mathbf{R}^\top=\mathbf{I}$，位置列 $-\mathbf{R}\mathbf{R}^\top\mathbf{r}+\mathbf{r}=\mathbf{0}$，速度列 $-\mathbf{R}\mathbf{R}^\top\mathbf{v}+\mathbf{v}=\mathbf{0}$，得 $\mathbf{I}_5$。注意逆的位置/速度列\emph{不耦合}（各自只含本列），这是 $\mathrm{SE}_2(3)$ 比一般群干净的地方。
\end{proof}

\begin{theorem}[$\mathrm{SE}_2(3)$ 指数映射]\label{thm:se23-exp}
$\boldsymbol{\xi}=[\boldsymbol{\rho};\boldsymbol{\nu};\boldsymbol{\phi}]\in\mathfrak{se}_2(3)$ 的指数映射为
\begin{equation}
  \mathbf{T}=\exp(\boldsymbol{\xi}^\wedge)
  =\begin{bmatrix}\mathbf{R}&\mathbf{J}\boldsymbol{\rho}&\mathbf{J}\boldsymbol{\nu}\\\mathbf{0}^\top&1&0\\\mathbf{0}^\top&0&1\end{bmatrix},
  \qquad \mathbf{R}=\exp(\boldsymbol{\phi}^\wedge),\quad \mathbf{J}=\mathbf{J}_l(\boldsymbol{\phi}),
  \label{eq:se23-exp}
\end{equation}
其中 $\mathbf{J}_l$ 正是 $\SO(3)$ 的\emph{左雅可比}（\cref{eq:left-jacobian}/\cref{eq:Jl}）。即真实位置 $\mathbf{r}=\mathbf{J}\boldsymbol{\rho}$、真实速度 $\mathbf{v}=\mathbf{J}\boldsymbol{\nu}$，二者被\emph{同一个} $\mathbf{J}$ 作用。对数映射为其逆：$\boldsymbol{\phi}=\mathrm{Log}(\mathbf{R})$、$\boldsymbol{\rho}=\mathbf{J}^{-1}\mathbf{r}$、$\boldsymbol{\nu}=\mathbf{J}^{-1}\mathbf{v}$。映射满射非单射（同 \cref{prop:surjective}）。
\end{theorem}
\begin{proof}
按 \cref{eq:se23-hat} 的 $5\times5$ 分块算 $\exp$。注意 $\boldsymbol{\xi}^\wedge$ 的幂只把左上 $\boldsymbol{\phi}^\wedge$ 自乘、把两个并排列各自左乘 $\boldsymbol{\phi}^\wedge$，而右下 $2\times2$ 零块自乘恒为零：
\[
  (\boldsymbol{\xi}^\wedge)^n=\begin{bmatrix}(\boldsymbol{\phi}^\wedge)^n&(\boldsymbol{\phi}^\wedge)^{n-1}\boldsymbol{\rho}&(\boldsymbol{\phi}^\wedge)^{n-1}\boldsymbol{\nu}\\\mathbf{0}^\top&0&0\\\mathbf{0}^\top&0&0\end{bmatrix}\quad(n\ge1).
\]
代入 $\exp(\boldsymbol{\xi}^\wedge)=\sum_{n\ge0}(\boldsymbol{\xi}^\wedge)^n/n!$：左上块归并为 $\mathbf{R}=\exp(\boldsymbol{\phi}^\wedge)$；位置列为 $\big(\sum_{n\ge0}\tfrac{1}{(n+1)!}(\boldsymbol{\phi}^\wedge)^n\big)\boldsymbol{\rho}=\mathbf{J}\boldsymbol{\rho}$（与 \cref{thm:se3-exp} 同一求和，故同一 $\mathbf{J}$），速度列同样为 $\mathbf{J}\boldsymbol{\nu}$。两列共用 $\mathbf{J}$ 正因它们在 \cref{eq:se23-hat} 中处于对称地位。
\end{proof}

\begin{insight}[$\mathrm{SE}_2(3)$ 为何“几乎就是两份 $\SE(3)$”]
比较 \cref{eq:se23-exp} 与 $\SE(3)$ 指数（\cref{thm:se3-exp}）：$\mathrm{SE}_2(3)$ 不过是让\emph{同一个旋转} $\mathbf{R}$ 同时搬运两个平移类量（位置 $\mathbf{r}=\mathbf{J}\boldsymbol{\rho}$、速度 $\mathbf{v}=\mathbf{J}\boldsymbol{\nu}$），且共用同一左雅可比 $\mathbf{J}$。所以本章为 $\SE(3)$ 建立的一切机器（Rodrigues、$\mathbf{J}$ 闭式、$\boldsymbol{\rho}\neq\mathbf{t}$ 的陷阱、最小多项式降阶 \cref{thm:minimal-poly}）几乎\emph{原样}迁移——只是把“一列平移”变成“两列平移”。这也是为什么遇到 $\mathrm{SE}_2(3)$ 不必从零推导：识别出它的生成元结构，闭式的“形状”立即可知。
\end{insight}

\paragraph{伴随 $\mathrm{Ad}(\mathbf{T})$ 与小伴随。}
$\mathrm{SE}_2(3)$ 的群伴随是 $9\times9$ 矩阵（按 $\boldsymbol{\xi}=[\boldsymbol{\rho};\boldsymbol{\nu};\boldsymbol{\phi}]$ 排序），同样满足移指数恒等式 $\mathbf{T}\,\mathrm{Exp}(\boldsymbol{\xi})=\mathrm{Exp}(\mathrm{Ad}(\mathbf{T})\boldsymbol{\xi})\,\mathbf{T}$（与 \cref{eq:Ad-move} 同形）：
\begin{equation}
  \mathrm{Ad}(\mathbf{T})=\begin{bmatrix}\mathbf{R}&\mathbf{0}&\mathbf{r}^\wedge\mathbf{R}\\\mathbf{0}&\mathbf{R}&\mathbf{v}^\wedge\mathbf{R}\\\mathbf{0}&\mathbf{0}&\mathbf{R}\end{bmatrix}\in\mathbb{R}^{9\times9},
  \qquad
  \boldsymbol{\xi}^\curlywedge=\begin{bmatrix}\boldsymbol{\phi}^\wedge&\mathbf{0}&\boldsymbol{\rho}^\wedge\\\mathbf{0}&\boldsymbol{\phi}^\wedge&\boldsymbol{\nu}^\wedge\\\mathbf{0}&\mathbf{0}&\boldsymbol{\phi}^\wedge\end{bmatrix}\in\mathbb{R}^{9\times9},
  \label{eq:se23-adjoint}
\end{equation}
其中 $\boldsymbol{\xi}^\curlywedge=\mathrm{ad}(\boldsymbol{\xi})$ 为小伴随，仍有 $\mathrm{Ad}(\exp(\boldsymbol{\xi}^\wedge))=\exp(\boldsymbol{\xi}^\curlywedge)$。$\mathrm{Ad}(\mathbf{T})$ 的两行 $\mathbf{r}^\wedge\mathbf{R}$、$\mathbf{v}^\wedge\mathbf{R}$ 正是位置/速度对旋转扰动的耦合，结构与 $\SE(3)$ 伴随（\cref{eq:Ad}）的 $\mathbf{t}^\wedge\mathbf{R}$ 块一脉相承。$9\times9$ 左雅可比 $\boldsymbol{\mathcal{J}}_l(\boldsymbol{\xi})$ 是分块上三角，对角为三个 $\mathbf{J}_l(\boldsymbol{\phi})$、两个耦合块与 $\SE(3)$ 的 $\mathbf{Q}_l$ 块（\cref{eq:se3-6x6,deriv:Ql}）同形：位置块 $\mathbf{Q}_l(\boldsymbol{\rho},\boldsymbol{\phi})$、速度块 $\mathbf{Q}_l(\boldsymbol{\nu},\boldsymbol{\phi})$ 共用同一闭式（只是把平移分量分别换成 $\boldsymbol{\rho},\boldsymbol{\nu}$）。

\begin{note}
\textbf{记号对照。}\;
Barfoot\cite{barfoot2024state} 在惯导一章用 $\mathrm{SE}_2(3)$ 并以 $\mathbf{C}$ 记旋转、$\mathbf{T}$ 记群元素，且\emph{唯独}在 $\mathrm{SE}_2(3)$ 的 IMU 运动模型处采用\emph{右扰动}（与本书全局主线一致）。本书统一 $\mathbf{R}$ 记旋转、右扰动 $\mathbf{T}\,\mathrm{Exp}(\delta\boldsymbol{\xi})$、$\boldsymbol{\xi}=[\boldsymbol{\rho};\boldsymbol{\nu};\boldsymbol{\phi}]$。文献中 $\mathrm{SE}_2(3)$ 亦称“双重直接等距群（double direct isometry group）”或径称“扩展位姿（extended pose）”。Barrau--Bonnabel 原始工作\cite{barrau2017invariant} 即在此群上建立不变 EKF。\textbf{分量排序提醒}：本书与 Barfoot 把\emph{旋转放最后}（$\boldsymbol{\xi}=[\boldsymbol{\rho};\boldsymbol{\nu};\boldsymbol{\phi}]$），故 \cref{eq:se23-adjoint} 的伴随是分块\emph{上}三角；而 InEKF 文献（Barrau--Bonnabel、Hartley 等\cite{hartley2020contact}）多把\emph{旋转放最前}（$[\boldsymbol{\phi};\dots]$），同一伴随写成分块\emph{下}三角——只是基排列不同的同一线性映射，对照外部论文时务必先核对排序。\textbf{还须当心位置/速度两列的顺序}：本书 \cref{eq:se23-group} 取“位置 $\mathbf{r}$ 列在前、速度 $\mathbf{v}$ 列在后”（与平移类分量序 $\boldsymbol{\rho}$ 在 $\boldsymbol{\nu}$ 前一致），\cref{ch:eskf} 的 InEKF 紧凑版（\cref{paper:inekf}）沿用此序；而 Hartley 等部分 InEKF 文献把\emph{速度列放在位置列之前}，对照时同样须先核对，两种排法只是基的置换。
\end{note}

\paragraph{为何要费这个劲：与不变滤波（InEKF）的关系。}
把状态放进 $\mathrm{SE}_2(3)$ 不是为了好看，而是为了\emph{一致性}。惯导/VIO 的连续运动学（姿态由陀螺、速度由加速度计 + 重力、位置由速度）在 $\mathrm{SE}_2(3)$ 上恰好构成 Barrau--Bonnabel 所称的\textbf{群仿射（group-affine）}动力学。群仿射系统有一个关键性质（Barrau--Bonnabel 的 \emph{log-linear} 定理\cite{barrau2017invariant}）：用\emph{不变误差}（右不变 $\mathbf{e}=\hat{\mathbf{T}}\mathbf{T}^{-1}$ 或左不变 $\mathbf{T}^{-1}\hat{\mathbf{T}}$，见 \cref{sec:lie-symmetry}）时，\emph{无噪声}误差动力学严格满足自治线性方程 $\dot{\boldsymbol{\eta}}=\mathbf{A}_t\boldsymbol{\eta}$，且其系统矩阵 $\mathbf{A}_t$ \emph{只依赖输入（IMU 读数），不依赖当前状态估计或轨迹}——这正是标准 EKF 做不到的（标准 EKF 的 $\mathbf{A}$ 含状态，估计有偏时线性化点错位，导致虚假信息增益、一致性退化）。误差动力学与状态解耦，使 InEKF 能给出与轨迹无关的收敛性/稳定性保证（含过程噪声时该自治性退化为近似，仍优于标准 EKF）。$\mathrm{SE}_2(3)$ 就是让“IMU 运动学 = 群仿射”这件事成立的那个群。完整的 InEKF 滤波推导见 \cref{ch:eskf}；这里只交代它为什么\emph{必须}站在 $\mathrm{SE}_2(3)$ 上，以及对称性的群论语言（下一节）如何支撑“线性化误差与状态无关”这一论断。

\begin{pitfall}[概念误区：把 $\mathbf{v}$ 当普通 $\mathbb{R}^3$ 状态拼在 $\SE(3)$ 旁]
“速度不就是个三维向量吗，直接和 $\SE(3)$ 并列放进状态向量不就行了？”——能跑，但会\emph{丢掉群仿射结构}：速度在重力/旋转下的耦合不再由群乘法自动携带，误差动力学重新变得依赖状态，InEKF 的一致性优势随之消失。$\mathrm{SE}_2(3)$ 的价值恰在于：让 $\mathbf{R},\mathbf{r},\mathbf{v}$ 的耦合成为\emph{群结构的一部分}，而非事后手工拼接。偏置 $\mathbf{b}$ 仍以普通向量附在群旁（biases 破坏严格群仿射，对应工程上的“不完美 InEKF / imperfect InEKF”，\cref{ch:eskf} 详述）。
\end{pitfall}

\begin{practice}
\begin{enumerate}
  \item 验证 \cref{prop:se23-group}：直接做 $5\times5$ 分块乘法核对 $\mathrm{SE}_2(3)$ 的封闭性与 \cref{eq:se23-inverse} 的逆（这正是 Barfoot 习题 8.21）。
  \item 由 \cref{thm:se23-exp} 证明：当 $\boldsymbol{\phi}\to\mathbf{0}$ 时 $\mathbf{J}\to\mathbf{I}$，故 $\mathbf{r}\to\boldsymbol{\rho}$、$\mathbf{v}\to\boldsymbol{\nu}$——李代数分量与真实位置/速度仅在零旋转处重合（对照 $\SE(3)$ 的 $\mathbf{t}=\mathbf{J}\boldsymbol{\rho}$）。
  \item 验证小伴随 \cref{eq:se23-adjoint} 满足 $\boldsymbol{\xi}^\curlywedge\mathbf{x}=(\boldsymbol{\xi}^\wedge\mathbf{x}^\wedge-\mathbf{x}^\wedge\boldsymbol{\xi}^\wedge)^\vee$（即它确是 $\mathfrak{se}_2(3)$ 李括号的矩阵表示）。
\end{enumerate}
\end{practice}

% ════════════════════════════════════════════════════════════════
\section{（选读）对称、不变与等变（通往 InEKF/EqF）}
\label{sec:lie-symmetry}

\paragraph{动机：换个参考系，答案该不该变？}
做定位时，你可以把地图路标的坐标表达在任意一个静止参考系里。直觉上，机器人估出来的位姿\emph{相对地图}应当一样——换参考系只是“换了把尺子量同一件事”。本节用群论语言把这条直觉精确化：群作用、\emph{不变性}与\emph{等变性}。它们既解释了“为何状态估计问题不依赖参考系选择”，也为上一节 InEKF“线性化误差与状态无关”的一致性论断提供了概念地基\cite{barfoot2024state}。这是一节概念性的收尾（选读），不引入新的闭式，但点透李群在估计中“好用”的深层原因——\emph{连续对称性}。

\paragraph{群作用。}
李群 $G$ 在集合 $X$ 上的（左）\emph{作用}是映射 $G\times X\to X$，记 $g\cdot x$，满足 $\mathcal{E}\cdot x=x$（单位元不动）与 $g_1\cdot(g_2\cdot x)=(g_1g_2)\cdot x$（与群乘法相容）。本章最常见的作用就是 $\SE(3)$ 经 $\mathbf{T}\mathbf{p}$ 作用在齐次点 $\mathbf{p}$ 上（刚体把点搬来搬去）。

\begin{definition}[不变性]\label{def:invariance}
称函数 $f$ \emph{在群 $G$ 作用下不变}（$G$ 是 $f$ 的一个\emph{对称}），若
\begin{equation}
  f(\mathbf{T}\mathbf{p})=f(\mathbf{p}),\qquad\forall\,\mathbf{T}\in\SE(3).
  \label{eq:invariance}
\end{equation}
即先用 $\mathbf{T}$ 变换输入，$f$ 的值不变。
\end{definition}

\begin{example}[定位代价的最优值是不变量]
考虑把路标观测对齐到地图的定位代价的\emph{最优值}（这是 \cref{ch:slam_est}/\cref{ch:vo} 定位问题的雏形）：
\begin{equation}
  J^\star(\mathbf{p}_j)=\min_{\mathbf{T}_1\in\SE(3)}\tfrac12\sum_j\big(\mathbf{y}_j-\mathbf{T}_1\mathbf{p}_j\big)^\top\big(\mathbf{y}_j-\mathbf{T}_1\mathbf{p}_j\big),
  \label{eq:loc-cost}
\end{equation}
其中 $\mathbf{p}_j$ 是静止系下的路标、$\mathbf{y}_j$ 是机体系下的（含噪）观测、$\mathbf{T}_1$ 为待估位姿。若把全部路标改表达在另一参考系 $\mathbf{p}_j\to\mathbf{T}_{12}\mathbf{p}_j$，代入并令 $\mathbf{T}_2=\mathbf{T}_1\mathbf{T}_{12}$，则对 $\mathbf{T}_2$ 取 $\min$ 的范围不变（$\SE(3)$ 群封闭），于是
\begin{equation}
  J^\star(\mathbf{T}_{12}\mathbf{p}_j)=J^\star(\mathbf{p}_j).
  \label{eq:loc-cost-invariant}
\end{equation}
\emph{最小可达代价不依赖路标表达在哪个参考系}——$J^\star$ 在 $\SE(3)$ 作用下不变，$\SE(3)$ 是它的对称。
\end{example}

\begin{definition}[等变性]\label{def:equivariance}
最优\emph{解}本身一般会随参考系改变，但改变方式\emph{可预测}。称矩阵值函数 $\mathbf{F}(\mathbf{p})$ 为\emph{右等变}，若
\begin{equation}
  \mathbf{F}(\mathbf{T}\mathbf{p})=\mathbf{F}(\mathbf{p})\,\mathbf{T};\qquad\text{\emph{左等变}：}\ \mathbf{F}(\mathbf{T}\mathbf{p})=\mathbf{T}\,\mathbf{F}(\mathbf{p}).
  \label{eq:equivariance}
\end{equation}
两种情形都称 $\SE(3)$ 是 $\mathbf{F}$ 的一个对称。
\end{definition}

\begin{example}[最优位姿是等变量]
\cref{eq:loc-cost} 取 $\arg\min$ 得最优位姿 $\mathbf{T}^\star(\mathbf{p}_j)$。同样把路标换系 $\mathbf{p}_j\to\mathbf{T}_{12}\mathbf{p}_j$，由上面的换元立得
\begin{equation}
  \mathbf{T}^\star(\mathbf{T}_{12}\mathbf{p}_j)=\mathbf{T}^\star(\mathbf{p}_j)\,\mathbf{T}_{12}.
  \label{eq:opt-pose-equivariant}
\end{equation}
最优位姿\emph{确实}随地图参考系变化，但只差一个右乘 $\mathbf{T}_{12}$——换言之，在一个系里解完，乘个已知变换就得到另一个系里的解。\emph{在哪个静止系建问题无所谓}，因为答案能轻松换算。这正是右等变（\cref{eq:equivariance}）。
\end{example}

\begin{definition}[不变误差]\label{def:invariant-error}
做滤波/观测器时，我们希望误差也以群为对称。设 $\mathbf{T}_{vi}$ 为待估位姿（$v$ 机体系、$i$ 惯性系）、$\tilde{\mathbf{T}}_{vi}$ 为某传感器对它的量测，定义 $\se(3)$ 中的误差：
\begin{align}
  \text{动系（左）不变误差：}\quad \mathbf{e}_v&=\mathrm{Log}\big(\mathbf{T}_{vi}\tilde{\mathbf{T}}_{vi}^{-1}\big),\label{eq:inv-err-moving}\\
  \text{静系（右）不变误差：}\quad \mathbf{e}_i&=\mathrm{Log}\big(\mathbf{T}_{vi}^{-1}\tilde{\mathbf{T}}_{vi}\big).\label{eq:inv-err-stationary}
\end{align}
二者都是\emph{相对量}，故对静止参考系的选择不变。
\end{definition}

\begin{insight}[对称性如何给 InEKF “线性化误差与状态无关”兜底]
把 \cref{def:invariant-error} 的不变误差用在 $\mathrm{SE}_2(3)$（\cref{sec:se23}）上的群仿射惯导动力学，对称性保证了误差是\emph{相对量}——它在群的左（或右）平移下不变。Barrau--Bonnabel\cite{barrau2017invariant} 正是据此证明：群仿射 + 不变误差 $\Rightarrow$ 误差动力学自治（系统矩阵不含状态），即上一节的 log-linear 性质。\emph{“换参考系答案可换算”（等变）与“误差是相对量”（不变）这两条对称性，就是 InEKF 一致性优势的群论根源。}标准 EKF 在原始坐标里线性化、不尊重这层对称，才会出现依赖估计的虚假信息增益。Bonnabel 等\cite{bonnabel2008symmetry} 最早把这套“对称保持观测器”思想引入；Barrau--Bonnabel\cite{barrau2017invariant} 据此提出 InEKF；Mahony--Trumpf\cite{mahony2021equivariant} 进一步提出\emph{等变滤波}（EqF），让任意李群上的运动学系统都能利用等变性设计滤波。InEKF/EqF 的完整滤波推导见 \cref{ch:eskf}，本节只给其概念地基。
\end{insight}

\begin{note}
\textbf{“左/右不变”命名的坑.}\;
\cref{eq:inv-err-moving,eq:inv-err-stationary} 有时分别叫“左不变/右不变误差”，但这套叫法\emph{易混}——左右取决于你估的是 $\mathbf{T}_{vi}$ 还是 $\mathbf{T}_{iv}$。本书随 Barfoot\cite{barfoot2024state} 改称\emph{动系不变误差}（\cref{eq:inv-err-moving}）与\emph{静系不变误差}（\cref{eq:inv-err-stationary}），以系（frame）而非左右来命名，避免歧义。
\end{note}

\begin{practice}
\begin{enumerate}
  \item 补全 \cref{eq:loc-cost-invariant} 的换元细节：写出 $\mathbf{T}_2=\mathbf{T}_1\mathbf{T}_{12}$ 代换后为何 $\min$ 的取值范围不变（用 $\SE(3)$ 群封闭）。
  \item 由 \cref{eq:opt-pose-equivariant} 说明：若在系 $A$ 解得最优位姿，如何\emph{无需重新优化}得到系 $B$（$\mathbf{p}^B=\mathbf{T}_{BA}\mathbf{p}^A$）下的最优位姿。
  \item （概念）解释为何“不变误差是相对量”这一点，使其对静止系的选择免疫——并把它与 InEKF 的一致性优势联系起来。
\end{enumerate}
\end{practice}

% ════════════════════════════════════════════════════════════════
\section{代码：Sophus 实践}
\label{sec:sophus}

Sophus（基于 Eigen 的头文件库）实现了 $\SO(3),\SE(3),\mathrm{Sim}(3)$。下面演示构造、指数/对数、$\wedge/\vee$ 与\emph{右扰动}更新（注意与本书约定一致地写成\textbf{右乘}）。

\begin{codebox}[C++]{Sophus：构造、exp/log、右扰动更新}
#include <sophus/se3.hpp>
#include <Eigen/Core>
// 沿 Z 轴 90 度
Eigen::Matrix3d R = Eigen::AngleAxisd(M_PI/2, Eigen::Vector3d::UnitZ()).toRotationMatrix();
Sophus::SO3d SO3_R(R);                 // 由旋转矩阵构造（也可 SO3d(q) 由四元数）
Eigen::Vector3d so3 = SO3_R.log();     // 对数映射 -> 李代数(3维)
Eigen::Matrix3d hat = Sophus::SO3d::hat(so3);   // ^ : 向量 -> 反对称阵

// 右扰动更新（本书主线）：X <- X * Exp(delta)
Eigen::Vector3d delta(1e-4, 0, 0);
Sophus::SO3d SO3_updated = SO3_R * Sophus::SO3d::exp(delta);   // 右乘！

// SE(3) 同理；注意 Sophus 的 se(3) 顺序是“平移在前、旋转在后”，与本书一致
Eigen::Vector3d t(1, 0, 0);
Sophus::SE3d SE3(R, t);
Eigen::Matrix<double,6,1> xi = SE3.log();        // [rho; phi]
Eigen::Matrix<double,6,1> dxi; dxi.setZero(); dxi(0)=1e-4;
Sophus::SE3d SE3_updated = SE3 * Sophus::SE3d::exp(dxi);       // 右乘更新
\end{codebox}

\begin{note}
不少实现（如视觉SLAM十四讲\cite{gaoxiang2019slam14} 配套代码）采用\emph{左}乘更新 $\mathrm{Exp}(\delta)\,\mathbf{X}$；本书统一\emph{右}乘 $\mathbf{X}\,\mathrm{Exp}(\delta)$ 以对齐右扰动约定。容器存放需对齐分配器：\texttt{std::vector<Sophus::SE3d, Eigen::aligned\_allocator<Sophus::SE3d>>}。轨迹误差的核心是一行 \texttt{(T\_gt.inverse()*T\_est).log().norm()}，即 $\|\mathrm{Log}(\bar{\mathbf{T}}_{\text{gt}}^{-1}\mathbf{T}_{\text{est}})\|$——\emph{用李代数把位姿误差变成可求范数的向量}。
\end{note}

\subsection{应用实例：用李代数评估轨迹误差}
\label{sec:traj-error}

前面把“位姿误差”一笔带过；下面用一个经典实例把它讲透\cite{gaoxiang2019slam14}——它也是本章理论最直接、最常用的落地：\emph{如何度量一条估计轨迹与真实轨迹的差距}。这是 SLAM/里程计论文里评测精度的标准动作，而它之所以能干净地写出来，正是因为我们有了\emph{对数映射}这把“把位姿差变成向量”的钥匙。

\paragraph{问题。}设估计轨迹 $\{\mathbf{T}_{\mathrm{est},i}\}$ 与真值轨迹 $\{\mathbf{T}_{\mathrm{gt},i}\}$（$i=1,\dots,N$，真值常由动捕等更高精度系统给出）。难点在于：位姿是 $\SE(3)$ 元素，\emph{不能}直接相减求误差。出路是先求两位姿的\emph{相对变换} $\mathbf{T}_{\mathrm{gt},i}^{-1}\mathbf{T}_{\mathrm{est},i}$（完全一致时它是单位阵 $\mathbf{I}$），再用 $\mathrm{Log}$ 把这个“偏离单位阵多远”的相对位姿压成一个六维向量，其范数即单帧误差。

\paragraph{绝对轨迹误差（ATE）。}对全程取均方根（RMSE），同时含旋转与平移：
\begin{equation}
  \mathrm{ATE}_{\text{all}}=\sqrt{\frac{1}{N}\sum_{i=1}^{N}\big\|\mathrm{Log}\big(\mathbf{T}_{\mathrm{gt},i}^{-1}\mathbf{T}_{\mathrm{est},i}\big)\big\|_2^2}.
  \label{eq:ate-all}
\end{equation}
若只关心平移漂移（如某些里程计文献），取相对位姿的平移部分 $\mathrm{trans}(\cdot)$：
\begin{equation}
  \mathrm{ATE}_{\text{trans}}=\sqrt{\frac{1}{N}\sum_{i=1}^{N}\big\|\mathrm{trans}\big(\mathbf{T}_{\mathrm{gt},i}^{-1}\mathbf{T}_{\mathrm{est},i}\big)\big\|_2^2}.
  \label{eq:ate-trans}
\end{equation}
注意旋转误差会沿轨迹\emph{累积}成平移误差，故二者在实践中都常用。

\paragraph{相对位姿误差（RPE）。}ATE 衡量与真值的\emph{全局}对齐；RPE 则衡量\emph{每段相对运动}的准确性，对 $i\to i+\Delta t$ 这一小段做差：
\begin{equation}
  \mathrm{RPE}_{\text{all}}=\sqrt{\frac{1}{N-\Delta t}\sum_{i=1}^{N-\Delta t}\Big\|\mathrm{Log}\Big(\big(\mathbf{T}_{\mathrm{gt},i}^{-1}\mathbf{T}_{\mathrm{gt},i+\Delta t}\big)^{-1}\big(\mathbf{T}_{\mathrm{est},i}^{-1}\mathbf{T}_{\mathrm{est},i+\Delta t}\big)\Big)\Big\|_2^2},
  \label{eq:rpe-all}
\end{equation}
同样可只取平移部分得 $\mathrm{RPE}_{\text{trans}}$。直觉：RPE 把每段真值相对运动与对应估计相对运动相比，对累积漂移不敏感，更能反映局部一致性。

\begin{insight}
四个指标 \cref{eq:ate-all,eq:ate-trans,eq:rpe-all} 都建立在同一动作上：\textbf{相对位姿 $\to$ $\mathrm{Log}$ $\to$ 取范数}。没有对数映射，“两个位姿差多少”根本无从定义为一个标量。这就是本章理论最朴素也最高频的用途——它把抽象的“流形上两点的距离”落成了一行可执行代码。
\end{insight}

下面用 Sophus 实现 ATE（RMSE）。核心仍是 \texttt{(p2.inverse()*p1).log().norm()} 一行；读取轨迹按 TUM 格式 \texttt{time tx ty tz qx qy qz qw}：

\begin{codebox}[C++]{评估轨迹误差：ATE 的 RMSE（trajectoryError.cpp 核心）}
#include <sophus/se3.hpp>
#include <Eigen/Core>
#include <fstream>
#include <vector>
using TrajectoryType =
    std::vector<Sophus::SE3d, Eigen::aligned_allocator<Sophus::SE3d>>;

// 读取 TUM 格式轨迹: time tx ty tz qx qy qz qw
TrajectoryType ReadTrajectory(const std::string &path) {
  std::ifstream fin(path);
  TrajectoryType traj;
  double time, tx, ty, tz, qx, qy, qz, qw;
  while (fin >> time >> tx >> ty >> tz >> qx >> qy >> qz >> qw)
    traj.emplace_back(Eigen::Quaterniond(qw, qx, qy, qz),  // (w,x,y,z)
                      Eigen::Vector3d(tx, ty, tz));
  return traj;
}

double ComputeATE(const TrajectoryType &est, const TrajectoryType &gt) {
  double sse = 0.0;                       // sum of squared errors
  for (size_t i = 0; i < est.size(); ++i) {
    // T_gt^-1 * T_est -> Log -> norm : single-frame error (6-vector norm)
    double e = (gt[i].inverse() * est[i]).log().norm();
    sse += e * e;
  }
  return std::sqrt(sse / static_cast<double>(est.size()));   // RMSE
}
\end{codebox}

\begin{note}
\textbf{实现要点.}\;(1) 轨迹容器必须用 \texttt{Eigen::aligned\_allocator} 以满足 Eigen 的内存对齐；(2) 此例的真值/估计已做过\emph{时间对齐}与\emph{外参预估}（坐标系对齐）——这两步本身是后续章节内容，工程里务必先做，否则 ATE 数值无意义；(3) 在该配套数据集上此例得到 RMSE $\approx 2.207$\cite{gaoxiang2019slam14}。把 \texttt{.log().norm()} 换成只取平移分量即得 $\mathrm{ATE}_{\text{trans}}$。
\end{note}

\paragraph{小结与过渡。}至此，本章的理论（指数/对数、扰动、伴随）已在“位姿更新”和“轨迹评估”两件最常见的工程任务上各落一子。接下来进入选读的 $\mathrm{Sim}(3)$ 与本章附录，正文主干到此完整；后续 \cref{ch:nlopt,ch:vo} 将把这些雅可比真正放进优化器里求解。

% ════════════════════════════════════════════════════════════════
\section{本章推导附录}
\label{sec:lie-appendix}

本章主要证明（$\so(3)$ 的引出、Rodrigues 公式、左雅可比 $\mathbf{J}$、左/右扰动求导）均已在正文相应定理处给出。此处收录若干\emph{较冗长或较进阶}的补充：两处补充闭式（$\mathbf{Q}_l$ 块、四阶协方差算子）；把正文反复用到的“降阶恒等式”统一到代数源头的\emph{特征/Jordan 分解与三个最小多项式}，及由其得到的 $\SE(3)$ 指数/伴随\emph{直接整体级数式}（\cref{deriv:direct-series}）；伴随的导数恒等式（\cref{deriv:Ad-deriv}）与运动学雅可比恒等式；以及点观测雅可比所依赖的\emph{齐次点算子 $\odot/\circledcirc$ 代数}（\cref{subsec:odot-algebra}）。均供正文按需取用，初读可跳过。

\begin{derivation}[$\SE(3)$ 的 $\mathbf{Q}_l$ 块闭式（Barfoot 式 8.91）]\label{deriv:Ql}
\cref{eq:se3-6x6} 中 $6\times6$ 左雅可比的耦合块 $\mathbf{Q}_l$ 由双重级数 $\mathbf{Q}_l=\sum_{n,m\ge0}\frac{(\boldsymbol{\phi}^\wedge)^n\boldsymbol{\rho}^\wedge(\boldsymbol{\phi}^\wedge)^m}{(n+m+2)!}$ 归并而得（$\theta=\|\boldsymbol{\phi}\|$）：
\begin{align*}
  \mathbf{Q}_l(\boldsymbol{\rho},\boldsymbol{\phi})
  &=\tfrac12\boldsymbol{\rho}^\wedge
   +\frac{\theta-\sin\theta}{\theta^3}\big(\boldsymbol{\phi}^\wedge\boldsymbol{\rho}^\wedge+\boldsymbol{\rho}^\wedge\boldsymbol{\phi}^\wedge+\boldsymbol{\phi}^\wedge\boldsymbol{\rho}^\wedge\boldsymbol{\phi}^\wedge\big)\\
  &\quad+\frac{\theta^2+2\cos\theta-2}{2\theta^4}\big(\boldsymbol{\phi}^\wedge\boldsymbol{\phi}^\wedge\boldsymbol{\rho}^\wedge+\boldsymbol{\rho}^\wedge\boldsymbol{\phi}^\wedge\boldsymbol{\phi}^\wedge-3\,\boldsymbol{\phi}^\wedge\boldsymbol{\rho}^\wedge\boldsymbol{\phi}^\wedge\big)\\
  &\quad+\frac{2\theta-3\sin\theta+\theta\cos\theta}{2\theta^5}\big(\boldsymbol{\phi}^\wedge\boldsymbol{\rho}^\wedge\boldsymbol{\phi}^\wedge\boldsymbol{\phi}^\wedge+\boldsymbol{\phi}^\wedge\boldsymbol{\phi}^\wedge\boldsymbol{\rho}^\wedge\boldsymbol{\phi}^\wedge\big).
\end{align*}
右版本 $\mathbf{Q}_r(\boldsymbol{\xi})=\mathbf{Q}_l(-\boldsymbol{\xi})$。常规 SLAM 优化旋转较小，取 $\boldsymbol{\mathcal{J}}\approx\mathbf{I}$ 即可；仅高精度协方差传播时用此式。
\end{derivation}

\begin{derivation}[位姿复合四阶协方差的算子（\cref{eq:fourth-order}）]
\cref{eq:fourth-order} 中的算子定义为（Barfoot\cite{barfoot2024state} §8.3，$(\cdot)^\curlywedge$ 为 \cref{eq:ad} 的小伴随）：
\[
  \boldsymbol{\mathcal{A}}_i=\mathbb{E}\!\big[\boldsymbol{\epsilon}_i^\curlywedge\boldsymbol{\epsilon}_i^\curlywedge\big],\qquad
  \boldsymbol{\mathcal{B}}=\mathbb{E}\!\big[\boldsymbol{\epsilon}_1^\curlywedge\,(\boldsymbol{\epsilon}_2'\boldsymbol{\epsilon}_2'^\top)\,\boldsymbol{\epsilon}_1^{\curlywedge\top}\big],\qquad
  \boldsymbol{\epsilon}_2'=\mathrm{Ad}(\bar{\mathbf{T}}_1)\boldsymbol{\epsilon}_2.
\]
该修正在扰动上为四阶、在协方差上为二阶；当两位姿不相关时与 sigma-point（无迹）法代数等价。
\end{derivation}

\subsection{（选读）特征 / Jordan 分解与三个最小多项式}
\label{subsec:minimal-poly}

\paragraph{动机：正文那些“降阶恒等式”从哪来？}本章正文多处靠一条恒等式把高阶幂\emph{降阶}，从而把无穷级数收敛成闭式：Rodrigues（\cref{thm:so3-exp}）与左雅可比（\cref{thm:se3-exp}）用 $(\mathbf{a}^\wedge)^3=-\mathbf{a}^\wedge$（即 \cref{eq:a-identities}）；$6\times6$ 雅可比 $\boldsymbol{\mathcal{J}}_l$、伴随级数 $\mathrm{Ad}=\exp(\mathrm{ad})$ 之所以是\emph{有限}次幂的组合，背后也各有一条降阶恒等式。这些恒等式并非互不相干的“技巧”，它们都是相应生成元的\textbf{最小多项式}——由特征值结构唯一决定。本小节把它们统一到代数源头：先做特征/Jordan 分解，读出特征值，再经 Cayley--Hamilton 写出最小多项式。

\paragraph{$\so(3)$ 的特征分解。}$\boldsymbol{\phi}^\wedge$ 的特征多项式为（直接展开 $3\times3$ 行列式，$\phi=\|\boldsymbol{\phi}\|$）
\begin{equation}
  \det(\lambda\mathbf{I}-\boldsymbol{\phi}^\wedge)=\lambda^3+\phi^2\lambda=\lambda(\lambda^2+\phi^2),
  \label{eq:so3-charpoly}
\end{equation}
故特征值为 $\{0,\,+i\phi,\,-i\phi\}$（实轴特征向量即转轴 $\mathbf{a}$，对应特征值 $0$，印证 $\boldsymbol{\phi}^\wedge\mathbf{a}=\mathbf{0}$）。$\boldsymbol{\phi}^\wedge$ 可对角化为 $\boldsymbol{\phi}^\wedge=\mathbf{V}\mathbf{D}\mathbf{V}^{H}$，$\mathbf{D}=\operatorname{diag}(i\phi,-i\phi,0)$，从而旋转矩阵
\begin{equation}
  \mathbf{R}=\exp(\boldsymbol{\phi}^\wedge)=\mathbf{V}\exp(\mathbf{D})\mathbf{V}^{H},\qquad
  \exp(\mathbf{D})=\operatorname{diag}(e^{i\phi},e^{-i\phi},1).
  \label{eq:so3-eigen}
\end{equation}
由此立得两条熟悉的不变量：$\operatorname{tr}\mathbf{R}=e^{i\phi}+e^{-i\phi}+1=2\cos\phi+1$（与 \cref{eq:so3-log} 反求转角一致）、$\det\mathbf{R}=e^{i\phi}e^{-i\phi}\cdot1=1$。

\begin{theorem}[三个最小多项式]\label{thm:minimal-poly}
$\SO(3)/\SE(3)$ 及其伴随的生成元满足如下\emph{最小多项式}（$\phi=\|\boldsymbol{\phi}\|$ 为转角）：
\begin{align}
  &\so(3):    &(\boldsymbol{\phi}^\wedge)^3+\phi^2\,\boldsymbol{\phi}^\wedge&=\mathbf{0},\label{eq:minpoly-so3}\\
  &\se(3):    &(\boldsymbol{\xi}^\wedge)^4+\phi^2\,(\boldsymbol{\xi}^\wedge)^2&=\mathbf{0},\label{eq:minpoly-se3}\\
  &\mathrm{ad}(\se(3)):&(\boldsymbol{\xi}^\curlywedge)^5+2\phi^2\,(\boldsymbol{\xi}^\curlywedge)^3+\phi^4\,\boldsymbol{\xi}^\curlywedge&=\mathbf{0}.\label{eq:minpoly-ad}
\end{align}
它们正是本章正文一切“级数收敛成闭式”的代数根据：\cref{eq:minpoly-so3} 即 \cref{eq:a-identities} 的 $(\mathbf{a}^\wedge)^3=-\mathbf{a}^\wedge$（取 $\phi=1$、$\boldsymbol{\phi}^\wedge=\mathbf{a}^\wedge$）；\cref{eq:minpoly-se3} 使 $\mathbf{T}=\exp(\boldsymbol{\xi}^\wedge)$ 与左雅可比 $\mathbf{J}$ 只含到 $(\boldsymbol{\xi}^\wedge)^3$ 的项；\cref{eq:minpoly-ad} 使伴随 $\mathrm{Ad}(\mathbf{T})=\exp(\boldsymbol{\xi}^\curlywedge)$ 与 $6\times6$ 雅可比 $\boldsymbol{\mathcal{J}}_l$ 只含到 $(\boldsymbol{\xi}^\curlywedge)^4$ 的项。
\end{theorem}
\begin{proof}
\emph{$\so(3)$.}\;由 Cayley--Hamilton 定理，矩阵满足自身的特征多项式 \cref{eq:so3-charpoly}，即 $(\boldsymbol{\phi}^\wedge)^3+\phi^2\boldsymbol{\phi}^\wedge=\mathbf{0}$；因三特征值互异，特征多项式即\emph{最小}多项式。这恰是 \cref{eq:a-identities}（$\phi=1$ 时 $(\mathbf{a}^\wedge)^3=-\mathbf{a}^\wedge$）。\par
\emph{$\se(3)$.}\;$\boldsymbol{\xi}^\wedge=\big[\begin{smallmatrix}\boldsymbol{\phi}^\wedge&\boldsymbol{\rho}\\\mathbf{0}^\top&0\end{smallmatrix}\big]$ 是 $4\times4$ 上三角分块，其特征值为 $\boldsymbol{\phi}^\wedge$ 的三个 $\{0,\pm i\phi\}$ 再加右下角的 $0$，即 $\{0,0,+i\phi,-i\phi\}$，特征多项式 $\det(\lambda\mathbf{I}-\boldsymbol{\xi}^\wedge)=\lambda^2(\lambda^2+\phi^2)$。但 $\lambda=0$ 的代数重数 $2$ 对应一个 $2\times2$ Jordan 块（一般 $\mathbf{a}^\top\boldsymbol{\rho}\neq0$ 时不可对角化，需 Jordan 分解 $\boldsymbol{\xi}^\wedge=\mathbf{U}\mathbf{E}\mathbf{U}^{-1}$），故\emph{最小}多项式仍含 $\lambda^2$：$\lambda^2(\lambda^2+\phi^2)$，即 \cref{eq:minpoly-se3}。验证用 Jordan 分解：$(\boldsymbol{\xi}^\wedge)^4+\phi^2(\boldsymbol{\xi}^\wedge)^2=\mathbf{U}\,\mathbf{E}^2(\mathbf{E}^2+\phi^2\mathbf{I})\,\mathbf{U}^{-1}$，因 $\mathbf{E}^2(\mathbf{E}^2+\phi^2\mathbf{I})=\mathbf{0}$（$2\times2$ 零 Jordan 块平方为零、$\pm i\phi$ 块被 $(\mathbf{E}^2+\phi^2\mathbf{I})$ 消去）而为零。\par
\emph{$\mathrm{ad}(\se(3))$.}\;$\boldsymbol{\xi}^\curlywedge=\big[\begin{smallmatrix}\boldsymbol{\phi}^\wedge&\boldsymbol{\rho}^\wedge\\\mathbf{0}&\boldsymbol{\phi}^\wedge\end{smallmatrix}\big]$（\cref{eq:ad}）是 $6\times6$ 分块上三角，特征值为 $\boldsymbol{\phi}^\wedge$ 的特征值各\emph{两份}：$\{0,0,+i\phi,+i\phi,-i\phi,-i\phi\}$，特征多项式 $\lambda^2(\lambda^2+\phi^2)^2$。Jordan 分析表明 $\lambda=0$ 的最大块为 $1\times1$（故最小多项式只保留\emph{一个} $\lambda$ 因子），而 $\pm i\phi$ 各保留二重：最小多项式为 $\lambda(\lambda^2+\phi^2)^2$，展开即 \cref{eq:minpoly-ad}。
\end{proof}

\begin{insight}[为什么“最小多项式”是闭式的总开关]
一个矩阵 $\mathbf{M}$ 的最小多项式次数 $=d$，意味着 $\mathbf{M}^d$ 可由 $\{\mathbf{I},\mathbf{M},\dots,\mathbf{M}^{d-1}\}$ 线性表出，于是\emph{任何}关于 $\mathbf{M}$ 的幂级数（$\exp$、$\mathbf{J}$、$\mathrm{Ad}$、$\boldsymbol{\mathcal{J}}$）都坍缩成这 $d$ 个基的有限组合，系数是把级数按基归并后的标量函数（正是 $\sin\phi/\phi$、$(1-\cos\phi)/\phi^2$ 这类）。这就是为什么 $\so(3)$ 的一切闭式都落在 $\{\mathbf{I},\boldsymbol{\phi}^\wedge,(\boldsymbol{\phi}^\wedge)^2\}$ 三族（$d=3$）、$\mathrm{Sim}(3)$ 的 $\mathbf{J}_s$（\cref{eq:sim3-Js}）也是三族——本章所有“神奇的闭式”本质是\emph{同一}件事：最小多项式把无穷降为有限。记住这条，遇到新李群（如 $\mathrm{SE}_2(3)$）也能照方抓药：先求特征值、定最小多项式次数，闭式的“基”和项数当即可知。
\end{insight}

\begin{derivation}[$\SE(3)$ 指数与伴随的\emph{直接整体}级数式（避开分块）]\label{deriv:direct-series}
\cref{thm:se3-exp} 把 $\mathbf{T}=\exp(\boldsymbol{\xi}^\wedge)$ 按 $4\times4$ 分块（左上 $\mathbf{R}$、右上 $\mathbf{J}\boldsymbol{\rho}$）求出，是\emph{结构清晰}的写法；但实现时若已有现成的 $\boldsymbol{\xi}^\wedge$（或 $6\times6$ 的 $\boldsymbol{\xi}^\curlywedge$）矩阵，更省事的是\emph{不拆块}、直接对整块旋量矩阵做有限项级数。\cref{thm:minimal-poly} 的两条最小多项式（\cref{eq:minpoly-se3,eq:minpoly-ad}）正好把无穷级数截断成有限项\cite{barfoot2024state}。

\textbf{$\SE(3)$ 指数（$4\times4$）。}用 \cref{eq:minpoly-se3} 即 $(\boldsymbol{\xi}^\wedge)^4=-\phi^2(\boldsymbol{\xi}^\wedge)^2$ 把四次及更高次幂逐级降到 $\{(\boldsymbol{\xi}^\wedge)^2,(\boldsymbol{\xi}^\wedge)^3\}$，再把 $\exp$ 级数按 $\mathbf{I},\boldsymbol{\xi}^\wedge,(\boldsymbol{\xi}^\wedge)^2,(\boldsymbol{\xi}^\wedge)^3$ 四族归并（系数与罗德里格斯同源：偶/奇次项各凑成 $\tfrac{1-\cos\phi}{\phi^2}$、$\tfrac{\phi-\sin\phi}{\phi^3}$）：
\begin{equation}
  \mathbf{T}=\exp(\boldsymbol{\xi}^\wedge)
  =\mathbf{I}+\boldsymbol{\xi}^\wedge
   +\frac{1-\cos\phi}{\phi^2}(\boldsymbol{\xi}^\wedge)^2
   +\frac{\phi-\sin\phi}{\phi^3}(\boldsymbol{\xi}^\wedge)^3,
  \qquad \phi=\|\boldsymbol{\phi}\|.
  \label{eq:se3-exp-direct}
\end{equation}
（核对：$(\boldsymbol{\xi}^\wedge)^2,(\boldsymbol{\xi}^\wedge)^3$ 的右上块经 \cref{eq:se23-hat} 类型的分块乘法回到 $\mathbf{J}\boldsymbol{\rho}$，与 \cref{thm:se3-exp} 完全一致；此式只是把同一结果写成\emph{整体}多项式。）

\textbf{伴随 $\boldsymbol{\mathcal{T}}=\mathrm{Ad}(\mathbf{T})$（$6\times6$）。}同理，\cref{sec:adjoint} 已得 $\boldsymbol{\mathcal{T}}=\exp(\boldsymbol{\xi}^\curlywedge)$。用 \cref{eq:minpoly-ad} 即 $(\boldsymbol{\xi}^\curlywedge)^5=-2\phi^2(\boldsymbol{\xi}^\curlywedge)^3-\phi^4\boldsymbol{\xi}^\curlywedge$ 把五次及更高次降到 $\{\boldsymbol{\xi}^\curlywedge,\dots,(\boldsymbol{\xi}^\curlywedge)^4\}$，归并得五项闭式
\begin{equation}
  \begin{aligned}
  \boldsymbol{\mathcal{T}}=\exp(\boldsymbol{\xi}^\curlywedge)
  =\mathbf{I}
  &+\frac{3\sin\phi-\phi\cos\phi}{2\phi}\,\boldsymbol{\xi}^\curlywedge
   +\frac{4-\phi\sin\phi-4\cos\phi}{2\phi^2}\,(\boldsymbol{\xi}^\curlywedge)^2\\
  &+\frac{\sin\phi-\phi\cos\phi}{2\phi^3}\,(\boldsymbol{\xi}^\curlywedge)^3
   +\frac{2-\phi\sin\phi-2\cos\phi}{2\phi^4}\,(\boldsymbol{\xi}^\curlywedge)^4.
  \end{aligned}
  \label{eq:Ad-series-direct}
\end{equation}
$6\times6$ 左雅可比 $\boldsymbol{\mathcal{J}}_l=\sum_{n\ge0}\tfrac{(\boldsymbol{\xi}^\curlywedge)^n}{(n+1)!}$（\cref{eq:se3-6x6}）同样被 \cref{eq:minpoly-ad} 截断到 $(\boldsymbol{\xi}^\curlywedge)^4$，其整体级数式只需把 \cref{eq:Ad-series-direct} 的标量系数换成把 $\tfrac1{n!}$ 改为 $\tfrac1{(n+1)!}$ 后的归并值（与 \cref{deriv:Ql} 的分块 $\mathbf{Q}_l$ 等价）。两式都\emph{不必处理分块}，便于直接用矩阵库求值；当转角小、又不追求高精度时仍可退回 $\boldsymbol{\mathcal{T}}\approx\mathbf{I}+\boldsymbol{\xi}^\curlywedge$、$\boldsymbol{\mathcal{J}}_l\approx\mathbf{I}$。
\end{derivation}

\begin{derivation}[$\mathrm{Ad}(\SE(3))$ 的导数恒等式（附录 B.1）]\label{deriv:Ad-deriv}
\cref{sec:adjoint} 的伴随 $\boldsymbol{\mathcal{T}}(\boldsymbol{\xi})=\mathrm{Ad}(\mathrm{Exp}(\boldsymbol{\xi}))$ 在协方差传播里常以“$\boldsymbol{\mathcal{T}}(\boldsymbol{\xi})\mathbf{x}$ 对 $\boldsymbol{\xi}$ 求导”的形式出现（$\mathbf{x}\in\mathbb{R}^6$ 为一个固定切向量）。这恰是 \cref{eq:right-pert-so3,eq:deriv-model} 那套“扰动求导”在 $6\times6$ 伴随上的翻版\cite{barfoot2024state}。沿第 $i$ 个坐标方向 $\mathbf{1}_i$ 取方向导数，并用 BCH 的左乘一阶式 \cref{eq:bch-se3-left}（$\mathrm{Exp}((\boldsymbol{\xi}+h\mathbf{1}_i)^\wedge)\approx\mathrm{Exp}((\boldsymbol{\mathcal{J}}_l h\mathbf{1}_i)^\wedge)\mathrm{Exp}(\boldsymbol{\xi}^\wedge)$）把扰动挪到左侧：
\begin{align*}
  \frac{\partial(\boldsymbol{\mathcal{T}}(\boldsymbol{\xi})\mathbf{x})}{\partial\xi_i}
  &=\lim_{h\to0}\frac{\mathrm{Ad}\big(\mathrm{Exp}((\boldsymbol{\xi}+h\mathbf{1}_i)^\wedge)\big)\mathbf{x}-\boldsymbol{\mathcal{T}}(\boldsymbol{\xi})\mathbf{x}}{h}
   =\lim_{h\to0}\frac{\big(\mathbf{I}+h(\boldsymbol{\mathcal{J}}_l\mathbf{1}_i)^\curlywedge\big)\boldsymbol{\mathcal{T}}(\boldsymbol{\xi})\mathbf{x}-\boldsymbol{\mathcal{T}}(\boldsymbol{\xi})\mathbf{x}}{h}\\
  &=(\boldsymbol{\mathcal{J}}_l\mathbf{1}_i)^\curlywedge\,\boldsymbol{\mathcal{T}}(\boldsymbol{\xi})\mathbf{x}
   =-\big(\boldsymbol{\mathcal{T}}(\boldsymbol{\xi})\mathbf{x}\big)^\curlywedge\boldsymbol{\mathcal{J}}_l\mathbf{1}_i,
\end{align*}
其中第二步用了 $\mathrm{Ad}(\mathrm{Exp}(\boldsymbol{\delta}))\approx\mathbf{I}+\boldsymbol{\delta}^\curlywedge$（小伴随是 $\mathrm{Ad}$ 在单位元的导数，\cref{eq:ad}），末步用了 $6$ 维“叉积”反交换 $\mathbf{a}^\curlywedge\mathbf{b}=-\mathbf{b}^\curlywedge\mathbf{a}$（$\boldsymbol{\mathcal{J}}_l\mathbf{1}_i$ 与 $\boldsymbol{\mathcal{T}}\mathbf{x}$ 互换）。把六列 $\mathbf{1}_i$ 横排即得
\begin{equation}
  \boxed{\;\frac{\partial\big(\boldsymbol{\mathcal{T}}(\boldsymbol{\xi})\,\mathbf{x}\big)}{\partial\boldsymbol{\xi}}
  =-\big(\boldsymbol{\mathcal{T}}(\boldsymbol{\xi})\,\mathbf{x}\big)^\curlywedge\,\boldsymbol{\mathcal{J}}_l(\boldsymbol{\xi})\;}
  \label{eq:Ad-deriv}
\end{equation}
（结构与点观测雅可比 \cref{eq:right-pert-se3} 的 $(\mathbf{T}\tilde{\mathbf{p}})^\odot\boldsymbol{\mathcal{J}}$ 完全平行：把 $4\times6$ 的 $(\cdot)^\odot$ 换成 $6\times6$ 的 $-(\cdot)^\curlywedge$、把 $\tilde{\mathbf{p}}$ 换成切向量 $\mathbf{x}$）。它在把“伴随作用过的切向量”再对位姿求导时直接套用，是连续时间协方差链式法则里的一块标准积木。
\end{derivation}

\begin{derivation}[运动学雅可比恒等式 $\dot{\mathbf{J}}-\boldsymbol{\omega}^\wedge\mathbf{J}=\partial\boldsymbol{\omega}/\partial\boldsymbol{\phi}$]
\cref{thm:so3-kinematics} 给了 $\boldsymbol{\omega}=\mathbf{J}\dot{\boldsymbol{\phi}}$。连续时间状态估计/预积分协方差推导里还要用到 $\mathbf{J}$ 的\emph{时间导数}与 $\boldsymbol{\omega}$ 对 $\boldsymbol{\phi}$ 的\emph{偏导}的关系\cite{barfoot2024state}：
\begin{equation}
  \dot{\mathbf{J}}(\boldsymbol{\phi})-\boldsymbol{\omega}^\wedge\mathbf{J}(\boldsymbol{\phi})\equiv\frac{\partial\boldsymbol{\omega}}{\partial\boldsymbol{\phi}},\qquad(\SE(3))\quad\dot{\boldsymbol{\mathcal{J}}}(\boldsymbol{\xi})-\boldsymbol{\varpi}^\curlywedge\boldsymbol{\mathcal{J}}(\boldsymbol{\xi})\equiv\frac{\partial\boldsymbol{\varpi}}{\partial\boldsymbol{\xi}}.
  \label{eq:kinematics-jac-identity}
\end{equation}
（此处取世界系 $\boldsymbol{\omega}=\boldsymbol{\omega}_w=\mathbf{J}_l\dot{\boldsymbol{\phi}}$、$\mathbf{J}=\mathbf{J}_l$。）\emph{推导}：对 $\boldsymbol{\omega}=\mathbf{J}\dot{\boldsymbol{\phi}}$ 用积分表示 $\mathbf{J}(\boldsymbol{\phi})=\int_0^1\exp(\alpha\boldsymbol{\phi}^\wedge)d\alpha$（\cref{eq:Jl-integral}），把 $\partial\boldsymbol{\omega}/\partial\boldsymbol{\phi}$ 写成对 $\boldsymbol{\phi}$ 的方向导数并分部积分（利用 $\tfrac{d}{d\alpha}\exp(\alpha\boldsymbol{\phi}^\wedge)=\boldsymbol{\phi}^\wedge\exp(\alpha\boldsymbol{\phi}^\wedge)$）：
\[
  \frac{\partial\boldsymbol{\omega}}{\partial\boldsymbol{\phi}}
  =-\int_0^1\big(\exp(\alpha\boldsymbol{\phi}^\wedge)\dot{\boldsymbol{\phi}}\big)^\wedge\,\alpha\,\mathbf{J}(\alpha\boldsymbol{\phi})\,d\alpha
  =-\boldsymbol{\omega}^\wedge\mathbf{J}(\boldsymbol{\phi})+\dot{\mathbf{J}}(\boldsymbol{\phi}),
\]
移项即 \cref{eq:kinematics-jac-identity} 第一式；$\SE(3)$ 版把 $\boldsymbol{\phi}^\wedge\to\boldsymbol{\xi}^\wedge$、$\boldsymbol{\omega}^\wedge\to\boldsymbol{\varpi}^\curlywedge$、$\mathbf{J}\to\boldsymbol{\mathcal{J}}$ 逐字平行。该恒等式在 IMU 预积分的连续时间协方差微分方程、连续时间 SLAM 的状态转移矩阵推导中作为基本工具出现。
\end{derivation}

% ────────────────────────────────────────────────
\subsection{齐次点算子 \texorpdfstring{$(\cdot)^\odot$}{odot} 与 \texorpdfstring{$(\cdot)^\circledcirc$}{circledcirc} 的代数}
\label{subsec:odot-algebra}

\paragraph{动机：把“旋量乘点”整理成“矩阵乘旋量”。}
\cref{sec:perturb} 推位姿作用点的右扰动雅可比时，关键一步是把 $\boldsymbol{\xi}^\wedge\tilde{\mathbf{p}}$（$\se(3)$ 元素乘齐次点）重写成 $\tilde{\mathbf{p}}^{\,\odot}\boldsymbol{\xi}$（一个只依赖点的 $4\times6$ 矩阵乘旋量向量），从而把对 $\boldsymbol{\xi}$ 的导数一眼读出。这套把 $\boldsymbol{\xi}$ 与 $\tilde{\mathbf{p}}$“换位”的算子，值得单独立其代数：它是点观测雅可比（\cref{eq:right-pert-se3}、以及 \cref{ch:vo} 重投影、\cref{ch:slam_est} 点-面残差）反复使用的地基\cite{barfoot2024state}。本小节给全 $(\cdot)^\odot$、与之配对的 $(\cdot)^\circledcirc$、以及把它们\emph{搬过位姿}的伴随恒等式。

\paragraph{两个算子：$4\times6$ 与 $6\times4$。}对齐次点 $\mathbf{p}=[\boldsymbol{\varepsilon}^\top,\eta]^\top\in\mathbb{R}^4$（$\boldsymbol{\varepsilon}\in\mathbb{R}^3$ 为前三维、$\eta$ 为齐次标量；普通点取 $\eta=1$、$\boldsymbol{\varepsilon}=\mathbf{p}$），定义
\begin{equation}
  \mathbf{p}^{\odot}
  =\begin{bmatrix}\boldsymbol{\varepsilon}\\\eta\end{bmatrix}^{\!\odot}
  =\begin{bmatrix}\eta\,\mathbf{I}&-\boldsymbol{\varepsilon}^\wedge\\\mathbf{0}^\top&\mathbf{0}^\top\end{bmatrix}\in\mathbb{R}^{4\times6},
  \qquad
  \mathbf{p}^{\circledcirc}
  =\begin{bmatrix}\boldsymbol{\varepsilon}\\\eta\end{bmatrix}^{\!\circledcirc}
  =\begin{bmatrix}\mathbf{0}&\boldsymbol{\varepsilon}\\-\boldsymbol{\varepsilon}^\wedge&\mathbf{0}\end{bmatrix}\in\mathbb{R}^{6\times4}.
  \label{eq:odot-circ-def}
\end{equation}
$(\cdot)^\odot$ 即 \cref{eq:right-pert-se3} 用过的 $4\times6$ 算子（取 $\eta=1$ 时右上块 $-\boldsymbol{\varepsilon}^\wedge=-\mathbf{p}^\wedge$、左上块 $\mathbf{I}$ 与彼处一致）；$(\cdot)^\circledcirc$ 是它的 $6\times4$ “转置位”伙伴。两者的列/行分块都按本书 $\boldsymbol{\xi}=[\boldsymbol{\rho};\boldsymbol{\phi}]$（平移在前、旋转在后）排布。

\paragraph{两条换位恒等式。}它们的用处全在下面这对恒等式：对任意 $\boldsymbol{\xi}\in\mathbb{R}^6$、$\mathbf{p}\in\mathbb{R}^4$，
\begin{equation}
  \boldsymbol{\xi}^\wedge\mathbf{p}\equiv\mathbf{p}^{\odot}\boldsymbol{\xi},
  \qquad
  \mathbf{p}^\top\boldsymbol{\xi}^\wedge\equiv\boldsymbol{\xi}^\top\mathbf{p}^{\circledcirc}.
  \label{eq:odot-identities}
\end{equation}
第一式把“$\se(3)$ 矩阵作用在点上”翻成“点的 $\odot$ 矩阵作用在旋量上”，正是求 $\partial(\boldsymbol{\xi}^\wedge\mathbf{p})/\partial\boldsymbol{\xi}=\mathbf{p}^\odot$ 的依据；第二式是其\emph{转置形态}，把“点的行向量左乘 $\se(3)$ 矩阵”翻成“旋量行向量乘点的 $\circledcirc$ 矩阵”，在标量型残差 $\mathbf{q}^\top(\mathbf{T}\mathbf{p})$（如点到平面距离、$\mathbf{q}$ 为法向）对 $\boldsymbol{\xi}$ 求梯度时直接登场——这就是 $(\cdot)^\circledcirc$ 存在的理由。两式都可由 \cref{eq:odot-circ-def} 的分块直接验证：例如第一式右上块给 $-\boldsymbol{\varepsilon}^\wedge\boldsymbol{\phi}=\boldsymbol{\phi}^\wedge\boldsymbol{\varepsilon}$、左上块给 $\eta\boldsymbol{\rho}$，正是 $\boldsymbol{\xi}^\wedge\mathbf{p}=[\boldsymbol{\phi}^\wedge\boldsymbol{\varepsilon}+\eta\boldsymbol{\rho};\,0]$。

\paragraph{把 $(\cdot)^\odot$ 搬过位姿：伴随恒等式。}最有用的是\emph{精确}恒等式（记 $\boldsymbol{\mathcal{T}}=\mathrm{Ad}(\mathbf{T})$）
\begin{equation}
  \boxed{\;(\mathbf{T}\mathbf{p})^{\odot}\equiv\mathbf{T}\,\mathbf{p}^{\odot}\,\boldsymbol{\mathcal{T}}^{-1}\;}
  \label{eq:Tp-odot}
\end{equation}
它把“先变换点、再取 $\odot$”等同于“先取 $\odot$、左乘 $\mathbf{T}$、右乘伴随之逆 $\boldsymbol{\mathcal{T}}^{-1}$”。\emph{推导}：对任意 $\boldsymbol{\xi}$，用 \cref{eq:odot-identities} 第一式与移指数恒等式 \cref{eq:Ad-move}（$\boldsymbol{\xi}^\wedge\mathbf{T}=\mathbf{T}(\mathrm{Ad}(\mathbf{T})^{-1}\boldsymbol{\xi})^\wedge=\mathbf{T}(\boldsymbol{\mathcal{T}}^{-1}\boldsymbol{\xi})^\wedge$，对 $4\times4$ 矩阵两边成立）有
\[
  (\mathbf{T}\mathbf{p})^{\odot}\boldsymbol{\xi}=\boldsymbol{\xi}^\wedge(\mathbf{T}\mathbf{p})=\mathbf{T}\,(\boldsymbol{\mathcal{T}}^{-1}\boldsymbol{\xi})^\wedge\mathbf{p}=\mathbf{T}\,\mathbf{p}^{\odot}\boldsymbol{\mathcal{T}}^{-1}\boldsymbol{\xi},
\]
对任意 $\boldsymbol{\xi}$ 成立即得 \cref{eq:Tp-odot}。它的直接后果就是 \cref{eq:right-pert-se3} 与左扰动版 \cref{eq:left-pert-se3} 之间的精确换算：右扰动雅可比 $\partial(\mathbf{T}\mathbf{p})/\partial\delta\boldsymbol{\xi}_{\text{右}}=\mathbf{T}\mathbf{p}^{\odot}$ 与左扰动雅可比 $(\mathbf{T}\mathbf{p})^{\odot}$ 恰好相差一个 $\boldsymbol{\mathcal{T}}^{-1}=\mathrm{Ad}(\mathbf{T})^{-1}$——与 \cref{sec:adjoint} 练习 2 的“左右扰动雅可比经伴随互换”逐字吻合。配套地还有度量型恒等式 $(\mathbf{T}\mathbf{p})^{\odot\top}(\mathbf{T}\mathbf{p})^{\odot}=\boldsymbol{\mathcal{T}}^{-\top}\mathbf{p}^{\odot\top}\mathbf{p}^{\odot}\boldsymbol{\mathcal{T}}^{-1}$（由 \cref{eq:Tp-odot} 两边各乘转置立得），在最小二乘的 Hessian $\sum(\mathbf{T}\mathbf{p})^{\odot\top}\mathbf{W}(\mathbf{T}\mathbf{p})^{\odot}$ 化简时常用。

\begin{insight}[$(\cdot)^\odot$ 之于点、犹如 $(\cdot)^\curlywedge$ 之于切向量]
对照 \cref{deriv:Ad-deriv} 的伴随导数 $\partial(\boldsymbol{\mathcal{T}}\mathbf{x})/\partial\boldsymbol{\xi}=-(\boldsymbol{\mathcal{T}}\mathbf{x})^\curlywedge\boldsymbol{\mathcal{J}}_l$ 与点观测导数 $\partial(\mathbf{T}\mathbf{p})/\partial\boldsymbol{\xi}=(\mathbf{T}\mathbf{p})^\odot\boldsymbol{\mathcal{J}}_l$（\cref{eq:right-pert-se3} 在李代数求导下的形态，残留 $\boldsymbol{\mathcal{J}}_l$；右扰动下 $\boldsymbol{\mathcal{J}}_l$ 消失为 \cref{eq:right-pert-se3}），可见 $(\cdot)^\odot$ 与 $(\cdot)^\curlywedge$ 在各自语境里扮演\emph{同一角色}：把“被作用对象”提到前面、把“扰动旋量”留在后面，使导数现成可读。\cref{eq:Tp-odot} 的 $(\mathbf{T}\mathbf{p})^\odot=\mathbf{T}\mathbf{p}^\odot\boldsymbol{\mathcal{T}}^{-1}$ 也与伴随的“移指数”$\mathbf{T}\,\mathrm{Exp}(\boldsymbol{\xi})=\mathrm{Exp}(\boldsymbol{\mathcal{T}}\boldsymbol{\xi})\mathbf{T}$ 同构——本章的伴随，无论作用在切向量、点、还是它们的算子上，都是同一个“跨越位姿搬运”的机制。
\end{insight}

% ════════════════════════════════════════════════════════════════
\section*{本章常见误解汇总}
\begin{table}[H]\centering\small
\begin{tabular}{p{0.46\textwidth} p{0.46\textwidth}}
\toprule
\textbf{常见误解} & \textbf{正确理解} \\
\midrule
$\SO(3)$ 是向量空间，可以加减 & 只对乘法封闭，是群非向量空间（\cref{eq:not-vs}）\\
$\mathbf{R}^{-1}=\mathbf{R}^\top$ 是巧合 & 正交性的定义性质，旋转保范的必然 \\
$\se(3)$ 的 $\boldsymbol{\rho}$ 就是平移 $\mathbf{t}$ & $\mathbf{t}=\mathbf{J}\boldsymbol{\rho}$，仅 $\boldsymbol{\phi}\to0$ 时相等（\cref{thm:se3-exp}）\\
指数映射一一对应 & 满射非单射，需限 $|\theta|<\pi$（\cref{prop:surjective}）\\
BCH 是精确等式 & 仅小扰动下的一阶近似（\cref{eq:bch-linear}）\\
左右扰动可随意混用 & 雅可比/更新/协方差须同侧，否则发散 \\
$\mathrm{Ad}(\mathbf{T})$ 是 $4\times4$ 位姿 & 是 $6\times6$、作用于切向量 $\boldsymbol{\xi}$ 的线性映射 \\
直接优化 $\mathbf{R}$ 的 $9$ 元素 & 过参数化 + 破坏约束；应用 $\boxplus$ 在 $3/6$ 维切空间更新 \\
轴角率就是角速度 $\dot{\boldsymbol{\phi}}=\boldsymbol{\omega}$ & 一般差一个 $\mathbf{J}$：$\boldsymbol{\omega}_b=\mathbf{J}_r\dot{\boldsymbol{\phi}}$，仅匀轴时相等（\cref{cor:fixed-axis}）\\
相关位姿可当独立复合/差分 & 须加/减交叉项 $\boldsymbol{\Sigma}_{12}$，符号相反（\cref{eq:corr-compound,eq:corr-diff}）\\
位姿可直接协方差加权平均 & $\SE(3)$ 无加法，须李群上 Gauss--Newton 融合（\cref{eq:fusion-cov}）\\
$(\mathbf{a}^\wedge)^3=-\mathbf{a}^\wedge$ 是孤立技巧 & 它是 $\so(3)$ 最小多项式，闭式收敛的统一源头（\cref{thm:minimal-poly}）\\
位姿可逐元素线性插值 & $(1-\alpha)\mathbf{R}_1+\alpha\mathbf{R}_2\notin\SO(3)$；须保群测地插值（\cref{thm:lie-interp}）\\
速度只是 $\mathbb{R}^3$，拼在 $\SE(3)$ 旁即可 & 丢群仿射结构、误差又依赖状态；应装进 $\mathrm{SE}_2(3)$（\cref{def:se23}）\\
\bottomrule
\end{tabular}
\end{table}

% ════════════════════════════════════════════════════════════════
\section*{本章小结}

\begin{table}[H]\centering\small
\caption{符号表}
\begin{tabular}{lll}
\toprule
符号 & 含义 & 首次出现 \\
\midrule
$\SO(3),\SE(3)$ & 三维旋转群 / 位姿群 & \cref{eq:so3,eq:se3} \\
$\so(3),\se(3)$ & 对应李代数 & \cref{def:so3-se3} \\
$(\cdot)^\wedge,(\cdot)^\vee$ & 反对称化 / 其逆（$\se(3)$ 为重载） & \cref{eq:hat} \\
$\mathrm{Exp},\mathrm{Log}$ & 吃/吐向量的指数 / 对数 & \cref{sec:exp-log} \\
$\boldsymbol{\phi}=\theta\mathbf{a},\ \boldsymbol{\xi}=[\boldsymbol{\rho};\boldsymbol{\phi}]$ & 轴角 / $\se(3)$ 坐标（平移在前） & \cref{def:so3-se3} \\
$\mathbf{J}_l,\mathbf{J}_r$ & 左 / 右雅可比（$3\times3$） & \cref{eq:Jl,eq:Jr} \\
$\boldsymbol{\mathcal{J}}_l,\mathbf{Q}_l$ & $\SE(3)$ 的 $6\times6$ 雅可比及其耦合块 & \cref{eq:se3-6x6} \\
$\mathbf{J}_s$ & $\mathrm{Sim}(3)$ 的尺度雅可比（含 $\sigma,\theta$） & \cref{eq:sim3-Js} \\
$\mathrm{SE}_2(3),\ \boldsymbol{\xi}=[\boldsymbol{\rho};\boldsymbol{\nu};\boldsymbol{\phi}]$ & 扩展位姿群（$5\times5$）/ 其 $9$ 维李代数坐标 & \cref{def:se23} \\
$\mathbf{A}(\alpha,\boldsymbol{\phi})$ & 李群插值的扰动权重矩阵 & \cref{eq:interp-perturb} \\
$\mathrm{Ad}(\mathbf{T}),\ \boldsymbol{\xi}^\curlywedge$ & 群伴随（$6\times6$）/ 小伴随 & \cref{eq:Ad,eq:ad} \\
$\boldsymbol{\omega}_b,\boldsymbol{\omega}_w,\dot{\boldsymbol{\phi}}$ & 机体/世界系角速度、李代数坐标率 & \cref{eq:so3-kinematics} \\
$(\cdot)^\odot,(\cdot)^\circledcirc$ & 齐次点 $\to 4\times6$ / $6\times4$ 的扰动算子 & \cref{eq:right-pert-se3,eq:odot-circ-def} \\
$\boxplus,\boxminus$ & 流形“加/减”（retraction/局部坐标） & \cref{eq:boxplus} \\
$|\det\mathbf{J}|$ & 李群体积元因子（Haar 测度） & \cref{eq:detJ} \\
$\bar{\mathbf{T}},\boldsymbol{\Sigma},\boldsymbol{\Sigma}_{12}$ & 位姿均值 / 切空间协方差 / 交叉协方差 & \cref{eq:lie-gaussian} \\
$\mathbf{G}_k$ & 位姿融合的雅可比 $-\boldsymbol{\mathcal{J}}_l^{-1}$ & \cref{eq:fusion-jacobian} \\
\bottomrule
\end{tabular}
\end{table}

\begin{table}[H]\centering\small
\caption{定理速查表}
\begin{tabular}{lll}
\toprule
定理 / 公式 & 一句话说明 & 对应 \\
\midrule
罗德里格斯 \cref{eq:rodrigues} & 轴角 $\to$ 旋转矩阵，避免无穷级数 & \cref{thm:so3-exp} \\
$\SE(3)$ 指数 \cref{eq:left-jacobian} & $\mathbf{t}=\mathbf{J}\boldsymbol{\rho}$，左雅可比 $\mathbf{J}$ 现身 & \cref{thm:se3-exp} \\
BCH 线性近似 \cref{eq:bch-linear} & 李代数加法 $\leftrightarrow$ 李群乘法，分左右 & \cref{sec:bch} \\
运动学 $\boldsymbol{\omega}=\mathbf{J}\dot{\boldsymbol{\phi}}$ \cref{eq:so3-kinematics} & $\dot{\boldsymbol{\phi}}=\mathbf{J}_r^{-1}\boldsymbol{\omega}_b$，连续预积分前置 & \cref{thm:so3-kinematics} \\
右扰动求导 \cref{eq:right-pert-so3} & $\partial(\mathbf{R}\mathbf{p})/\partial\boldsymbol{\varphi}=-\mathbf{R}\mathbf{p}^\wedge$，无 $\mathbf{J}$ & \cref{sec:perturb} \\
伴随移指数 \cref{eq:Ad-move} & 左右扰动精确互换的桥 & \cref{sec:adjoint} \\
齐次点算子 \cref{eq:Tp-odot} & $(\mathbf{T}\mathbf{p})^\odot=\mathbf{T}\mathbf{p}^\odot\boldsymbol{\mathcal{T}}^{-1}$，点雅可比地基 & \cref{subsec:odot-algebra} \\
直接级数 \cref{eq:se3-exp-direct} & 不分块整体求 $\mathbf{T}=\exp(\boldsymbol{\xi}^\wedge)$ & \cref{deriv:direct-series} \\
$\boxplus/\boxminus$ \cref{eq:boxplus} & 流形优化 = 切空间欧氏优化 + 收缩 & \cref{sec:manifold} \\
体积元 \cref{eq:detJ} & $d\mathbf{R}=|\det\mathbf{J}|d\boldsymbol{\phi}$，群上积分/密度归一化 & \cref{thm:detJ} \\
复合协方差 \cref{eq:compound-cov} & 伴随搬运切空间协方差 & \cref{sec:uncertainty} \\
位姿融合 \cref{eq:fusion-cov} & 李群上 Gauss--Newton = 信息相加 & \cref{subsec:pose-fusion} \\
最小多项式 \cref{thm:minimal-poly} & 一切闭式收敛的代数源头 & \cref{subsec:minimal-poly} \\
李群插值 \cref{eq:interp-so3} & $(\mathbf{R}_2\mathbf{R}_1^\top)^\alpha\mathbf{R}_1$ 保群 = 匀速运动 & \cref{thm:lie-interp} \\
$\mathrm{SE}_2(3)$ 指数 \cref{eq:se23-exp} & $\mathbf{r}=\mathbf{J}\boldsymbol{\rho},\mathbf{v}=\mathbf{J}\boldsymbol{\nu}$，姿/速/位同群 & \cref{thm:se23-exp} \\
不变/等变 \cref{eq:invariance,eq:equivariance} & 换系答案可换算；InEKF 一致性根源 & \cref{sec:lie-symmetry} \\
\bottomrule
\end{tabular}
\end{table}

\begin{table}[H]\centering\small
\caption{知识点总表}
\begin{tabular}{clll}
\toprule
\# & 知识点 & 核心要点 & 难度 \\
\midrule
1 & 群与矩阵李群 & 加法不封闭、乘法封闭；封结幺逆 & \(\bigstar\bigstar\) \\
2 & 李代数 & $\mathbf{R}\mathbf{R}^\top=\mathbf{I}$ 求导挤出反对称 & \(\bigstar\bigstar\bigstar\) \\
3 & 指数/对数 & Rodrigues、$\SE(3)$ 指数、$\mathbf{J}$ & \(\bigstar\bigstar\bigstar\) \\
4 & BCH 与 $\mathbf{J}_l,\mathbf{J}_r$ & 李代数加 $\leftrightarrow$ 群乘的线性近似 & \(\bigstar\bigstar\bigstar\bigstar\) \\
5 & 运动学 $\boldsymbol{\omega}=\mathbf{J}\dot{\boldsymbol{\phi}}$ & $\mathbf{J}$ 的运动学身份；$\dot{\boldsymbol{\phi}}=\mathbf{J}_r^{-1}\boldsymbol{\omega}_b$ & \(\bigstar\bigstar\bigstar\bigstar\) \\
6 & 右扰动模型 & $-\mathbf{R}\mathbf{p}^\wedge$；更新右乘指数 & \(\bigstar\bigstar\bigstar\bigstar\) \\
7 & 伴随 $\mathrm{Ad},\mathrm{ad}$ & 左右扰动/坐标系/协方差的桥 & \(\bigstar\bigstar\bigstar\bigstar\) \\
8 & 流形/$\boxplus\boxminus$ & 对接因子图、GTSAM/Ceres & \(\bigstar\bigstar\bigstar\bigstar\) \\
9 & 距离/体积/积分 & $|\det\mathbf{J}|$ 体积元；群上密度归一化 & \(\bigstar\bigstar\bigstar\bigstar\) \\
10 & 李群不确定度 & 切空间高斯 + 伴随传播 & \(\bigstar\bigstar\bigstar\bigstar\bigstar\) \\
11 & 位姿融合/相关复合 & 李群 GN；信息相加；交叉项 & \(\bigstar\bigstar\bigstar\bigstar\bigstar\) \\
12 & 最小多项式/Jordan & 三个降阶恒等式的统一出处 & \(\bigstar\bigstar\bigstar\bigstar\bigstar\) \\
13 & 李群插值 & 保群（测地）插值 = 匀速运动 & \(\bigstar\bigstar\bigstar\) \\
14 & 扩展位姿 $\mathrm{SE}_2(3)$ & 姿/速/位同群；InEKF/预积分地基 & \(\bigstar\bigstar\bigstar\bigstar\) \\
15 & 对称/不变/等变 & 换系不变；InEKF 一致性的群论语言 & \(\bigstar\bigstar\bigstar\bigstar\) \\
\bottomrule
\end{tabular}
\end{table}

% ════════════════════════════════════════════════════════════════
\section*{延伸阅读}
\begin{itemize}
  \item \textbf{严谨主线}：Barfoot《State Estimation for Robotics》\cite{barfoot2024state} 第 7--8 章——矩阵李群、$6\times6$ 雅可比、李群上的不确定度与位姿复合（本章不确定度一节的完整出处）。
  \item \textbf{几何直觉与代码}：视觉 SLAM 十四讲\cite{gaoxiang2019slam14} 第 4 章——招牌的“从 $\mathbf{R}\mathbf{R}^\top=\mathbf{I}$ 引出李代数”、Sophus 实践。
  \item \textbf{现代统一记号}：Solà 等《A micro Lie theory for state estimation in robotics》\cite{sola2018micro}——右/左 $\oplus\,\ominus$、函数的右雅可比与链式法则，强烈推荐对照本书右扰动约定阅读。
  \item \textbf{流形优化接口}：Hertzberg 等\cite{hertzberg2013integrating}（$\boxplus/\boxminus$ 封装的原始论文）、SLAM Handbook\cite{carlone2026handbook} 第 2 章（流形/retraction/因子图衔接）。
  \item \textbf{随机李群理论}：Chirikjian\cite{chirikjian2011stochastic}——李群上概率与信息论的奠基性参考。
  \item \textbf{$\mathrm{SE}_2(3)$ 与不变滤波}：Barrau--Bonnabel\cite{barrau2017invariant}（InEKF 与群仿射/log-linear 性质的奠基论文）、Hartley 等\cite{hartley2020contact}（足式机器人接触辅助 InEKF、“不完美 InEKF”）、Brossard 等\cite{brossard2021associating}（$\mathrm{SE}_2(3)$ 扩展位姿上的预积分与不确定度）；等变滤波见 Mahony--Trumpf\cite{mahony2021equivariant}。本章 \cref{sec:se23,sec:lie-symmetry} 为这条线提供群结构与对称性地基。
\end{itemize}

\section*{本章与后续章节的关系}
\begin{table}[H]\centering\small
\begin{tabular}{lll}
\toprule
后续章节 & 关系 & 本章铺垫 \\
\midrule
\cref{ch:nlopt} 非线性优化 & 位姿作为流形变量做 GN/LM & 右扰动雅可比、$\boxplus$ 更新 \\
\cref{ch:slam_est} 状态估计 & 携带协方差的位姿估计 & 切空间高斯、伴随传播 \\
\cref{ch:vo} 视觉里程计 & 重投影误差对位姿求导 & \cref{eq:right-pert-se3} 的 $\mathbf{T}\tilde{\mathbf{p}}^\odot$ \\
\cref{ch:vio} VIO & IMU 预积分在 $\SE(3)$/$\SO(3)$ 上 & 伴随、$\SE(3)$ 雅可比、不确定度 \\
\cref{ch:eskf} 误差状态滤波 & InEKF / 不变滤波一致性 & $\mathrm{SE}_2(3)$（\cref{sec:se23}）、对称/不变/等变（\cref{sec:lie-symmetry}）\\
\cref{ch:imu} IMU 模型 & 扩展位姿预积分、连续轨迹 & $\mathrm{SE}_2(3)$、李群插值（\cref{sec:lie-interp}）\\
\bottomrule
\end{tabular}
\end{table}

\section*{故障排查手册}
\begin{enumerate}
  \item \textbf{症状}：BA / 位姿图\emph{不收敛或发散}。\textbf{排查}：检查雅可比扰动侧与更新乘法侧是否同为右侧（\cref{sec:perturb} 陷阱）；检查是否给求解器注册了 $\SO(3)/\SE(3)$ 的流形/局部参数化（\cref{sec:manifold} 陷阱）。
  \item \textbf{症状}：大旋转处残差/雅可比\emph{突变}。\textbf{原因}：$\mathrm{Log}$ 在 $|\theta|\to\pi$ 的 $2\pi$ 跳变（\cref{prop:surjective}）。\textbf{对策}：对相对旋转取 $\mathrm{Log}$、保持每步增量小。
  \item \textbf{症状}：协方差传播后\emph{不对称或非正定}。\textbf{排查}：确认用 $\mathrm{Ad}(\bar{\mathbf{T}}_1)\boldsymbol{\Sigma}_2\mathrm{Ad}(\bar{\mathbf{T}}_1)^\top$（\cref{eq:compound-cov}）而非漏掉转置；大不确定度需四阶修正（\cref{eq:fourth-order}）。
  \item \textbf{症状}：$\mathbf{J}$ 或 $\mathbf{J}^{-1}$ 计算出 \texttt{NaN}。\textbf{原因}：$\theta\to0$ 时 $\sin\theta/\theta$、$\theta\cot(\theta/2)$ 的 $0/0$。\textbf{对策}：小角度用泰勒展开分支（$\mathbf{J}\approx\mathbf{I}+\tfrac12\boldsymbol{\phi}^\wedge$）。
  \item \textbf{症状}：Sophus 里 $\se(3)$ 分量顺序与公式对不上。\textbf{原因}：Sophus/本书为“平移在前、旋转在后”，部分文献（JPL、某些 GTSAM 约定）相反。\textbf{对策}：核对 $\boldsymbol{\xi}=[\boldsymbol{\rho};\boldsymbol{\phi}]$ 顺序后再套雅可比。
\end{enumerate}
